\subsection{Orchestration Frameworks}
\paragraph{MCQ-0521 [medium]} A design requires this behavior: LangGraph models long-running agent workflows as stateful graphs with checkpoints, interrupts, streaming, and durable execution. Which mechanism should the team study?
\begin{enumerate}[label=\Alph*.]
\item LlamaIndex workflows
\item DSPy
\item LangGraph
\item Airflow, Dagster, and Prefect
\end{enumerate}
\noindent\textbf{Answer.} LangGraph
\par\textit{LangGraph models long-running agent workflows as stateful graphs with checkpoints, interrupts, streaming, and durable execution. Tradeoff: Fine-grained control aids production reliability but exposes more state and graph design decisions. Watch for: Non-idempotent node side effects make checkpoint replay unsafe.}
\paragraph{MCQ-0522 [hard]} The team proposes LangGraph for a real workload. Which finding gives the correct decision boundary?
\begin{enumerate}[label=\Alph*.]
\item Systematic optimization can beat manual prompting, while it depends on representative data and stable metrics.
\item They excel at batch pipelines but are not automatically low-latency interactive agent runtimes.
\item Fine-grained control aids production reliability but exposes more state and graph design decisions.
\item Strong data and indexing integrations speed RAG work but can couple application flow to framework abstractions.
\end{enumerate}
\noindent\textbf{Answer.} Fine-grained control aids production reliability but exposes more state and graph design decisions.
\par\textit{LangGraph models long-running agent workflows as stateful graphs with checkpoints, interrupts, streaming, and durable execution. Tradeoff: Fine-grained control aids production reliability but exposes more state and graph design decisions. Watch for: Non-idempotent node side effects make checkpoint replay unsafe.}
\paragraph{MCQ-0523 [hard]} Before releasing LangGraph, which failure deserves a dedicated regression test?
\begin{enumerate}[label=\Alph*.]
\item Running per-token or per-chat-turn work as heavy batch tasks adds unacceptable overhead.
\item Optimizing against a leaky judge overfits wording instead of capability.
\item Non-idempotent node side effects make checkpoint replay unsafe.
\item Hiding all retrieval behind defaults makes evaluation and migration difficult.
\end{enumerate}
\noindent\textbf{Answer.} Non-idempotent node side effects make checkpoint replay unsafe.
\par\textit{LangGraph models long-running agent workflows as stateful graphs with checkpoints, interrupts, streaming, and durable execution. Tradeoff: Fine-grained control aids production reliability but exposes more state and graph design decisions. Watch for: Non-idempotent node side effects make checkpoint replay unsafe.}
\paragraph{MCQ-0524 [medium]} A design review reveals this mistake: Non-idempotent node side effects make checkpoint replay unsafe. Which mechanism should be traced first?
\begin{enumerate}[label=\Alph*.]
\item LlamaIndex workflows
\item LangGraph
\item Airflow, Dagster, and Prefect
\item DSPy
\end{enumerate}
\noindent\textbf{Answer.} LangGraph
\par\textit{LangGraph models long-running agent workflows as stateful graphs with checkpoints, interrupts, streaming, and durable execution. Tradeoff: Fine-grained control aids production reliability but exposes more state and graph design decisions. Watch for: Non-idempotent node side effects make checkpoint replay unsafe.}
\paragraph{MCQ-0525 [hard]} Which causal explanation of LangGraph connects what it does to when it is worth its cost?
\begin{enumerate}[label=\Alph*.]
\item Mechanism: DSPy treats prompts and modules as optimizable programs evaluated against examples and metrics. Consequence: Systematic optimization can beat manual prompting, while it depends on representative data and stable metrics.
\item Mechanism: Event-driven workflows coordinate retrieval, tools, agents, and structured steps in the LlamaIndex ecosystem. Consequence: Strong data and indexing integrations speed RAG work but can couple application flow to framework abstractions.
\item Mechanism: LangGraph models long-running agent workflows as stateful graphs with checkpoints, interrupts, streaming, and durable execution. Consequence: Fine-grained control aids production reliability but exposes more state and graph design decisions.
\item Mechanism: General-purpose data orchestrators schedule DAGs, assets, or flows with retries, lineage, and operational UIs. Consequence: They excel at batch pipelines but are not automatically low-latency interactive agent runtimes.
\end{enumerate}
\noindent\textbf{Answer.} Mechanism: LangGraph models long-running agent workflows as stateful graphs with checkpoints, interrupts, streaming, and durable execution. Consequence: Fine-grained control aids production reliability but exposes more state and graph design decisions.
\par\textit{LangGraph models long-running agent workflows as stateful graphs with checkpoints, interrupts, streaming, and durable execution. Tradeoff: Fine-grained control aids production reliability but exposes more state and graph design decisions. Watch for: Non-idempotent node side effects make checkpoint replay unsafe.}
\paragraph{FC-0526 [medium]} Define LangGraph in interview-ready language.
\noindent\textbf{Answer.} LangGraph models long-running agent workflows as stateful graphs with checkpoints, interrupts, streaming, and durable execution.
\par\textit{LangGraph models long-running agent workflows as stateful graphs with checkpoints, interrupts, streaming, and durable execution.}
\paragraph{FC-0527 [hard]} State the main tradeoff and production pitfall for LangGraph.
\noindent\textbf{Answer.} Tradeoff: Fine-grained control aids production reliability but exposes more state and graph design decisions. Pitfall: Non-idempotent node side effects make checkpoint replay unsafe.
\par\textit{Tradeoff: Fine-grained control aids production reliability but exposes more state and graph design decisions. Pitfall: Non-idempotent node side effects make checkpoint replay unsafe.}
\paragraph{LA-0528 [hard]} Explain LangGraph from first principles, then show how you would evaluate and operate it in production.
\noindent\textbf{Answer.} Start with the mechanism: LangGraph models long-running agent workflows as stateful graphs with checkpoints, interrupts, streaming, and durable execution. Make the tradeoff explicit: Fine-grained control aids production reliability but exposes more state and graph design decisions. Define workload-specific offline metrics and a latency/cost budget, test important slices, instrument the critical path, and create a rollback criterion. The production review must explicitly guard against this failure: Non-idempotent node side effects make checkpoint replay unsafe.
\par\textit{A strong answer connects mechanism, measurable tradeoffs, failure isolation, observability, and rollback rather than listing product features.}
\paragraph{MCQ-0529 [medium]} A design requires this behavior: Event-driven workflows coordinate retrieval, tools, agents, and structured steps in the LlamaIndex ecosystem. Which mechanism should the team study?
\begin{enumerate}[label=\Alph*.]
\item DSPy
\item AutoGen
\item LlamaIndex workflows
\item Airflow, Dagster, and Prefect
\end{enumerate}
\noindent\textbf{Answer.} LlamaIndex workflows
\par\textit{Event-driven workflows coordinate retrieval, tools, agents, and structured steps in the LlamaIndex ecosystem. Tradeoff: Strong data and indexing integrations speed RAG work but can couple application flow to framework abstractions. Watch for: Hiding all retrieval behind defaults makes evaluation and migration difficult.}
\paragraph{MCQ-0530 [hard]} The team proposes LlamaIndex workflows for a real workload. Which finding gives the correct decision boundary?
\begin{enumerate}[label=\Alph*.]
\item Role decomposition enables complex workflows but multiplies calls, coordination failures, and evaluation cost.
\item Systematic optimization can beat manual prompting, while it depends on representative data and stable metrics.
\item Strong data and indexing integrations speed RAG work but can couple application flow to framework abstractions.
\item They excel at batch pipelines but are not automatically low-latency interactive agent runtimes.
\end{enumerate}
\noindent\textbf{Answer.} Strong data and indexing integrations speed RAG work but can couple application flow to framework abstractions.
\par\textit{Event-driven workflows coordinate retrieval, tools, agents, and structured steps in the LlamaIndex ecosystem. Tradeoff: Strong data and indexing integrations speed RAG work but can couple application flow to framework abstractions. Watch for: Hiding all retrieval behind defaults makes evaluation and migration difficult.}
\paragraph{MCQ-0531 [hard]} Before releasing LlamaIndex workflows, which failure deserves a dedicated regression test?
\begin{enumerate}[label=\Alph*.]
\item Hiding all retrieval behind defaults makes evaluation and migration difficult.
\item Optimizing against a leaky judge overfits wording instead of capability.
\item Adding agents without independent capabilities creates expensive echo chambers.
\item Running per-token or per-chat-turn work as heavy batch tasks adds unacceptable overhead.
\end{enumerate}
\noindent\textbf{Answer.} Hiding all retrieval behind defaults makes evaluation and migration difficult.
\par\textit{Event-driven workflows coordinate retrieval, tools, agents, and structured steps in the LlamaIndex ecosystem. Tradeoff: Strong data and indexing integrations speed RAG work but can couple application flow to framework abstractions. Watch for: Hiding all retrieval behind defaults makes evaluation and migration difficult.}
\paragraph{MCQ-0532 [medium]} A design review reveals this mistake: Hiding all retrieval behind defaults makes evaluation and migration difficult. Which mechanism should be traced first?
\begin{enumerate}[label=\Alph*.]
\item AutoGen
\item DSPy
\item Airflow, Dagster, and Prefect
\item LlamaIndex workflows
\end{enumerate}
\noindent\textbf{Answer.} LlamaIndex workflows
\par\textit{Event-driven workflows coordinate retrieval, tools, agents, and structured steps in the LlamaIndex ecosystem. Tradeoff: Strong data and indexing integrations speed RAG work but can couple application flow to framework abstractions. Watch for: Hiding all retrieval behind defaults makes evaluation and migration difficult.}
\paragraph{MCQ-0533 [hard]} Which causal explanation of LlamaIndex workflows connects what it does to when it is worth its cost?
\begin{enumerate}[label=\Alph*.]
\item Mechanism: AutoGen structures multi-agent conversations and tool-using participants for collaborative task execution. Consequence: Role decomposition enables complex workflows but multiplies calls, coordination failures, and evaluation cost.
\item Mechanism: DSPy treats prompts and modules as optimizable programs evaluated against examples and metrics. Consequence: Systematic optimization can beat manual prompting, while it depends on representative data and stable metrics.
\item Mechanism: General-purpose data orchestrators schedule DAGs, assets, or flows with retries, lineage, and operational UIs. Consequence: They excel at batch pipelines but are not automatically low-latency interactive agent runtimes.
\item Mechanism: Event-driven workflows coordinate retrieval, tools, agents, and structured steps in the LlamaIndex ecosystem. Consequence: Strong data and indexing integrations speed RAG work but can couple application flow to framework abstractions.
\end{enumerate}
\noindent\textbf{Answer.} Mechanism: Event-driven workflows coordinate retrieval, tools, agents, and structured steps in the LlamaIndex ecosystem. Consequence: Strong data and indexing integrations speed RAG work but can couple application flow to framework abstractions.
\par\textit{Event-driven workflows coordinate retrieval, tools, agents, and structured steps in the LlamaIndex ecosystem. Tradeoff: Strong data and indexing integrations speed RAG work but can couple application flow to framework abstractions. Watch for: Hiding all retrieval behind defaults makes evaluation and migration difficult.}
\paragraph{FC-0534 [medium]} Define LlamaIndex workflows in interview-ready language.
\noindent\textbf{Answer.} Event-driven workflows coordinate retrieval, tools, agents, and structured steps in the LlamaIndex ecosystem.
\par\textit{Event-driven workflows coordinate retrieval, tools, agents, and structured steps in the LlamaIndex ecosystem.}
\paragraph{FC-0535 [hard]} State the main tradeoff and production pitfall for LlamaIndex workflows.
\noindent\textbf{Answer.} Tradeoff: Strong data and indexing integrations speed RAG work but can couple application flow to framework abstractions. Pitfall: Hiding all retrieval behind defaults makes evaluation and migration difficult.
\par\textit{Tradeoff: Strong data and indexing integrations speed RAG work but can couple application flow to framework abstractions. Pitfall: Hiding all retrieval behind defaults makes evaluation and migration difficult.}
\paragraph{LA-0536 [hard]} Explain LlamaIndex workflows from first principles, then show how you would evaluate and operate it in production.
\noindent\textbf{Answer.} Start with the mechanism: Event-driven workflows coordinate retrieval, tools, agents, and structured steps in the LlamaIndex ecosystem. Make the tradeoff explicit: Strong data and indexing integrations speed RAG work but can couple application flow to framework abstractions. Define workload-specific offline metrics and a latency/cost budget, test important slices, instrument the critical path, and create a rollback criterion. The production review must explicitly guard against this failure: Hiding all retrieval behind defaults makes evaluation and migration difficult.
\par\textit{A strong answer connects mechanism, measurable tradeoffs, failure isolation, observability, and rollback rather than listing product features.}
\paragraph{MCQ-0537 [medium]} A design requires this behavior: Haystack connects typed components into retrieval, generation, routing, and agent pipelines. Which mechanism should the team study?
\begin{enumerate}[label=\Alph*.]
\item CrewAI
\item LlamaIndex workflows
\item Airflow, Dagster, and Prefect
\item Haystack pipelines
\end{enumerate}
\noindent\textbf{Answer.} Haystack pipelines
\par\textit{Haystack connects typed components into retrieval, generation, routing, and agent pipelines. Tradeoff: Explicit components support testing, while custom graphs still require lifecycle discipline. Watch for: Allowing incompatible document schemas across components creates runtime surprises.}
\paragraph{MCQ-0538 [hard]} The team proposes Haystack pipelines for a real workload. Which finding gives the correct decision boundary?
\begin{enumerate}[label=\Alph*.]
\item They excel at batch pipelines but are not automatically low-latency interactive agent runtimes.
\item Explicit components support testing, while custom graphs still require lifecycle discipline.
\item It gives accessible orchestration but still needs explicit state, permissions, and observability.
\item Strong data and indexing integrations speed RAG work but can couple application flow to framework abstractions.
\end{enumerate}
\noindent\textbf{Answer.} Explicit components support testing, while custom graphs still require lifecycle discipline.
\par\textit{Haystack connects typed components into retrieval, generation, routing, and agent pipelines. Tradeoff: Explicit components support testing, while custom graphs still require lifecycle discipline. Watch for: Allowing incompatible document schemas across components creates runtime surprises.}
\paragraph{MCQ-0539 [hard]} Before releasing Haystack pipelines, which failure deserves a dedicated regression test?
\begin{enumerate}[label=\Alph*.]
\item Running per-token or per-chat-turn work as heavy batch tasks adds unacceptable overhead.
\item Allowing incompatible document schemas across components creates runtime surprises.
\item Hiding all retrieval behind defaults makes evaluation and migration difficult.
\item Persona descriptions cannot substitute for tool-level access control.
\end{enumerate}
\noindent\textbf{Answer.} Allowing incompatible document schemas across components creates runtime surprises.
\par\textit{Haystack connects typed components into retrieval, generation, routing, and agent pipelines. Tradeoff: Explicit components support testing, while custom graphs still require lifecycle discipline. Watch for: Allowing incompatible document schemas across components creates runtime surprises.}
\paragraph{MCQ-0540 [medium]} A design review reveals this mistake: Allowing incompatible document schemas across components creates runtime surprises. Which mechanism should be traced first?
\begin{enumerate}[label=\Alph*.]
\item Airflow, Dagster, and Prefect
\item CrewAI
\item LlamaIndex workflows
\item Haystack pipelines
\end{enumerate}
\noindent\textbf{Answer.} Haystack pipelines
\par\textit{Haystack connects typed components into retrieval, generation, routing, and agent pipelines. Tradeoff: Explicit components support testing, while custom graphs still require lifecycle discipline. Watch for: Allowing incompatible document schemas across components creates runtime surprises.}
\paragraph{MCQ-0541 [hard]} Which causal explanation of Haystack pipelines connects what it does to when it is worth its cost?
\begin{enumerate}[label=\Alph*.]
\item Mechanism: CrewAI organizes role-based agents, tasks, crews, and flows for sequential or coordinated work. Consequence: It gives accessible orchestration but still needs explicit state, permissions, and observability.
\item Mechanism: Haystack connects typed components into retrieval, generation, routing, and agent pipelines. Consequence: Explicit components support testing, while custom graphs still require lifecycle discipline.
\item Mechanism: Event-driven workflows coordinate retrieval, tools, agents, and structured steps in the LlamaIndex ecosystem. Consequence: Strong data and indexing integrations speed RAG work but can couple application flow to framework abstractions.
\item Mechanism: General-purpose data orchestrators schedule DAGs, assets, or flows with retries, lineage, and operational UIs. Consequence: They excel at batch pipelines but are not automatically low-latency interactive agent runtimes.
\end{enumerate}
\noindent\textbf{Answer.} Mechanism: Haystack connects typed components into retrieval, generation, routing, and agent pipelines. Consequence: Explicit components support testing, while custom graphs still require lifecycle discipline.
\par\textit{Haystack connects typed components into retrieval, generation, routing, and agent pipelines. Tradeoff: Explicit components support testing, while custom graphs still require lifecycle discipline. Watch for: Allowing incompatible document schemas across components creates runtime surprises.}
\paragraph{FC-0542 [medium]} Define Haystack pipelines in interview-ready language.
\noindent\textbf{Answer.} Haystack connects typed components into retrieval, generation, routing, and agent pipelines.
\par\textit{Haystack connects typed components into retrieval, generation, routing, and agent pipelines.}
\paragraph{FC-0543 [hard]} State the main tradeoff and production pitfall for Haystack pipelines.
\noindent\textbf{Answer.} Tradeoff: Explicit components support testing, while custom graphs still require lifecycle discipline. Pitfall: Allowing incompatible document schemas across components creates runtime surprises.
\par\textit{Tradeoff: Explicit components support testing, while custom graphs still require lifecycle discipline. Pitfall: Allowing incompatible document schemas across components creates runtime surprises.}
\paragraph{LA-0544 [hard]} Explain Haystack pipelines from first principles, then show how you would evaluate and operate it in production.
\noindent\textbf{Answer.} Start with the mechanism: Haystack connects typed components into retrieval, generation, routing, and agent pipelines. Make the tradeoff explicit: Explicit components support testing, while custom graphs still require lifecycle discipline. Define workload-specific offline metrics and a latency/cost budget, test important slices, instrument the critical path, and create a rollback criterion. The production review must explicitly guard against this failure: Allowing incompatible document schemas across components creates runtime surprises.
\par\textit{A strong answer connects mechanism, measurable tradeoffs, failure isolation, observability, and rollback rather than listing product features.}
\paragraph{MCQ-0545 [medium]} A design requires this behavior: PydanticAI uses typed dependencies, tools, results, and validation to build model-agnostic Python agents. Which mechanism should the team study?
\begin{enumerate}[label=\Alph*.]
\item Temporal
\item PydanticAI
\item LangGraph
\item DSPy
\end{enumerate}
\noindent\textbf{Answer.} PydanticAI
\par\textit{PydanticAI uses typed dependencies, tools, results, and validation to build model-agnostic Python agents. Tradeoff: Type safety improves integration but cannot validate semantic truth. Watch for: Confusing schema conformance with a correct answer produces false confidence.}
\paragraph{MCQ-0546 [hard]} The team proposes PydanticAI for a real workload. Which finding gives the correct decision boundary?
\begin{enumerate}[label=\Alph*.]
\item Type safety improves integration but cannot validate semantic truth.
\item Fine-grained control aids production reliability but exposes more state and graph design decisions.
\item It provides robust retries and long-running execution but requires deterministic workflow discipline.
\item Systematic optimization can beat manual prompting, while it depends on representative data and stable metrics.
\end{enumerate}
\noindent\textbf{Answer.} Type safety improves integration but cannot validate semantic truth.
\par\textit{PydanticAI uses typed dependencies, tools, results, and validation to build model-agnostic Python agents. Tradeoff: Type safety improves integration but cannot validate semantic truth. Watch for: Confusing schema conformance with a correct answer produces false confidence.}
\paragraph{MCQ-0547 [hard]} Before releasing PydanticAI, which failure deserves a dedicated regression test?
\begin{enumerate}[label=\Alph*.]
\item Reading wall-clock time or random state directly in replayed code causes nondeterminism.
\item Non-idempotent node side effects make checkpoint replay unsafe.
\item Optimizing against a leaky judge overfits wording instead of capability.
\item Confusing schema conformance with a correct answer produces false confidence.
\end{enumerate}
\noindent\textbf{Answer.} Confusing schema conformance with a correct answer produces false confidence.
\par\textit{PydanticAI uses typed dependencies, tools, results, and validation to build model-agnostic Python agents. Tradeoff: Type safety improves integration but cannot validate semantic truth. Watch for: Confusing schema conformance with a correct answer produces false confidence.}
\paragraph{MCQ-0548 [medium]} A design review reveals this mistake: Confusing schema conformance with a correct answer produces false confidence. Which mechanism should be traced first?
\begin{enumerate}[label=\Alph*.]
\item LangGraph
\item Temporal
\item PydanticAI
\item DSPy
\end{enumerate}
\noindent\textbf{Answer.} PydanticAI
\par\textit{PydanticAI uses typed dependencies, tools, results, and validation to build model-agnostic Python agents. Tradeoff: Type safety improves integration but cannot validate semantic truth. Watch for: Confusing schema conformance with a correct answer produces false confidence.}
\paragraph{MCQ-0549 [hard]} Which causal explanation of PydanticAI connects what it does to when it is worth its cost?
\begin{enumerate}[label=\Alph*.]
\item Mechanism: PydanticAI uses typed dependencies, tools, results, and validation to build model-agnostic Python agents. Consequence: Type safety improves integration but cannot validate semantic truth.
\item Mechanism: DSPy treats prompts and modules as optimizable programs evaluated against examples and metrics. Consequence: Systematic optimization can beat manual prompting, while it depends on representative data and stable metrics.
\item Mechanism: LangGraph models long-running agent workflows as stateful graphs with checkpoints, interrupts, streaming, and durable execution. Consequence: Fine-grained control aids production reliability but exposes more state and graph design decisions.
\item Mechanism: Temporal persists workflow history and replays deterministic workflow code around durable activities. Consequence: It provides robust retries and long-running execution but requires deterministic workflow discipline.
\end{enumerate}
\noindent\textbf{Answer.} Mechanism: PydanticAI uses typed dependencies, tools, results, and validation to build model-agnostic Python agents. Consequence: Type safety improves integration but cannot validate semantic truth.
\par\textit{PydanticAI uses typed dependencies, tools, results, and validation to build model-agnostic Python agents. Tradeoff: Type safety improves integration but cannot validate semantic truth. Watch for: Confusing schema conformance with a correct answer produces false confidence.}
\paragraph{FC-0550 [medium]} Define PydanticAI in interview-ready language.
\noindent\textbf{Answer.} PydanticAI uses typed dependencies, tools, results, and validation to build model-agnostic Python agents.
\par\textit{PydanticAI uses typed dependencies, tools, results, and validation to build model-agnostic Python agents.}
\paragraph{FC-0551 [hard]} State the main tradeoff and production pitfall for PydanticAI.
\noindent\textbf{Answer.} Tradeoff: Type safety improves integration but cannot validate semantic truth. Pitfall: Confusing schema conformance with a correct answer produces false confidence.
\par\textit{Tradeoff: Type safety improves integration but cannot validate semantic truth. Pitfall: Confusing schema conformance with a correct answer produces false confidence.}
\paragraph{LA-0552 [hard]} Explain PydanticAI from first principles, then show how you would evaluate and operate it in production.
\noindent\textbf{Answer.} Start with the mechanism: PydanticAI uses typed dependencies, tools, results, and validation to build model-agnostic Python agents. Make the tradeoff explicit: Type safety improves integration but cannot validate semantic truth. Define workload-specific offline metrics and a latency/cost budget, test important slices, instrument the critical path, and create a rollback criterion. The production review must explicitly guard against this failure: Confusing schema conformance with a correct answer produces false confidence.
\par\textit{A strong answer connects mechanism, measurable tradeoffs, failure isolation, observability, and rollback rather than listing product features.}
\paragraph{MCQ-0553 [medium]} A design requires this behavior: Semantic Kernel offers planners, agents, plugins, memory connectors, and multi-language enterprise integration. Which mechanism should the team study?
\begin{enumerate}[label=\Alph*.]
\item DSPy
\item AutoGen
\item Semantic Kernel
\item PydanticAI
\end{enumerate}
\noindent\textbf{Answer.} Semantic Kernel
\par\textit{Semantic Kernel offers planners, agents, plugins, memory connectors, and multi-language enterprise integration. Tradeoff: Its broad integration surface suits platform teams but introduces framework concepts and version coupling. Watch for: Registering overly broad plugins grants the model unnecessary authority.}
\paragraph{MCQ-0554 [hard]} The team proposes Semantic Kernel for a real workload. Which finding gives the correct decision boundary?
\begin{enumerate}[label=\Alph*.]
\item Its broad integration surface suits platform teams but introduces framework concepts and version coupling.
\item Role decomposition enables complex workflows but multiplies calls, coordination failures, and evaluation cost.
\item Type safety improves integration but cannot validate semantic truth.
\item Systematic optimization can beat manual prompting, while it depends on representative data and stable metrics.
\end{enumerate}
\noindent\textbf{Answer.} Its broad integration surface suits platform teams but introduces framework concepts and version coupling.
\par\textit{Semantic Kernel offers planners, agents, plugins, memory connectors, and multi-language enterprise integration. Tradeoff: Its broad integration surface suits platform teams but introduces framework concepts and version coupling. Watch for: Registering overly broad plugins grants the model unnecessary authority.}
\paragraph{MCQ-0555 [hard]} Before releasing Semantic Kernel, which failure deserves a dedicated regression test?
\begin{enumerate}[label=\Alph*.]
\item Registering overly broad plugins grants the model unnecessary authority.
\item Adding agents without independent capabilities creates expensive echo chambers.
\item Confusing schema conformance with a correct answer produces false confidence.
\item Optimizing against a leaky judge overfits wording instead of capability.
\end{enumerate}
\noindent\textbf{Answer.} Registering overly broad plugins grants the model unnecessary authority.
\par\textit{Semantic Kernel offers planners, agents, plugins, memory connectors, and multi-language enterprise integration. Tradeoff: Its broad integration surface suits platform teams but introduces framework concepts and version coupling. Watch for: Registering overly broad plugins grants the model unnecessary authority.}
\paragraph{MCQ-0556 [medium]} A design review reveals this mistake: Registering overly broad plugins grants the model unnecessary authority. Which mechanism should be traced first?
\begin{enumerate}[label=\Alph*.]
\item Semantic Kernel
\item AutoGen
\item PydanticAI
\item DSPy
\end{enumerate}
\noindent\textbf{Answer.} Semantic Kernel
\par\textit{Semantic Kernel offers planners, agents, plugins, memory connectors, and multi-language enterprise integration. Tradeoff: Its broad integration surface suits platform teams but introduces framework concepts and version coupling. Watch for: Registering overly broad plugins grants the model unnecessary authority.}
\paragraph{MCQ-0557 [hard]} Which causal explanation of Semantic Kernel connects what it does to when it is worth its cost?
\begin{enumerate}[label=\Alph*.]
\item Mechanism: DSPy treats prompts and modules as optimizable programs evaluated against examples and metrics. Consequence: Systematic optimization can beat manual prompting, while it depends on representative data and stable metrics.
\item Mechanism: AutoGen structures multi-agent conversations and tool-using participants for collaborative task execution. Consequence: Role decomposition enables complex workflows but multiplies calls, coordination failures, and evaluation cost.
\item Mechanism: PydanticAI uses typed dependencies, tools, results, and validation to build model-agnostic Python agents. Consequence: Type safety improves integration but cannot validate semantic truth.
\item Mechanism: Semantic Kernel offers planners, agents, plugins, memory connectors, and multi-language enterprise integration. Consequence: Its broad integration surface suits platform teams but introduces framework concepts and version coupling.
\end{enumerate}
\noindent\textbf{Answer.} Mechanism: Semantic Kernel offers planners, agents, plugins, memory connectors, and multi-language enterprise integration. Consequence: Its broad integration surface suits platform teams but introduces framework concepts and version coupling.
\par\textit{Semantic Kernel offers planners, agents, plugins, memory connectors, and multi-language enterprise integration. Tradeoff: Its broad integration surface suits platform teams but introduces framework concepts and version coupling. Watch for: Registering overly broad plugins grants the model unnecessary authority.}
\paragraph{FC-0558 [medium]} Define Semantic Kernel in interview-ready language.
\noindent\textbf{Answer.} Semantic Kernel offers planners, agents, plugins, memory connectors, and multi-language enterprise integration.
\par\textit{Semantic Kernel offers planners, agents, plugins, memory connectors, and multi-language enterprise integration.}
\paragraph{FC-0559 [hard]} State the main tradeoff and production pitfall for Semantic Kernel.
\noindent\textbf{Answer.} Tradeoff: Its broad integration surface suits platform teams but introduces framework concepts and version coupling. Pitfall: Registering overly broad plugins grants the model unnecessary authority.
\par\textit{Tradeoff: Its broad integration surface suits platform teams but introduces framework concepts and version coupling. Pitfall: Registering overly broad plugins grants the model unnecessary authority.}
\paragraph{LA-0560 [hard]} Explain Semantic Kernel from first principles, then show how you would evaluate and operate it in production.
\noindent\textbf{Answer.} Start with the mechanism: Semantic Kernel offers planners, agents, plugins, memory connectors, and multi-language enterprise integration. Make the tradeoff explicit: Its broad integration surface suits platform teams but introduces framework concepts and version coupling. Define workload-specific offline metrics and a latency/cost budget, test important slices, instrument the critical path, and create a rollback criterion. The production review must explicitly guard against this failure: Registering overly broad plugins grants the model unnecessary authority.
\par\textit{A strong answer connects mechanism, measurable tradeoffs, failure isolation, observability, and rollback rather than listing product features.}
\paragraph{MCQ-0561 [medium]} A design requires this behavior: AutoGen structures multi-agent conversations and tool-using participants for collaborative task execution. Which mechanism should the team study?
\begin{enumerate}[label=\Alph*.]
\item AutoGen
\item Haystack pipelines
\item Semantic Kernel
\item CrewAI
\end{enumerate}
\noindent\textbf{Answer.} AutoGen
\par\textit{AutoGen structures multi-agent conversations and tool-using participants for collaborative task execution. Tradeoff: Role decomposition enables complex workflows but multiplies calls, coordination failures, and evaluation cost. Watch for: Adding agents without independent capabilities creates expensive echo chambers.}
\paragraph{MCQ-0562 [hard]} The team proposes AutoGen for a real workload. Which finding gives the correct decision boundary?
\begin{enumerate}[label=\Alph*.]
\item It gives accessible orchestration but still needs explicit state, permissions, and observability.
\item Its broad integration surface suits platform teams but introduces framework concepts and version coupling.
\item Explicit components support testing, while custom graphs still require lifecycle discipline.
\item Role decomposition enables complex workflows but multiplies calls, coordination failures, and evaluation cost.
\end{enumerate}
\noindent\textbf{Answer.} Role decomposition enables complex workflows but multiplies calls, coordination failures, and evaluation cost.
\par\textit{AutoGen structures multi-agent conversations and tool-using participants for collaborative task execution. Tradeoff: Role decomposition enables complex workflows but multiplies calls, coordination failures, and evaluation cost. Watch for: Adding agents without independent capabilities creates expensive echo chambers.}
\paragraph{MCQ-0563 [hard]} Before releasing AutoGen, which failure deserves a dedicated regression test?
\begin{enumerate}[label=\Alph*.]
\item Registering overly broad plugins grants the model unnecessary authority.
\item Allowing incompatible document schemas across components creates runtime surprises.
\item Persona descriptions cannot substitute for tool-level access control.
\item Adding agents without independent capabilities creates expensive echo chambers.
\end{enumerate}
\noindent\textbf{Answer.} Adding agents without independent capabilities creates expensive echo chambers.
\par\textit{AutoGen structures multi-agent conversations and tool-using participants for collaborative task execution. Tradeoff: Role decomposition enables complex workflows but multiplies calls, coordination failures, and evaluation cost. Watch for: Adding agents without independent capabilities creates expensive echo chambers.}
\paragraph{MCQ-0564 [medium]} A design review reveals this mistake: Adding agents without independent capabilities creates expensive echo chambers. Which mechanism should be traced first?
\begin{enumerate}[label=\Alph*.]
\item Haystack pipelines
\item CrewAI
\item AutoGen
\item Semantic Kernel
\end{enumerate}
\noindent\textbf{Answer.} AutoGen
\par\textit{AutoGen structures multi-agent conversations and tool-using participants for collaborative task execution. Tradeoff: Role decomposition enables complex workflows but multiplies calls, coordination failures, and evaluation cost. Watch for: Adding agents without independent capabilities creates expensive echo chambers.}
\paragraph{MCQ-0565 [hard]} Which causal explanation of AutoGen connects what it does to when it is worth its cost?
\begin{enumerate}[label=\Alph*.]
\item Mechanism: Haystack connects typed components into retrieval, generation, routing, and agent pipelines. Consequence: Explicit components support testing, while custom graphs still require lifecycle discipline.
\item Mechanism: AutoGen structures multi-agent conversations and tool-using participants for collaborative task execution. Consequence: Role decomposition enables complex workflows but multiplies calls, coordination failures, and evaluation cost.
\item Mechanism: CrewAI organizes role-based agents, tasks, crews, and flows for sequential or coordinated work. Consequence: It gives accessible orchestration but still needs explicit state, permissions, and observability.
\item Mechanism: Semantic Kernel offers planners, agents, plugins, memory connectors, and multi-language enterprise integration. Consequence: Its broad integration surface suits platform teams but introduces framework concepts and version coupling.
\end{enumerate}
\noindent\textbf{Answer.} Mechanism: AutoGen structures multi-agent conversations and tool-using participants for collaborative task execution. Consequence: Role decomposition enables complex workflows but multiplies calls, coordination failures, and evaluation cost.
\par\textit{AutoGen structures multi-agent conversations and tool-using participants for collaborative task execution. Tradeoff: Role decomposition enables complex workflows but multiplies calls, coordination failures, and evaluation cost. Watch for: Adding agents without independent capabilities creates expensive echo chambers.}
\paragraph{FC-0566 [medium]} Define AutoGen in interview-ready language.
\noindent\textbf{Answer.} AutoGen structures multi-agent conversations and tool-using participants for collaborative task execution.
\par\textit{AutoGen structures multi-agent conversations and tool-using participants for collaborative task execution.}
\paragraph{FC-0567 [hard]} State the main tradeoff and production pitfall for AutoGen.
\noindent\textbf{Answer.} Tradeoff: Role decomposition enables complex workflows but multiplies calls, coordination failures, and evaluation cost. Pitfall: Adding agents without independent capabilities creates expensive echo chambers.
\par\textit{Tradeoff: Role decomposition enables complex workflows but multiplies calls, coordination failures, and evaluation cost. Pitfall: Adding agents without independent capabilities creates expensive echo chambers.}
\paragraph{LA-0568 [hard]} Explain AutoGen from first principles, then show how you would evaluate and operate it in production.
\noindent\textbf{Answer.} Start with the mechanism: AutoGen structures multi-agent conversations and tool-using participants for collaborative task execution. Make the tradeoff explicit: Role decomposition enables complex workflows but multiplies calls, coordination failures, and evaluation cost. Define workload-specific offline metrics and a latency/cost budget, test important slices, instrument the critical path, and create a rollback criterion. The production review must explicitly guard against this failure: Adding agents without independent capabilities creates expensive echo chambers.
\par\textit{A strong answer connects mechanism, measurable tradeoffs, failure isolation, observability, and rollback rather than listing product features.}
\paragraph{MCQ-0569 [medium]} A design requires this behavior: CrewAI organizes role-based agents, tasks, crews, and flows for sequential or coordinated work. Which mechanism should the team study?
\begin{enumerate}[label=\Alph*.]
\item DSPy
\item CrewAI
\item AutoGen
\item Airflow, Dagster, and Prefect
\end{enumerate}
\noindent\textbf{Answer.} CrewAI
\par\textit{CrewAI organizes role-based agents, tasks, crews, and flows for sequential or coordinated work. Tradeoff: It gives accessible orchestration but still needs explicit state, permissions, and observability. Watch for: Persona descriptions cannot substitute for tool-level access control.}
\paragraph{MCQ-0570 [hard]} The team proposes CrewAI for a real workload. Which finding gives the correct decision boundary?
\begin{enumerate}[label=\Alph*.]
\item It gives accessible orchestration but still needs explicit state, permissions, and observability.
\item Systematic optimization can beat manual prompting, while it depends on representative data and stable metrics.
\item They excel at batch pipelines but are not automatically low-latency interactive agent runtimes.
\item Role decomposition enables complex workflows but multiplies calls, coordination failures, and evaluation cost.
\end{enumerate}
\noindent\textbf{Answer.} It gives accessible orchestration but still needs explicit state, permissions, and observability.
\par\textit{CrewAI organizes role-based agents, tasks, crews, and flows for sequential or coordinated work. Tradeoff: It gives accessible orchestration but still needs explicit state, permissions, and observability. Watch for: Persona descriptions cannot substitute for tool-level access control.}
\paragraph{MCQ-0571 [hard]} Before releasing CrewAI, which failure deserves a dedicated regression test?
\begin{enumerate}[label=\Alph*.]
\item Optimizing against a leaky judge overfits wording instead of capability.
\item Running per-token or per-chat-turn work as heavy batch tasks adds unacceptable overhead.
\item Persona descriptions cannot substitute for tool-level access control.
\item Adding agents without independent capabilities creates expensive echo chambers.
\end{enumerate}
\noindent\textbf{Answer.} Persona descriptions cannot substitute for tool-level access control.
\par\textit{CrewAI organizes role-based agents, tasks, crews, and flows for sequential or coordinated work. Tradeoff: It gives accessible orchestration but still needs explicit state, permissions, and observability. Watch for: Persona descriptions cannot substitute for tool-level access control.}
\paragraph{MCQ-0572 [medium]} A design review reveals this mistake: Persona descriptions cannot substitute for tool-level access control. Which mechanism should be traced first?
\begin{enumerate}[label=\Alph*.]
\item CrewAI
\item DSPy
\item AutoGen
\item Airflow, Dagster, and Prefect
\end{enumerate}
\noindent\textbf{Answer.} CrewAI
\par\textit{CrewAI organizes role-based agents, tasks, crews, and flows for sequential or coordinated work. Tradeoff: It gives accessible orchestration but still needs explicit state, permissions, and observability. Watch for: Persona descriptions cannot substitute for tool-level access control.}
\paragraph{MCQ-0573 [hard]} Which causal explanation of CrewAI connects what it does to when it is worth its cost?
\begin{enumerate}[label=\Alph*.]
\item Mechanism: AutoGen structures multi-agent conversations and tool-using participants for collaborative task execution. Consequence: Role decomposition enables complex workflows but multiplies calls, coordination failures, and evaluation cost.
\item Mechanism: DSPy treats prompts and modules as optimizable programs evaluated against examples and metrics. Consequence: Systematic optimization can beat manual prompting, while it depends on representative data and stable metrics.
\item Mechanism: CrewAI organizes role-based agents, tasks, crews, and flows for sequential or coordinated work. Consequence: It gives accessible orchestration but still needs explicit state, permissions, and observability.
\item Mechanism: General-purpose data orchestrators schedule DAGs, assets, or flows with retries, lineage, and operational UIs. Consequence: They excel at batch pipelines but are not automatically low-latency interactive agent runtimes.
\end{enumerate}
\noindent\textbf{Answer.} Mechanism: CrewAI organizes role-based agents, tasks, crews, and flows for sequential or coordinated work. Consequence: It gives accessible orchestration but still needs explicit state, permissions, and observability.
\par\textit{CrewAI organizes role-based agents, tasks, crews, and flows for sequential or coordinated work. Tradeoff: It gives accessible orchestration but still needs explicit state, permissions, and observability. Watch for: Persona descriptions cannot substitute for tool-level access control.}
\paragraph{FC-0574 [medium]} Define CrewAI in interview-ready language.
\noindent\textbf{Answer.} CrewAI organizes role-based agents, tasks, crews, and flows for sequential or coordinated work.
\par\textit{CrewAI organizes role-based agents, tasks, crews, and flows for sequential or coordinated work.}
\paragraph{FC-0575 [hard]} State the main tradeoff and production pitfall for CrewAI.
\noindent\textbf{Answer.} Tradeoff: It gives accessible orchestration but still needs explicit state, permissions, and observability. Pitfall: Persona descriptions cannot substitute for tool-level access control.
\par\textit{Tradeoff: It gives accessible orchestration but still needs explicit state, permissions, and observability. Pitfall: Persona descriptions cannot substitute for tool-level access control.}
\paragraph{LA-0576 [hard]} Explain CrewAI from first principles, then show how you would evaluate and operate it in production.
\noindent\textbf{Answer.} Start with the mechanism: CrewAI organizes role-based agents, tasks, crews, and flows for sequential or coordinated work. Make the tradeoff explicit: It gives accessible orchestration but still needs explicit state, permissions, and observability. Define workload-specific offline metrics and a latency/cost budget, test important slices, instrument the critical path, and create a rollback criterion. The production review must explicitly guard against this failure: Persona descriptions cannot substitute for tool-level access control.
\par\textit{A strong answer connects mechanism, measurable tradeoffs, failure isolation, observability, and rollback rather than listing product features.}
\paragraph{MCQ-0577 [medium]} A design requires this behavior: DSPy treats prompts and modules as optimizable programs evaluated against examples and metrics. Which mechanism should the team study?
\begin{enumerate}[label=\Alph*.]
\item Airflow, Dagster, and Prefect
\item DSPy
\item LangGraph
\item AutoGen
\end{enumerate}
\noindent\textbf{Answer.} DSPy
\par\textit{DSPy treats prompts and modules as optimizable programs evaluated against examples and metrics. Tradeoff: Systematic optimization can beat manual prompting, while it depends on representative data and stable metrics. Watch for: Optimizing against a leaky judge overfits wording instead of capability.}
\paragraph{MCQ-0578 [hard]} The team proposes DSPy for a real workload. Which finding gives the correct decision boundary?
\begin{enumerate}[label=\Alph*.]
\item They excel at batch pipelines but are not automatically low-latency interactive agent runtimes.
\item Systematic optimization can beat manual prompting, while it depends on representative data and stable metrics.
\item Role decomposition enables complex workflows but multiplies calls, coordination failures, and evaluation cost.
\item Fine-grained control aids production reliability but exposes more state and graph design decisions.
\end{enumerate}
\noindent\textbf{Answer.} Systematic optimization can beat manual prompting, while it depends on representative data and stable metrics.
\par\textit{DSPy treats prompts and modules as optimizable programs evaluated against examples and metrics. Tradeoff: Systematic optimization can beat manual prompting, while it depends on representative data and stable metrics. Watch for: Optimizing against a leaky judge overfits wording instead of capability.}
\paragraph{MCQ-0579 [hard]} Before releasing DSPy, which failure deserves a dedicated regression test?
\begin{enumerate}[label=\Alph*.]
\item Running per-token or per-chat-turn work as heavy batch tasks adds unacceptable overhead.
\item Non-idempotent node side effects make checkpoint replay unsafe.
\item Optimizing against a leaky judge overfits wording instead of capability.
\item Adding agents without independent capabilities creates expensive echo chambers.
\end{enumerate}
\noindent\textbf{Answer.} Optimizing against a leaky judge overfits wording instead of capability.
\par\textit{DSPy treats prompts and modules as optimizable programs evaluated against examples and metrics. Tradeoff: Systematic optimization can beat manual prompting, while it depends on representative data and stable metrics. Watch for: Optimizing against a leaky judge overfits wording instead of capability.}
\paragraph{MCQ-0580 [medium]} A design review reveals this mistake: Optimizing against a leaky judge overfits wording instead of capability. Which mechanism should be traced first?
\begin{enumerate}[label=\Alph*.]
\item DSPy
\item LangGraph
\item Airflow, Dagster, and Prefect
\item AutoGen
\end{enumerate}
\noindent\textbf{Answer.} DSPy
\par\textit{DSPy treats prompts and modules as optimizable programs evaluated against examples and metrics. Tradeoff: Systematic optimization can beat manual prompting, while it depends on representative data and stable metrics. Watch for: Optimizing against a leaky judge overfits wording instead of capability.}
\paragraph{MCQ-0581 [hard]} Which causal explanation of DSPy connects what it does to when it is worth its cost?
\begin{enumerate}[label=\Alph*.]
\item Mechanism: General-purpose data orchestrators schedule DAGs, assets, or flows with retries, lineage, and operational UIs. Consequence: They excel at batch pipelines but are not automatically low-latency interactive agent runtimes.
\item Mechanism: DSPy treats prompts and modules as optimizable programs evaluated against examples and metrics. Consequence: Systematic optimization can beat manual prompting, while it depends on representative data and stable metrics.
\item Mechanism: LangGraph models long-running agent workflows as stateful graphs with checkpoints, interrupts, streaming, and durable execution. Consequence: Fine-grained control aids production reliability but exposes more state and graph design decisions.
\item Mechanism: AutoGen structures multi-agent conversations and tool-using participants for collaborative task execution. Consequence: Role decomposition enables complex workflows but multiplies calls, coordination failures, and evaluation cost.
\end{enumerate}
\noindent\textbf{Answer.} Mechanism: DSPy treats prompts and modules as optimizable programs evaluated against examples and metrics. Consequence: Systematic optimization can beat manual prompting, while it depends on representative data and stable metrics.
\par\textit{DSPy treats prompts and modules as optimizable programs evaluated against examples and metrics. Tradeoff: Systematic optimization can beat manual prompting, while it depends on representative data and stable metrics. Watch for: Optimizing against a leaky judge overfits wording instead of capability.}
\paragraph{FC-0582 [medium]} Define DSPy in interview-ready language.
\noindent\textbf{Answer.} DSPy treats prompts and modules as optimizable programs evaluated against examples and metrics.
\par\textit{DSPy treats prompts and modules as optimizable programs evaluated against examples and metrics.}
\paragraph{FC-0583 [hard]} State the main tradeoff and production pitfall for DSPy.
\noindent\textbf{Answer.} Tradeoff: Systematic optimization can beat manual prompting, while it depends on representative data and stable metrics. Pitfall: Optimizing against a leaky judge overfits wording instead of capability.
\par\textit{Tradeoff: Systematic optimization can beat manual prompting, while it depends on representative data and stable metrics. Pitfall: Optimizing against a leaky judge overfits wording instead of capability.}
\paragraph{LA-0584 [hard]} Explain DSPy from first principles, then show how you would evaluate and operate it in production.
\noindent\textbf{Answer.} Start with the mechanism: DSPy treats prompts and modules as optimizable programs evaluated against examples and metrics. Make the tradeoff explicit: Systematic optimization can beat manual prompting, while it depends on representative data and stable metrics. Define workload-specific offline metrics and a latency/cost budget, test important slices, instrument the critical path, and create a rollback criterion. The production review must explicitly guard against this failure: Optimizing against a leaky judge overfits wording instead of capability.
\par\textit{A strong answer connects mechanism, measurable tradeoffs, failure isolation, observability, and rollback rather than listing product features.}
\paragraph{MCQ-0585 [medium]} A design requires this behavior: Temporal persists workflow history and replays deterministic workflow code around durable activities. Which mechanism should the team study?
\begin{enumerate}[label=\Alph*.]
\item Haystack pipelines
\item Airflow, Dagster, and Prefect
\item Temporal
\item Semantic Kernel
\end{enumerate}
\noindent\textbf{Answer.} Temporal
\par\textit{Temporal persists workflow history and replays deterministic workflow code around durable activities. Tradeoff: It provides robust retries and long-running execution but requires deterministic workflow discipline. Watch for: Reading wall-clock time or random state directly in replayed code causes nondeterminism.}
\paragraph{MCQ-0586 [hard]} The team proposes Temporal for a real workload. Which finding gives the correct decision boundary?
\begin{enumerate}[label=\Alph*.]
\item Explicit components support testing, while custom graphs still require lifecycle discipline.
\item They excel at batch pipelines but are not automatically low-latency interactive agent runtimes.
\item Its broad integration surface suits platform teams but introduces framework concepts and version coupling.
\item It provides robust retries and long-running execution but requires deterministic workflow discipline.
\end{enumerate}
\noindent\textbf{Answer.} It provides robust retries and long-running execution but requires deterministic workflow discipline.
\par\textit{Temporal persists workflow history and replays deterministic workflow code around durable activities. Tradeoff: It provides robust retries and long-running execution but requires deterministic workflow discipline. Watch for: Reading wall-clock time or random state directly in replayed code causes nondeterminism.}
\paragraph{MCQ-0587 [hard]} Before releasing Temporal, which failure deserves a dedicated regression test?
\begin{enumerate}[label=\Alph*.]
\item Reading wall-clock time or random state directly in replayed code causes nondeterminism.
\item Running per-token or per-chat-turn work as heavy batch tasks adds unacceptable overhead.
\item Allowing incompatible document schemas across components creates runtime surprises.
\item Registering overly broad plugins grants the model unnecessary authority.
\end{enumerate}
\noindent\textbf{Answer.} Reading wall-clock time or random state directly in replayed code causes nondeterminism.
\par\textit{Temporal persists workflow history and replays deterministic workflow code around durable activities. Tradeoff: It provides robust retries and long-running execution but requires deterministic workflow discipline. Watch for: Reading wall-clock time or random state directly in replayed code causes nondeterminism.}
\paragraph{MCQ-0588 [medium]} A design review reveals this mistake: Reading wall-clock time or random state directly in replayed code causes nondeterminism. Which mechanism should be traced first?
\begin{enumerate}[label=\Alph*.]
\item Haystack pipelines
\item Semantic Kernel
\item Temporal
\item Airflow, Dagster, and Prefect
\end{enumerate}
\noindent\textbf{Answer.} Temporal
\par\textit{Temporal persists workflow history and replays deterministic workflow code around durable activities. Tradeoff: It provides robust retries and long-running execution but requires deterministic workflow discipline. Watch for: Reading wall-clock time or random state directly in replayed code causes nondeterminism.}
\paragraph{MCQ-0589 [hard]} Which causal explanation of Temporal connects what it does to when it is worth its cost?
\begin{enumerate}[label=\Alph*.]
\item Mechanism: Temporal persists workflow history and replays deterministic workflow code around durable activities. Consequence: It provides robust retries and long-running execution but requires deterministic workflow discipline.
\item Mechanism: Semantic Kernel offers planners, agents, plugins, memory connectors, and multi-language enterprise integration. Consequence: Its broad integration surface suits platform teams but introduces framework concepts and version coupling.
\item Mechanism: Haystack connects typed components into retrieval, generation, routing, and agent pipelines. Consequence: Explicit components support testing, while custom graphs still require lifecycle discipline.
\item Mechanism: General-purpose data orchestrators schedule DAGs, assets, or flows with retries, lineage, and operational UIs. Consequence: They excel at batch pipelines but are not automatically low-latency interactive agent runtimes.
\end{enumerate}
\noindent\textbf{Answer.} Mechanism: Temporal persists workflow history and replays deterministic workflow code around durable activities. Consequence: It provides robust retries and long-running execution but requires deterministic workflow discipline.
\par\textit{Temporal persists workflow history and replays deterministic workflow code around durable activities. Tradeoff: It provides robust retries and long-running execution but requires deterministic workflow discipline. Watch for: Reading wall-clock time or random state directly in replayed code causes nondeterminism.}
\paragraph{FC-0590 [medium]} Define Temporal in interview-ready language.
\noindent\textbf{Answer.} Temporal persists workflow history and replays deterministic workflow code around durable activities.
\par\textit{Temporal persists workflow history and replays deterministic workflow code around durable activities.}
\paragraph{FC-0591 [hard]} State the main tradeoff and production pitfall for Temporal.
\noindent\textbf{Answer.} Tradeoff: It provides robust retries and long-running execution but requires deterministic workflow discipline. Pitfall: Reading wall-clock time or random state directly in replayed code causes nondeterminism.
\par\textit{Tradeoff: It provides robust retries and long-running execution but requires deterministic workflow discipline. Pitfall: Reading wall-clock time or random state directly in replayed code causes nondeterminism.}
\paragraph{LA-0592 [hard]} Explain Temporal from first principles, then show how you would evaluate and operate it in production.
\noindent\textbf{Answer.} Start with the mechanism: Temporal persists workflow history and replays deterministic workflow code around durable activities. Make the tradeoff explicit: It provides robust retries and long-running execution but requires deterministic workflow discipline. Define workload-specific offline metrics and a latency/cost budget, test important slices, instrument the critical path, and create a rollback criterion. The production review must explicitly guard against this failure: Reading wall-clock time or random state directly in replayed code causes nondeterminism.
\par\textit{A strong answer connects mechanism, measurable tradeoffs, failure isolation, observability, and rollback rather than listing product features.}
\paragraph{MCQ-0593 [medium]} A design requires this behavior: General-purpose data orchestrators schedule DAGs, assets, or flows with retries, lineage, and operational UIs. Which mechanism should the team study?
\begin{enumerate}[label=\Alph*.]
\item Temporal
\item Airflow, Dagster, and Prefect
\item LangGraph
\item LlamaIndex workflows
\end{enumerate}
\noindent\textbf{Answer.} Airflow, Dagster, and Prefect
\par\textit{General-purpose data orchestrators schedule DAGs, assets, or flows with retries, lineage, and operational UIs. Tradeoff: They excel at batch pipelines but are not automatically low-latency interactive agent runtimes. Watch for: Running per-token or per-chat-turn work as heavy batch tasks adds unacceptable overhead.}
\paragraph{MCQ-0594 [hard]} The team proposes Airflow, Dagster, and Prefect for a real workload. Which finding gives the correct decision boundary?
\begin{enumerate}[label=\Alph*.]
\item It provides robust retries and long-running execution but requires deterministic workflow discipline.
\item Fine-grained control aids production reliability but exposes more state and graph design decisions.
\item Strong data and indexing integrations speed RAG work but can couple application flow to framework abstractions.
\item They excel at batch pipelines but are not automatically low-latency interactive agent runtimes.
\end{enumerate}
\noindent\textbf{Answer.} They excel at batch pipelines but are not automatically low-latency interactive agent runtimes.
\par\textit{General-purpose data orchestrators schedule DAGs, assets, or flows with retries, lineage, and operational UIs. Tradeoff: They excel at batch pipelines but are not automatically low-latency interactive agent runtimes. Watch for: Running per-token or per-chat-turn work as heavy batch tasks adds unacceptable overhead.}
\paragraph{MCQ-0595 [hard]} Before releasing Airflow, Dagster, and Prefect, which failure deserves a dedicated regression test?
\begin{enumerate}[label=\Alph*.]
\item Running per-token or per-chat-turn work as heavy batch tasks adds unacceptable overhead.
\item Non-idempotent node side effects make checkpoint replay unsafe.
\item Hiding all retrieval behind defaults makes evaluation and migration difficult.
\item Reading wall-clock time or random state directly in replayed code causes nondeterminism.
\end{enumerate}
\noindent\textbf{Answer.} Running per-token or per-chat-turn work as heavy batch tasks adds unacceptable overhead.
\par\textit{General-purpose data orchestrators schedule DAGs, assets, or flows with retries, lineage, and operational UIs. Tradeoff: They excel at batch pipelines but are not automatically low-latency interactive agent runtimes. Watch for: Running per-token or per-chat-turn work as heavy batch tasks adds unacceptable overhead.}
\paragraph{MCQ-0596 [medium]} A design review reveals this mistake: Running per-token or per-chat-turn work as heavy batch tasks adds unacceptable overhead. Which mechanism should be traced first?
\begin{enumerate}[label=\Alph*.]
\item LangGraph
\item Airflow, Dagster, and Prefect
\item Temporal
\item LlamaIndex workflows
\end{enumerate}
\noindent\textbf{Answer.} Airflow, Dagster, and Prefect
\par\textit{General-purpose data orchestrators schedule DAGs, assets, or flows with retries, lineage, and operational UIs. Tradeoff: They excel at batch pipelines but are not automatically low-latency interactive agent runtimes. Watch for: Running per-token or per-chat-turn work as heavy batch tasks adds unacceptable overhead.}
\paragraph{MCQ-0597 [hard]} Which causal explanation of Airflow, Dagster, and Prefect connects what it does to when it is worth its cost?
\begin{enumerate}[label=\Alph*.]
\item Mechanism: Event-driven workflows coordinate retrieval, tools, agents, and structured steps in the LlamaIndex ecosystem. Consequence: Strong data and indexing integrations speed RAG work but can couple application flow to framework abstractions.
\item Mechanism: LangGraph models long-running agent workflows as stateful graphs with checkpoints, interrupts, streaming, and durable execution. Consequence: Fine-grained control aids production reliability but exposes more state and graph design decisions.
\item Mechanism: Temporal persists workflow history and replays deterministic workflow code around durable activities. Consequence: It provides robust retries and long-running execution but requires deterministic workflow discipline.
\item Mechanism: General-purpose data orchestrators schedule DAGs, assets, or flows with retries, lineage, and operational UIs. Consequence: They excel at batch pipelines but are not automatically low-latency interactive agent runtimes.
\end{enumerate}
\noindent\textbf{Answer.} Mechanism: General-purpose data orchestrators schedule DAGs, assets, or flows with retries, lineage, and operational UIs. Consequence: They excel at batch pipelines but are not automatically low-latency interactive agent runtimes.
\par\textit{General-purpose data orchestrators schedule DAGs, assets, or flows with retries, lineage, and operational UIs. Tradeoff: They excel at batch pipelines but are not automatically low-latency interactive agent runtimes. Watch for: Running per-token or per-chat-turn work as heavy batch tasks adds unacceptable overhead.}
\paragraph{FC-0598 [medium]} Define Airflow, Dagster, and Prefect in interview-ready language.
\noindent\textbf{Answer.} General-purpose data orchestrators schedule DAGs, assets, or flows with retries, lineage, and operational UIs.
\par\textit{General-purpose data orchestrators schedule DAGs, assets, or flows with retries, lineage, and operational UIs.}
\paragraph{FC-0599 [hard]} State the main tradeoff and production pitfall for Airflow, Dagster, and Prefect.
\noindent\textbf{Answer.} Tradeoff: They excel at batch pipelines but are not automatically low-latency interactive agent runtimes. Pitfall: Running per-token or per-chat-turn work as heavy batch tasks adds unacceptable overhead.
\par\textit{Tradeoff: They excel at batch pipelines but are not automatically low-latency interactive agent runtimes. Pitfall: Running per-token or per-chat-turn work as heavy batch tasks adds unacceptable overhead.}
\paragraph{LA-0600 [hard]} Explain Airflow, Dagster, and Prefect from first principles, then show how you would evaluate and operate it in production.
\noindent\textbf{Answer.} Start with the mechanism: General-purpose data orchestrators schedule DAGs, assets, or flows with retries, lineage, and operational UIs. Make the tradeoff explicit: They excel at batch pipelines but are not automatically low-latency interactive agent runtimes. Define workload-specific offline metrics and a latency/cost budget, test important slices, instrument the critical path, and create a rollback criterion. The production review must explicitly guard against this failure: Running per-token or per-chat-turn work as heavy batch tasks adds unacceptable overhead.
\par\textit{A strong answer connects mechanism, measurable tradeoffs, failure isolation, observability, and rollback rather than listing product features.}
\subsection{Observability and Monitoring}
\paragraph{MCQ-0601 [medium]} A design requires this behavior: OpenTelemetry provides vendor-neutral APIs, SDKs, collectors, and semantic conventions for traces, metrics, logs, and profiles. Which mechanism should the team study?
\begin{enumerate}[label=\Alph*.]
\item OpenTelemetry
\item LLM tracing
\item LangSmith
\item Prompt and model drift
\end{enumerate}
\noindent\textbf{Answer.} OpenTelemetry
\par\textit{OpenTelemetry provides vendor-neutral APIs, SDKs, collectors, and semantic conventions for traces, metrics, logs, and profiles. Tradeoff: Portability improves, while instrumentation quality and backend cost remain engineering work. Watch for: High-cardinality prompt or user attributes can overwhelm telemetry systems.}
\paragraph{MCQ-0602 [hard]} The team proposes OpenTelemetry for a real workload. Which finding gives the correct decision boundary?
\begin{enumerate}[label=\Alph*.]
\item Tight LangChain integration accelerates work but requires portability planning.
\item Debugging improves, while content capture raises privacy and retention risk.
\item Alerts catch regressions but thresholds must distinguish seasonality from harm.
\item Portability improves, while instrumentation quality and backend cost remain engineering work.
\end{enumerate}
\noindent\textbf{Answer.} Portability improves, while instrumentation quality and backend cost remain engineering work.
\par\textit{OpenTelemetry provides vendor-neutral APIs, SDKs, collectors, and semantic conventions for traces, metrics, logs, and profiles. Tradeoff: Portability improves, while instrumentation quality and backend cost remain engineering work. Watch for: High-cardinality prompt or user attributes can overwhelm telemetry systems.}
\paragraph{MCQ-0603 [hard]} Before releasing OpenTelemetry, which failure deserves a dedicated regression test?
\begin{enumerate}[label=\Alph*.]
\item Logging secrets or personal data by default creates a second sensitive datastore.
\item Depending only on vendor UI annotations leaves evaluation logic unreproducible.
\item Measuring embedding distribution shift alone does not prove answer quality changed.
\item High-cardinality prompt or user attributes can overwhelm telemetry systems.
\end{enumerate}
\noindent\textbf{Answer.} High-cardinality prompt or user attributes can overwhelm telemetry systems.
\par\textit{OpenTelemetry provides vendor-neutral APIs, SDKs, collectors, and semantic conventions for traces, metrics, logs, and profiles. Tradeoff: Portability improves, while instrumentation quality and backend cost remain engineering work. Watch for: High-cardinality prompt or user attributes can overwhelm telemetry systems.}
\paragraph{MCQ-0604 [medium]} A design review reveals this mistake: High-cardinality prompt or user attributes can overwhelm telemetry systems. Which mechanism should be traced first?
\begin{enumerate}[label=\Alph*.]
\item Prompt and model drift
\item LLM tracing
\item OpenTelemetry
\item LangSmith
\end{enumerate}
\noindent\textbf{Answer.} OpenTelemetry
\par\textit{OpenTelemetry provides vendor-neutral APIs, SDKs, collectors, and semantic conventions for traces, metrics, logs, and profiles. Tradeoff: Portability improves, while instrumentation quality and backend cost remain engineering work. Watch for: High-cardinality prompt or user attributes can overwhelm telemetry systems.}
\paragraph{MCQ-0605 [hard]} Which causal explanation of OpenTelemetry connects what it does to when it is worth its cost?
\begin{enumerate}[label=\Alph*.]
\item Mechanism: OpenTelemetry provides vendor-neutral APIs, SDKs, collectors, and semantic conventions for traces, metrics, logs, and profiles. Consequence: Portability improves, while instrumentation quality and backend cost remain engineering work.
\item Mechanism: LLM traces capture prompts, outputs, model settings, token usage, latency, tool calls, retrieval, and evaluations. Consequence: Debugging improves, while content capture raises privacy and retention risk.
\item Mechanism: Drift monitoring detects changes in inputs, retrieved evidence, outputs, model versions, and user behavior. Consequence: Alerts catch regressions but thresholds must distinguish seasonality from harm.
\item Mechanism: LangSmith provides tracing, datasets, evaluation, prompt management, and deployment-oriented tooling for LLM applications. Consequence: Tight LangChain integration accelerates work but requires portability planning.
\end{enumerate}
\noindent\textbf{Answer.} Mechanism: OpenTelemetry provides vendor-neutral APIs, SDKs, collectors, and semantic conventions for traces, metrics, logs, and profiles. Consequence: Portability improves, while instrumentation quality and backend cost remain engineering work.
\par\textit{OpenTelemetry provides vendor-neutral APIs, SDKs, collectors, and semantic conventions for traces, metrics, logs, and profiles. Tradeoff: Portability improves, while instrumentation quality and backend cost remain engineering work. Watch for: High-cardinality prompt or user attributes can overwhelm telemetry systems.}
\paragraph{FC-0606 [medium]} Define OpenTelemetry in interview-ready language.
\noindent\textbf{Answer.} OpenTelemetry provides vendor-neutral APIs, SDKs, collectors, and semantic conventions for traces, metrics, logs, and profiles.
\par\textit{OpenTelemetry provides vendor-neutral APIs, SDKs, collectors, and semantic conventions for traces, metrics, logs, and profiles.}
\paragraph{FC-0607 [hard]} State the main tradeoff and production pitfall for OpenTelemetry.
\noindent\textbf{Answer.} Tradeoff: Portability improves, while instrumentation quality and backend cost remain engineering work. Pitfall: High-cardinality prompt or user attributes can overwhelm telemetry systems.
\par\textit{Tradeoff: Portability improves, while instrumentation quality and backend cost remain engineering work. Pitfall: High-cardinality prompt or user attributes can overwhelm telemetry systems.}
\paragraph{LA-0608 [hard]} Explain OpenTelemetry from first principles, then show how you would evaluate and operate it in production.
\noindent\textbf{Answer.} Start with the mechanism: OpenTelemetry provides vendor-neutral APIs, SDKs, collectors, and semantic conventions for traces, metrics, logs, and profiles. Make the tradeoff explicit: Portability improves, while instrumentation quality and backend cost remain engineering work. Define workload-specific offline metrics and a latency/cost budget, test important slices, instrument the critical path, and create a rollback criterion. The production review must explicitly guard against this failure: High-cardinality prompt or user attributes can overwhelm telemetry systems.
\par\textit{A strong answer connects mechanism, measurable tradeoffs, failure isolation, observability, and rollback rather than listing product features.}
\paragraph{MCQ-0609 [medium]} A design requires this behavior: Traces connect request spans across gateways, retrievers, models, tools, queues, and databases. Which mechanism should the team study?
\begin{enumerate}[label=\Alph*.]
\item LLM tracing
\item OpenTelemetry
\item Distributed tracing
\item MLflow tracing
\end{enumerate}
\noindent\textbf{Answer.} Distributed tracing
\par\textit{Traces connect request spans across gateways, retrievers, models, tools, queues, and databases. Tradeoff: They reveal critical paths but sampling can hide rare quality failures. Watch for: Reusing trace identifiers across independent user requests breaks isolation.}
\paragraph{MCQ-0610 [hard]} The team proposes Distributed tracing for a real workload. Which finding gives the correct decision boundary?
\begin{enumerate}[label=\Alph*.]
\item Debugging improves, while content capture raises privacy and retention risk.
\item They reveal critical paths but sampling can hide rare quality failures.
\item Portability improves, while instrumentation quality and backend cost remain engineering work.
\item A common ML platform improves lineage but may be heavier than a focused tracing tool.
\end{enumerate}
\noindent\textbf{Answer.} They reveal critical paths but sampling can hide rare quality failures.
\par\textit{Traces connect request spans across gateways, retrievers, models, tools, queues, and databases. Tradeoff: They reveal critical paths but sampling can hide rare quality failures. Watch for: Reusing trace identifiers across independent user requests breaks isolation.}
\paragraph{MCQ-0611 [hard]} Before releasing Distributed tracing, which failure deserves a dedicated regression test?
\begin{enumerate}[label=\Alph*.]
\item Mixing incomparable scorer versions in one trend chart produces false regressions.
\item Logging secrets or personal data by default creates a second sensitive datastore.
\item Reusing trace identifiers across independent user requests breaks isolation.
\item High-cardinality prompt or user attributes can overwhelm telemetry systems.
\end{enumerate}
\noindent\textbf{Answer.} Reusing trace identifiers across independent user requests breaks isolation.
\par\textit{Traces connect request spans across gateways, retrievers, models, tools, queues, and databases. Tradeoff: They reveal critical paths but sampling can hide rare quality failures. Watch for: Reusing trace identifiers across independent user requests breaks isolation.}
\paragraph{MCQ-0612 [medium]} A design review reveals this mistake: Reusing trace identifiers across independent user requests breaks isolation. Which mechanism should be traced first?
\begin{enumerate}[label=\Alph*.]
\item OpenTelemetry
\item LLM tracing
\item MLflow tracing
\item Distributed tracing
\end{enumerate}
\noindent\textbf{Answer.} Distributed tracing
\par\textit{Traces connect request spans across gateways, retrievers, models, tools, queues, and databases. Tradeoff: They reveal critical paths but sampling can hide rare quality failures. Watch for: Reusing trace identifiers across independent user requests breaks isolation.}
\paragraph{MCQ-0613 [hard]} Which causal explanation of Distributed tracing connects what it does to when it is worth its cost?
\begin{enumerate}[label=\Alph*.]
\item Mechanism: OpenTelemetry provides vendor-neutral APIs, SDKs, collectors, and semantic conventions for traces, metrics, logs, and profiles. Consequence: Portability improves, while instrumentation quality and backend cost remain engineering work.
\item Mechanism: MLflow connects traces, evaluation datasets, scorers, experiments, and production monitoring. Consequence: A common ML platform improves lineage but may be heavier than a focused tracing tool.
\item Mechanism: Traces connect request spans across gateways, retrievers, models, tools, queues, and databases. Consequence: They reveal critical paths but sampling can hide rare quality failures.
\item Mechanism: LLM traces capture prompts, outputs, model settings, token usage, latency, tool calls, retrieval, and evaluations. Consequence: Debugging improves, while content capture raises privacy and retention risk.
\end{enumerate}
\noindent\textbf{Answer.} Mechanism: Traces connect request spans across gateways, retrievers, models, tools, queues, and databases. Consequence: They reveal critical paths but sampling can hide rare quality failures.
\par\textit{Traces connect request spans across gateways, retrievers, models, tools, queues, and databases. Tradeoff: They reveal critical paths but sampling can hide rare quality failures. Watch for: Reusing trace identifiers across independent user requests breaks isolation.}
\paragraph{FC-0614 [medium]} Define Distributed tracing in interview-ready language.
\noindent\textbf{Answer.} Traces connect request spans across gateways, retrievers, models, tools, queues, and databases.
\par\textit{Traces connect request spans across gateways, retrievers, models, tools, queues, and databases.}
\paragraph{FC-0615 [hard]} State the main tradeoff and production pitfall for Distributed tracing.
\noindent\textbf{Answer.} Tradeoff: They reveal critical paths but sampling can hide rare quality failures. Pitfall: Reusing trace identifiers across independent user requests breaks isolation.
\par\textit{Tradeoff: They reveal critical paths but sampling can hide rare quality failures. Pitfall: Reusing trace identifiers across independent user requests breaks isolation.}
\paragraph{LA-0616 [hard]} Explain Distributed tracing from first principles, then show how you would evaluate and operate it in production.
\noindent\textbf{Answer.} Start with the mechanism: Traces connect request spans across gateways, retrievers, models, tools, queues, and databases. Make the tradeoff explicit: They reveal critical paths but sampling can hide rare quality failures. Define workload-specific offline metrics and a latency/cost budget, test important slices, instrument the critical path, and create a rollback criterion. The production review must explicitly guard against this failure: Reusing trace identifiers across independent user requests breaks isolation.
\par\textit{A strong answer connects mechanism, measurable tradeoffs, failure isolation, observability, and rollback rather than listing product features.}
\paragraph{MCQ-0617 [medium]} A design requires this behavior: LLM traces capture prompts, outputs, model settings, token usage, latency, tool calls, retrieval, and evaluations. Which mechanism should the team study?
\begin{enumerate}[label=\Alph*.]
\item LLM tracing
\item Prompt and model drift
\item Arize Phoenix
\item Evaluation in production
\end{enumerate}
\noindent\textbf{Answer.} LLM tracing
\par\textit{LLM traces capture prompts, outputs, model settings, token usage, latency, tool calls, retrieval, and evaluations. Tradeoff: Debugging improves, while content capture raises privacy and retention risk. Watch for: Logging secrets or personal data by default creates a second sensitive datastore.}
\paragraph{MCQ-0618 [hard]} The team proposes LLM tracing for a real workload. Which finding gives the correct decision boundary?
\begin{enumerate}[label=\Alph*.]
\item Live coverage improves realism but evaluations must not expose or delay users.
\item Alerts catch regressions but thresholds must distinguish seasonality from harm.
\item Debugging improves, while content capture raises privacy and retention risk.
\item It unifies qualitative traces and quantitative evals, while production operations still need design.
\end{enumerate}
\noindent\textbf{Answer.} Debugging improves, while content capture raises privacy and retention risk.
\par\textit{LLM traces capture prompts, outputs, model settings, token usage, latency, tool calls, retrieval, and evaluations. Tradeoff: Debugging improves, while content capture raises privacy and retention risk. Watch for: Logging secrets or personal data by default creates a second sensitive datastore.}
\paragraph{MCQ-0619 [hard]} Before releasing LLM tracing, which failure deserves a dedicated regression test?
\begin{enumerate}[label=\Alph*.]
\item Judge scores without slice analysis can mask failures for important cohorts.
\item Treating implicit clicks as unambiguous quality labels creates bias.
\item Measuring embedding distribution shift alone does not prove answer quality changed.
\item Logging secrets or personal data by default creates a second sensitive datastore.
\end{enumerate}
\noindent\textbf{Answer.} Logging secrets or personal data by default creates a second sensitive datastore.
\par\textit{LLM traces capture prompts, outputs, model settings, token usage, latency, tool calls, retrieval, and evaluations. Tradeoff: Debugging improves, while content capture raises privacy and retention risk. Watch for: Logging secrets or personal data by default creates a second sensitive datastore.}
\paragraph{MCQ-0620 [medium]} A design review reveals this mistake: Logging secrets or personal data by default creates a second sensitive datastore. Which mechanism should be traced first?
\begin{enumerate}[label=\Alph*.]
\item LLM tracing
\item Evaluation in production
\item Arize Phoenix
\item Prompt and model drift
\end{enumerate}
\noindent\textbf{Answer.} LLM tracing
\par\textit{LLM traces capture prompts, outputs, model settings, token usage, latency, tool calls, retrieval, and evaluations. Tradeoff: Debugging improves, while content capture raises privacy and retention risk. Watch for: Logging secrets or personal data by default creates a second sensitive datastore.}
\paragraph{MCQ-0621 [hard]} Which causal explanation of LLM tracing connects what it does to when it is worth its cost?
\begin{enumerate}[label=\Alph*.]
\item Mechanism: Phoenix provides open-source tracing and evaluation for LLM, retrieval, and agent applications with OpenTelemetry support. Consequence: It unifies qualitative traces and quantitative evals, while production operations still need design.
\item Mechanism: Drift monitoring detects changes in inputs, retrieved evidence, outputs, model versions, and user behavior. Consequence: Alerts catch regressions but thresholds must distinguish seasonality from harm.
\item Mechanism: Online evaluators combine sampled judges, deterministic checks, human feedback, and business outcomes. Consequence: Live coverage improves realism but evaluations must not expose or delay users.
\item Mechanism: LLM traces capture prompts, outputs, model settings, token usage, latency, tool calls, retrieval, and evaluations. Consequence: Debugging improves, while content capture raises privacy and retention risk.
\end{enumerate}
\noindent\textbf{Answer.} Mechanism: LLM traces capture prompts, outputs, model settings, token usage, latency, tool calls, retrieval, and evaluations. Consequence: Debugging improves, while content capture raises privacy and retention risk.
\par\textit{LLM traces capture prompts, outputs, model settings, token usage, latency, tool calls, retrieval, and evaluations. Tradeoff: Debugging improves, while content capture raises privacy and retention risk. Watch for: Logging secrets or personal data by default creates a second sensitive datastore.}
\paragraph{FC-0622 [medium]} Define LLM tracing in interview-ready language.
\noindent\textbf{Answer.} LLM traces capture prompts, outputs, model settings, token usage, latency, tool calls, retrieval, and evaluations.
\par\textit{LLM traces capture prompts, outputs, model settings, token usage, latency, tool calls, retrieval, and evaluations.}
\paragraph{FC-0623 [hard]} State the main tradeoff and production pitfall for LLM tracing.
\noindent\textbf{Answer.} Tradeoff: Debugging improves, while content capture raises privacy and retention risk. Pitfall: Logging secrets or personal data by default creates a second sensitive datastore.
\par\textit{Tradeoff: Debugging improves, while content capture raises privacy and retention risk. Pitfall: Logging secrets or personal data by default creates a second sensitive datastore.}
\paragraph{LA-0624 [hard]} Explain LLM tracing from first principles, then show how you would evaluate and operate it in production.
\noindent\textbf{Answer.} Start with the mechanism: LLM traces capture prompts, outputs, model settings, token usage, latency, tool calls, retrieval, and evaluations. Make the tradeoff explicit: Debugging improves, while content capture raises privacy and retention risk. Define workload-specific offline metrics and a latency/cost budget, test important slices, instrument the critical path, and create a rollback criterion. The production review must explicitly guard against this failure: Logging secrets or personal data by default creates a second sensitive datastore.
\par\textit{A strong answer connects mechanism, measurable tradeoffs, failure isolation, observability, and rollback rather than listing product features.}
\paragraph{MCQ-0625 [medium]} A design requires this behavior: Counters, histograms, gauges, service-level indicators, and objectives quantify reliability and performance. Which mechanism should the team study?
\begin{enumerate}[label=\Alph*.]
\item Structured logging
\item Metrics and SLOs
\item Arize Phoenix
\item Prompt and model drift
\end{enumerate}
\noindent\textbf{Answer.} Metrics and SLOs
\par\textit{Counters, histograms, gauges, service-level indicators, and objectives quantify reliability and performance. Tradeoff: Aggregate signals scale well but cannot explain individual semantic failures. Watch for: Averaging latency hides long-tail p95 and p99 user pain.}
\paragraph{MCQ-0626 [hard]} The team proposes Metrics and SLOs for a real workload. Which finding gives the correct decision boundary?
\begin{enumerate}[label=\Alph*.]
\item Aggregate signals scale well but cannot explain individual semantic failures.
\item It unifies qualitative traces and quantitative evals, while production operations still need design.
\item Search and aggregation improve, while schemas must evolve consistently.
\item Alerts catch regressions but thresholds must distinguish seasonality from harm.
\end{enumerate}
\noindent\textbf{Answer.} Aggregate signals scale well but cannot explain individual semantic failures.
\par\textit{Counters, histograms, gauges, service-level indicators, and objectives quantify reliability and performance. Tradeoff: Aggregate signals scale well but cannot explain individual semantic failures. Watch for: Averaging latency hides long-tail p95 and p99 user pain.}
\paragraph{MCQ-0627 [hard]} Before releasing Metrics and SLOs, which failure deserves a dedicated regression test?
\begin{enumerate}[label=\Alph*.]
\item Judge scores without slice analysis can mask failures for important cohorts.
\item Averaging latency hides long-tail p95 and p99 user pain.
\item Measuring embedding distribution shift alone does not prove answer quality changed.
\item Placing unbounded model output in one log field causes cost and ingestion failures.
\end{enumerate}
\noindent\textbf{Answer.} Averaging latency hides long-tail p95 and p99 user pain.
\par\textit{Counters, histograms, gauges, service-level indicators, and objectives quantify reliability and performance. Tradeoff: Aggregate signals scale well but cannot explain individual semantic failures. Watch for: Averaging latency hides long-tail p95 and p99 user pain.}
\paragraph{MCQ-0628 [medium]} A design review reveals this mistake: Averaging latency hides long-tail p95 and p99 user pain. Which mechanism should be traced first?
\begin{enumerate}[label=\Alph*.]
\item Prompt and model drift
\item Structured logging
\item Arize Phoenix
\item Metrics and SLOs
\end{enumerate}
\noindent\textbf{Answer.} Metrics and SLOs
\par\textit{Counters, histograms, gauges, service-level indicators, and objectives quantify reliability and performance. Tradeoff: Aggregate signals scale well but cannot explain individual semantic failures. Watch for: Averaging latency hides long-tail p95 and p99 user pain.}
\paragraph{MCQ-0629 [hard]} Which causal explanation of Metrics and SLOs connects what it does to when it is worth its cost?
\begin{enumerate}[label=\Alph*.]
\item Mechanism: Counters, histograms, gauges, service-level indicators, and objectives quantify reliability and performance. Consequence: Aggregate signals scale well but cannot explain individual semantic failures.
\item Mechanism: Phoenix provides open-source tracing and evaluation for LLM, retrieval, and agent applications with OpenTelemetry support. Consequence: It unifies qualitative traces and quantitative evals, while production operations still need design.
\item Mechanism: Structured events use stable fields, severity, timestamps, correlation IDs, and redaction. Consequence: Search and aggregation improve, while schemas must evolve consistently.
\item Mechanism: Drift monitoring detects changes in inputs, retrieved evidence, outputs, model versions, and user behavior. Consequence: Alerts catch regressions but thresholds must distinguish seasonality from harm.
\end{enumerate}
\noindent\textbf{Answer.} Mechanism: Counters, histograms, gauges, service-level indicators, and objectives quantify reliability and performance. Consequence: Aggregate signals scale well but cannot explain individual semantic failures.
\par\textit{Counters, histograms, gauges, service-level indicators, and objectives quantify reliability and performance. Tradeoff: Aggregate signals scale well but cannot explain individual semantic failures. Watch for: Averaging latency hides long-tail p95 and p99 user pain.}
\paragraph{FC-0630 [medium]} Define Metrics and SLOs in interview-ready language.
\noindent\textbf{Answer.} Counters, histograms, gauges, service-level indicators, and objectives quantify reliability and performance.
\par\textit{Counters, histograms, gauges, service-level indicators, and objectives quantify reliability and performance.}
\paragraph{FC-0631 [hard]} State the main tradeoff and production pitfall for Metrics and SLOs.
\noindent\textbf{Answer.} Tradeoff: Aggregate signals scale well but cannot explain individual semantic failures. Pitfall: Averaging latency hides long-tail p95 and p99 user pain.
\par\textit{Tradeoff: Aggregate signals scale well but cannot explain individual semantic failures. Pitfall: Averaging latency hides long-tail p95 and p99 user pain.}
\paragraph{LA-0632 [hard]} Explain Metrics and SLOs from first principles, then show how you would evaluate and operate it in production.
\noindent\textbf{Answer.} Start with the mechanism: Counters, histograms, gauges, service-level indicators, and objectives quantify reliability and performance. Make the tradeoff explicit: Aggregate signals scale well but cannot explain individual semantic failures. Define workload-specific offline metrics and a latency/cost budget, test important slices, instrument the critical path, and create a rollback criterion. The production review must explicitly guard against this failure: Averaging latency hides long-tail p95 and p99 user pain.
\par\textit{A strong answer connects mechanism, measurable tradeoffs, failure isolation, observability, and rollback rather than listing product features.}
\paragraph{MCQ-0633 [medium]} A design requires this behavior: Structured events use stable fields, severity, timestamps, correlation IDs, and redaction. Which mechanism should the team study?
\begin{enumerate}[label=\Alph*.]
\item LLM tracing
\item Cost observability
\item Structured logging
\item Evaluation in production
\end{enumerate}
\noindent\textbf{Answer.} Structured logging
\par\textit{Structured events use stable fields, severity, timestamps, correlation IDs, and redaction. Tradeoff: Search and aggregation improve, while schemas must evolve consistently. Watch for: Placing unbounded model output in one log field causes cost and ingestion failures.}
\paragraph{MCQ-0634 [hard]} The team proposes Structured logging for a real workload. Which finding gives the correct decision boundary?
\begin{enumerate}[label=\Alph*.]
\item Live coverage improves realism but evaluations must not expose or delay users.
\item Debugging improves, while content capture raises privacy and retention risk.
\item Search and aggregation improve, while schemas must evolve consistently.
\item Unit economics become actionable, while provider pricing and caching rules evolve.
\end{enumerate}
\noindent\textbf{Answer.} Search and aggregation improve, while schemas must evolve consistently.
\par\textit{Structured events use stable fields, severity, timestamps, correlation IDs, and redaction. Tradeoff: Search and aggregation improve, while schemas must evolve consistently. Watch for: Placing unbounded model output in one log field causes cost and ingestion failures.}
\paragraph{MCQ-0635 [hard]} Before releasing Structured logging, which failure deserves a dedicated regression test?
\begin{enumerate}[label=\Alph*.]
\item Logging secrets or personal data by default creates a second sensitive datastore.
\item Placing unbounded model output in one log field causes cost and ingestion failures.
\item Tracking only average cost hides pathological long-running agents.
\item Treating implicit clicks as unambiguous quality labels creates bias.
\end{enumerate}
\noindent\textbf{Answer.} Placing unbounded model output in one log field causes cost and ingestion failures.
\par\textit{Structured events use stable fields, severity, timestamps, correlation IDs, and redaction. Tradeoff: Search and aggregation improve, while schemas must evolve consistently. Watch for: Placing unbounded model output in one log field causes cost and ingestion failures.}
\paragraph{MCQ-0636 [medium]} A design review reveals this mistake: Placing unbounded model output in one log field causes cost and ingestion failures. Which mechanism should be traced first?
\begin{enumerate}[label=\Alph*.]
\item Structured logging
\item Evaluation in production
\item Cost observability
\item LLM tracing
\end{enumerate}
\noindent\textbf{Answer.} Structured logging
\par\textit{Structured events use stable fields, severity, timestamps, correlation IDs, and redaction. Tradeoff: Search and aggregation improve, while schemas must evolve consistently. Watch for: Placing unbounded model output in one log field causes cost and ingestion failures.}
\paragraph{MCQ-0637 [hard]} Which causal explanation of Structured logging connects what it does to when it is worth its cost?
\begin{enumerate}[label=\Alph*.]
\item Mechanism: Structured events use stable fields, severity, timestamps, correlation IDs, and redaction. Consequence: Search and aggregation improve, while schemas must evolve consistently.
\item Mechanism: LLM traces capture prompts, outputs, model settings, token usage, latency, tool calls, retrieval, and evaluations. Consequence: Debugging improves, while content capture raises privacy and retention risk.
\item Mechanism: Cost telemetry attributes input, output, cache, retrieval, tool, and infrastructure spend to requests and tenants. Consequence: Unit economics become actionable, while provider pricing and caching rules evolve.
\item Mechanism: Online evaluators combine sampled judges, deterministic checks, human feedback, and business outcomes. Consequence: Live coverage improves realism but evaluations must not expose or delay users.
\end{enumerate}
\noindent\textbf{Answer.} Mechanism: Structured events use stable fields, severity, timestamps, correlation IDs, and redaction. Consequence: Search and aggregation improve, while schemas must evolve consistently.
\par\textit{Structured events use stable fields, severity, timestamps, correlation IDs, and redaction. Tradeoff: Search and aggregation improve, while schemas must evolve consistently. Watch for: Placing unbounded model output in one log field causes cost and ingestion failures.}
\paragraph{FC-0638 [medium]} Define Structured logging in interview-ready language.
\noindent\textbf{Answer.} Structured events use stable fields, severity, timestamps, correlation IDs, and redaction.
\par\textit{Structured events use stable fields, severity, timestamps, correlation IDs, and redaction.}
\paragraph{FC-0639 [hard]} State the main tradeoff and production pitfall for Structured logging.
\noindent\textbf{Answer.} Tradeoff: Search and aggregation improve, while schemas must evolve consistently. Pitfall: Placing unbounded model output in one log field causes cost and ingestion failures.
\par\textit{Tradeoff: Search and aggregation improve, while schemas must evolve consistently. Pitfall: Placing unbounded model output in one log field causes cost and ingestion failures.}
\paragraph{LA-0640 [hard]} Explain Structured logging from first principles, then show how you would evaluate and operate it in production.
\noindent\textbf{Answer.} Start with the mechanism: Structured events use stable fields, severity, timestamps, correlation IDs, and redaction. Make the tradeoff explicit: Search and aggregation improve, while schemas must evolve consistently. Define workload-specific offline metrics and a latency/cost budget, test important slices, instrument the critical path, and create a rollback criterion. The production review must explicitly guard against this failure: Placing unbounded model output in one log field causes cost and ingestion failures.
\par\textit{A strong answer connects mechanism, measurable tradeoffs, failure isolation, observability, and rollback rather than listing product features.}
\paragraph{MCQ-0641 [medium]} A design requires this behavior: LangSmith provides tracing, datasets, evaluation, prompt management, and deployment-oriented tooling for LLM applications. Which mechanism should the team study?
\begin{enumerate}[label=\Alph*.]
\item LangSmith
\item Cost observability
\item Evaluation in production
\item Prompt and model drift
\end{enumerate}
\noindent\textbf{Answer.} LangSmith
\par\textit{LangSmith provides tracing, datasets, evaluation, prompt management, and deployment-oriented tooling for LLM applications. Tradeoff: Tight LangChain integration accelerates work but requires portability planning. Watch for: Depending only on vendor UI annotations leaves evaluation logic unreproducible.}
\paragraph{MCQ-0642 [hard]} The team proposes LangSmith for a real workload. Which finding gives the correct decision boundary?
\begin{enumerate}[label=\Alph*.]
\item Alerts catch regressions but thresholds must distinguish seasonality from harm.
\item Tight LangChain integration accelerates work but requires portability planning.
\item Live coverage improves realism but evaluations must not expose or delay users.
\item Unit economics become actionable, while provider pricing and caching rules evolve.
\end{enumerate}
\noindent\textbf{Answer.} Tight LangChain integration accelerates work but requires portability planning.
\par\textit{LangSmith provides tracing, datasets, evaluation, prompt management, and deployment-oriented tooling for LLM applications. Tradeoff: Tight LangChain integration accelerates work but requires portability planning. Watch for: Depending only on vendor UI annotations leaves evaluation logic unreproducible.}
\paragraph{MCQ-0643 [hard]} Before releasing LangSmith, which failure deserves a dedicated regression test?
\begin{enumerate}[label=\Alph*.]
\item Depending only on vendor UI annotations leaves evaluation logic unreproducible.
\item Measuring embedding distribution shift alone does not prove answer quality changed.
\item Tracking only average cost hides pathological long-running agents.
\item Treating implicit clicks as unambiguous quality labels creates bias.
\end{enumerate}
\noindent\textbf{Answer.} Depending only on vendor UI annotations leaves evaluation logic unreproducible.
\par\textit{LangSmith provides tracing, datasets, evaluation, prompt management, and deployment-oriented tooling for LLM applications. Tradeoff: Tight LangChain integration accelerates work but requires portability planning. Watch for: Depending only on vendor UI annotations leaves evaluation logic unreproducible.}
\paragraph{MCQ-0644 [medium]} A design review reveals this mistake: Depending only on vendor UI annotations leaves evaluation logic unreproducible. Which mechanism should be traced first?
\begin{enumerate}[label=\Alph*.]
\item Evaluation in production
\item LangSmith
\item Cost observability
\item Prompt and model drift
\end{enumerate}
\noindent\textbf{Answer.} LangSmith
\par\textit{LangSmith provides tracing, datasets, evaluation, prompt management, and deployment-oriented tooling for LLM applications. Tradeoff: Tight LangChain integration accelerates work but requires portability planning. Watch for: Depending only on vendor UI annotations leaves evaluation logic unreproducible.}
\paragraph{MCQ-0645 [hard]} Which causal explanation of LangSmith connects what it does to when it is worth its cost?
\begin{enumerate}[label=\Alph*.]
\item Mechanism: Online evaluators combine sampled judges, deterministic checks, human feedback, and business outcomes. Consequence: Live coverage improves realism but evaluations must not expose or delay users.
\item Mechanism: LangSmith provides tracing, datasets, evaluation, prompt management, and deployment-oriented tooling for LLM applications. Consequence: Tight LangChain integration accelerates work but requires portability planning.
\item Mechanism: Cost telemetry attributes input, output, cache, retrieval, tool, and infrastructure spend to requests and tenants. Consequence: Unit economics become actionable, while provider pricing and caching rules evolve.
\item Mechanism: Drift monitoring detects changes in inputs, retrieved evidence, outputs, model versions, and user behavior. Consequence: Alerts catch regressions but thresholds must distinguish seasonality from harm.
\end{enumerate}
\noindent\textbf{Answer.} Mechanism: LangSmith provides tracing, datasets, evaluation, prompt management, and deployment-oriented tooling for LLM applications. Consequence: Tight LangChain integration accelerates work but requires portability planning.
\par\textit{LangSmith provides tracing, datasets, evaluation, prompt management, and deployment-oriented tooling for LLM applications. Tradeoff: Tight LangChain integration accelerates work but requires portability planning. Watch for: Depending only on vendor UI annotations leaves evaluation logic unreproducible.}
\paragraph{FC-0646 [medium]} Define LangSmith in interview-ready language.
\noindent\textbf{Answer.} LangSmith provides tracing, datasets, evaluation, prompt management, and deployment-oriented tooling for LLM applications.
\par\textit{LangSmith provides tracing, datasets, evaluation, prompt management, and deployment-oriented tooling for LLM applications.}
\paragraph{FC-0647 [hard]} State the main tradeoff and production pitfall for LangSmith.
\noindent\textbf{Answer.} Tradeoff: Tight LangChain integration accelerates work but requires portability planning. Pitfall: Depending only on vendor UI annotations leaves evaluation logic unreproducible.
\par\textit{Tradeoff: Tight LangChain integration accelerates work but requires portability planning. Pitfall: Depending only on vendor UI annotations leaves evaluation logic unreproducible.}
\paragraph{LA-0648 [hard]} Explain LangSmith from first principles, then show how you would evaluate and operate it in production.
\noindent\textbf{Answer.} Start with the mechanism: LangSmith provides tracing, datasets, evaluation, prompt management, and deployment-oriented tooling for LLM applications. Make the tradeoff explicit: Tight LangChain integration accelerates work but requires portability planning. Define workload-specific offline metrics and a latency/cost budget, test important slices, instrument the critical path, and create a rollback criterion. The production review must explicitly guard against this failure: Depending only on vendor UI annotations leaves evaluation logic unreproducible.
\par\textit{A strong answer connects mechanism, measurable tradeoffs, failure isolation, observability, and rollback rather than listing product features.}
\paragraph{MCQ-0649 [medium]} A design requires this behavior: Langfuse is an open-source platform for LLM traces, sessions, prompts, scores, datasets, and cost analytics. Which mechanism should the team study?
\begin{enumerate}[label=\Alph*.]
\item MLflow tracing
\item Structured logging
\item Cost observability
\item Langfuse
\end{enumerate}
\noindent\textbf{Answer.} Langfuse
\par\textit{Langfuse is an open-source platform for LLM traces, sessions, prompts, scores, datasets, and cost analytics. Tradeoff: Self-hosting offers control but shifts scaling, upgrades, and security to the team. Watch for: Capturing every token forever creates expensive sensitive retention.}
\paragraph{MCQ-0650 [hard]} The team proposes Langfuse for a real workload. Which finding gives the correct decision boundary?
\begin{enumerate}[label=\Alph*.]
\item Search and aggregation improve, while schemas must evolve consistently.
\item Unit economics become actionable, while provider pricing and caching rules evolve.
\item Self-hosting offers control but shifts scaling, upgrades, and security to the team.
\item A common ML platform improves lineage but may be heavier than a focused tracing tool.
\end{enumerate}
\noindent\textbf{Answer.} Self-hosting offers control but shifts scaling, upgrades, and security to the team.
\par\textit{Langfuse is an open-source platform for LLM traces, sessions, prompts, scores, datasets, and cost analytics. Tradeoff: Self-hosting offers control but shifts scaling, upgrades, and security to the team. Watch for: Capturing every token forever creates expensive sensitive retention.}
\paragraph{MCQ-0651 [hard]} Before releasing Langfuse, which failure deserves a dedicated regression test?
\begin{enumerate}[label=\Alph*.]
\item Mixing incomparable scorer versions in one trend chart produces false regressions.
\item Capturing every token forever creates expensive sensitive retention.
\item Placing unbounded model output in one log field causes cost and ingestion failures.
\item Tracking only average cost hides pathological long-running agents.
\end{enumerate}
\noindent\textbf{Answer.} Capturing every token forever creates expensive sensitive retention.
\par\textit{Langfuse is an open-source platform for LLM traces, sessions, prompts, scores, datasets, and cost analytics. Tradeoff: Self-hosting offers control but shifts scaling, upgrades, and security to the team. Watch for: Capturing every token forever creates expensive sensitive retention.}
\paragraph{MCQ-0652 [medium]} A design review reveals this mistake: Capturing every token forever creates expensive sensitive retention. Which mechanism should be traced first?
\begin{enumerate}[label=\Alph*.]
\item Cost observability
\item Langfuse
\item MLflow tracing
\item Structured logging
\end{enumerate}
\noindent\textbf{Answer.} Langfuse
\par\textit{Langfuse is an open-source platform for LLM traces, sessions, prompts, scores, datasets, and cost analytics. Tradeoff: Self-hosting offers control but shifts scaling, upgrades, and security to the team. Watch for: Capturing every token forever creates expensive sensitive retention.}
\paragraph{MCQ-0653 [hard]} Which causal explanation of Langfuse connects what it does to when it is worth its cost?
\begin{enumerate}[label=\Alph*.]
\item Mechanism: Langfuse is an open-source platform for LLM traces, sessions, prompts, scores, datasets, and cost analytics. Consequence: Self-hosting offers control but shifts scaling, upgrades, and security to the team.
\item Mechanism: Cost telemetry attributes input, output, cache, retrieval, tool, and infrastructure spend to requests and tenants. Consequence: Unit economics become actionable, while provider pricing and caching rules evolve.
\item Mechanism: Structured events use stable fields, severity, timestamps, correlation IDs, and redaction. Consequence: Search and aggregation improve, while schemas must evolve consistently.
\item Mechanism: MLflow connects traces, evaluation datasets, scorers, experiments, and production monitoring. Consequence: A common ML platform improves lineage but may be heavier than a focused tracing tool.
\end{enumerate}
\noindent\textbf{Answer.} Mechanism: Langfuse is an open-source platform for LLM traces, sessions, prompts, scores, datasets, and cost analytics. Consequence: Self-hosting offers control but shifts scaling, upgrades, and security to the team.
\par\textit{Langfuse is an open-source platform for LLM traces, sessions, prompts, scores, datasets, and cost analytics. Tradeoff: Self-hosting offers control but shifts scaling, upgrades, and security to the team. Watch for: Capturing every token forever creates expensive sensitive retention.}
\paragraph{FC-0654 [medium]} Define Langfuse in interview-ready language.
\noindent\textbf{Answer.} Langfuse is an open-source platform for LLM traces, sessions, prompts, scores, datasets, and cost analytics.
\par\textit{Langfuse is an open-source platform for LLM traces, sessions, prompts, scores, datasets, and cost analytics.}
\paragraph{FC-0655 [hard]} State the main tradeoff and production pitfall for Langfuse.
\noindent\textbf{Answer.} Tradeoff: Self-hosting offers control but shifts scaling, upgrades, and security to the team. Pitfall: Capturing every token forever creates expensive sensitive retention.
\par\textit{Tradeoff: Self-hosting offers control but shifts scaling, upgrades, and security to the team. Pitfall: Capturing every token forever creates expensive sensitive retention.}
\paragraph{LA-0656 [hard]} Explain Langfuse from first principles, then show how you would evaluate and operate it in production.
\noindent\textbf{Answer.} Start with the mechanism: Langfuse is an open-source platform for LLM traces, sessions, prompts, scores, datasets, and cost analytics. Make the tradeoff explicit: Self-hosting offers control but shifts scaling, upgrades, and security to the team. Define workload-specific offline metrics and a latency/cost budget, test important slices, instrument the critical path, and create a rollback criterion. The production review must explicitly guard against this failure: Capturing every token forever creates expensive sensitive retention.
\par\textit{A strong answer connects mechanism, measurable tradeoffs, failure isolation, observability, and rollback rather than listing product features.}
\paragraph{MCQ-0657 [medium]} A design requires this behavior: Phoenix provides open-source tracing and evaluation for LLM, retrieval, and agent applications with OpenTelemetry support. Which mechanism should the team study?
\begin{enumerate}[label=\Alph*.]
\item LLM tracing
\item W\&B Weave
\item Arize Phoenix
\item OpenTelemetry
\end{enumerate}
\noindent\textbf{Answer.} Arize Phoenix
\par\textit{Phoenix provides open-source tracing and evaluation for LLM, retrieval, and agent applications with OpenTelemetry support. Tradeoff: It unifies qualitative traces and quantitative evals, while production operations still need design. Watch for: Judge scores without slice analysis can mask failures for important cohorts.}
\paragraph{MCQ-0658 [hard]} The team proposes Arize Phoenix for a real workload. Which finding gives the correct decision boundary?
\begin{enumerate}[label=\Alph*.]
\item Rich experiment lineage helps collaboration but introduces another governance surface.
\item Portability improves, while instrumentation quality and backend cost remain engineering work.
\item Debugging improves, while content capture raises privacy and retention risk.
\item It unifies qualitative traces and quantitative evals, while production operations still need design.
\end{enumerate}
\noindent\textbf{Answer.} It unifies qualitative traces and quantitative evals, while production operations still need design.
\par\textit{Phoenix provides open-source tracing and evaluation for LLM, retrieval, and agent applications with OpenTelemetry support. Tradeoff: It unifies qualitative traces and quantitative evals, while production operations still need design. Watch for: Judge scores without slice analysis can mask failures for important cohorts.}
\paragraph{MCQ-0659 [hard]} Before releasing Arize Phoenix, which failure deserves a dedicated regression test?
\begin{enumerate}[label=\Alph*.]
\item Logging secrets or personal data by default creates a second sensitive datastore.
\item Saving mutable datasets without immutable versions undermines reproducibility.
\item High-cardinality prompt or user attributes can overwhelm telemetry systems.
\item Judge scores without slice analysis can mask failures for important cohorts.
\end{enumerate}
\noindent\textbf{Answer.} Judge scores without slice analysis can mask failures for important cohorts.
\par\textit{Phoenix provides open-source tracing and evaluation for LLM, retrieval, and agent applications with OpenTelemetry support. Tradeoff: It unifies qualitative traces and quantitative evals, while production operations still need design. Watch for: Judge scores without slice analysis can mask failures for important cohorts.}
\paragraph{MCQ-0660 [medium]} A design review reveals this mistake: Judge scores without slice analysis can mask failures for important cohorts. Which mechanism should be traced first?
\begin{enumerate}[label=\Alph*.]
\item OpenTelemetry
\item Arize Phoenix
\item LLM tracing
\item W\&B Weave
\end{enumerate}
\noindent\textbf{Answer.} Arize Phoenix
\par\textit{Phoenix provides open-source tracing and evaluation for LLM, retrieval, and agent applications with OpenTelemetry support. Tradeoff: It unifies qualitative traces and quantitative evals, while production operations still need design. Watch for: Judge scores without slice analysis can mask failures for important cohorts.}
\paragraph{MCQ-0661 [hard]} Which causal explanation of Arize Phoenix connects what it does to when it is worth its cost?
\begin{enumerate}[label=\Alph*.]
\item Mechanism: Phoenix provides open-source tracing and evaluation for LLM, retrieval, and agent applications with OpenTelemetry support. Consequence: It unifies qualitative traces and quantitative evals, while production operations still need design.
\item Mechanism: Weave tracks calls, objects, datasets, evaluations, and model application experiments. Consequence: Rich experiment lineage helps collaboration but introduces another governance surface.
\item Mechanism: OpenTelemetry provides vendor-neutral APIs, SDKs, collectors, and semantic conventions for traces, metrics, logs, and profiles. Consequence: Portability improves, while instrumentation quality and backend cost remain engineering work.
\item Mechanism: LLM traces capture prompts, outputs, model settings, token usage, latency, tool calls, retrieval, and evaluations. Consequence: Debugging improves, while content capture raises privacy and retention risk.
\end{enumerate}
\noindent\textbf{Answer.} Mechanism: Phoenix provides open-source tracing and evaluation for LLM, retrieval, and agent applications with OpenTelemetry support. Consequence: It unifies qualitative traces and quantitative evals, while production operations still need design.
\par\textit{Phoenix provides open-source tracing and evaluation for LLM, retrieval, and agent applications with OpenTelemetry support. Tradeoff: It unifies qualitative traces and quantitative evals, while production operations still need design. Watch for: Judge scores without slice analysis can mask failures for important cohorts.}
\paragraph{FC-0662 [medium]} Define Arize Phoenix in interview-ready language.
\noindent\textbf{Answer.} Phoenix provides open-source tracing and evaluation for LLM, retrieval, and agent applications with OpenTelemetry support.
\par\textit{Phoenix provides open-source tracing and evaluation for LLM, retrieval, and agent applications with OpenTelemetry support.}
\paragraph{FC-0663 [hard]} State the main tradeoff and production pitfall for Arize Phoenix.
\noindent\textbf{Answer.} Tradeoff: It unifies qualitative traces and quantitative evals, while production operations still need design. Pitfall: Judge scores without slice analysis can mask failures for important cohorts.
\par\textit{Tradeoff: It unifies qualitative traces and quantitative evals, while production operations still need design. Pitfall: Judge scores without slice analysis can mask failures for important cohorts.}
\paragraph{LA-0664 [hard]} Explain Arize Phoenix from first principles, then show how you would evaluate and operate it in production.
\noindent\textbf{Answer.} Start with the mechanism: Phoenix provides open-source tracing and evaluation for LLM, retrieval, and agent applications with OpenTelemetry support. Make the tradeoff explicit: It unifies qualitative traces and quantitative evals, while production operations still need design. Define workload-specific offline metrics and a latency/cost budget, test important slices, instrument the critical path, and create a rollback criterion. The production review must explicitly guard against this failure: Judge scores without slice analysis can mask failures for important cohorts.
\par\textit{A strong answer connects mechanism, measurable tradeoffs, failure isolation, observability, and rollback rather than listing product features.}
\paragraph{MCQ-0665 [medium]} A design requires this behavior: MLflow connects traces, evaluation datasets, scorers, experiments, and production monitoring. Which mechanism should the team study?
\begin{enumerate}[label=\Alph*.]
\item LangSmith
\item Langfuse
\item Evaluation in production
\item MLflow tracing
\end{enumerate}
\noindent\textbf{Answer.} MLflow tracing
\par\textit{MLflow connects traces, evaluation datasets, scorers, experiments, and production monitoring. Tradeoff: A common ML platform improves lineage but may be heavier than a focused tracing tool. Watch for: Mixing incomparable scorer versions in one trend chart produces false regressions.}
\paragraph{MCQ-0666 [hard]} The team proposes MLflow tracing for a real workload. Which finding gives the correct decision boundary?
\begin{enumerate}[label=\Alph*.]
\item Live coverage improves realism but evaluations must not expose or delay users.
\item Tight LangChain integration accelerates work but requires portability planning.
\item Self-hosting offers control but shifts scaling, upgrades, and security to the team.
\item A common ML platform improves lineage but may be heavier than a focused tracing tool.
\end{enumerate}
\noindent\textbf{Answer.} A common ML platform improves lineage but may be heavier than a focused tracing tool.
\par\textit{MLflow connects traces, evaluation datasets, scorers, experiments, and production monitoring. Tradeoff: A common ML platform improves lineage but may be heavier than a focused tracing tool. Watch for: Mixing incomparable scorer versions in one trend chart produces false regressions.}
\paragraph{MCQ-0667 [hard]} Before releasing MLflow tracing, which failure deserves a dedicated regression test?
\begin{enumerate}[label=\Alph*.]
\item Mixing incomparable scorer versions in one trend chart produces false regressions.
\item Treating implicit clicks as unambiguous quality labels creates bias.
\item Capturing every token forever creates expensive sensitive retention.
\item Depending only on vendor UI annotations leaves evaluation logic unreproducible.
\end{enumerate}
\noindent\textbf{Answer.} Mixing incomparable scorer versions in one trend chart produces false regressions.
\par\textit{MLflow connects traces, evaluation datasets, scorers, experiments, and production monitoring. Tradeoff: A common ML platform improves lineage but may be heavier than a focused tracing tool. Watch for: Mixing incomparable scorer versions in one trend chart produces false regressions.}
\paragraph{MCQ-0668 [medium]} A design review reveals this mistake: Mixing incomparable scorer versions in one trend chart produces false regressions. Which mechanism should be traced first?
\begin{enumerate}[label=\Alph*.]
\item Evaluation in production
\item MLflow tracing
\item Langfuse
\item LangSmith
\end{enumerate}
\noindent\textbf{Answer.} MLflow tracing
\par\textit{MLflow connects traces, evaluation datasets, scorers, experiments, and production monitoring. Tradeoff: A common ML platform improves lineage but may be heavier than a focused tracing tool. Watch for: Mixing incomparable scorer versions in one trend chart produces false regressions.}
\paragraph{MCQ-0669 [hard]} Which causal explanation of MLflow tracing connects what it does to when it is worth its cost?
\begin{enumerate}[label=\Alph*.]
\item Mechanism: MLflow connects traces, evaluation datasets, scorers, experiments, and production monitoring. Consequence: A common ML platform improves lineage but may be heavier than a focused tracing tool.
\item Mechanism: Online evaluators combine sampled judges, deterministic checks, human feedback, and business outcomes. Consequence: Live coverage improves realism but evaluations must not expose or delay users.
\item Mechanism: Langfuse is an open-source platform for LLM traces, sessions, prompts, scores, datasets, and cost analytics. Consequence: Self-hosting offers control but shifts scaling, upgrades, and security to the team.
\item Mechanism: LangSmith provides tracing, datasets, evaluation, prompt management, and deployment-oriented tooling for LLM applications. Consequence: Tight LangChain integration accelerates work but requires portability planning.
\end{enumerate}
\noindent\textbf{Answer.} Mechanism: MLflow connects traces, evaluation datasets, scorers, experiments, and production monitoring. Consequence: A common ML platform improves lineage but may be heavier than a focused tracing tool.
\par\textit{MLflow connects traces, evaluation datasets, scorers, experiments, and production monitoring. Tradeoff: A common ML platform improves lineage but may be heavier than a focused tracing tool. Watch for: Mixing incomparable scorer versions in one trend chart produces false regressions.}
\paragraph{FC-0670 [medium]} Define MLflow tracing in interview-ready language.
\noindent\textbf{Answer.} MLflow connects traces, evaluation datasets, scorers, experiments, and production monitoring.
\par\textit{MLflow connects traces, evaluation datasets, scorers, experiments, and production monitoring.}
\paragraph{FC-0671 [hard]} State the main tradeoff and production pitfall for MLflow tracing.
\noindent\textbf{Answer.} Tradeoff: A common ML platform improves lineage but may be heavier than a focused tracing tool. Pitfall: Mixing incomparable scorer versions in one trend chart produces false regressions.
\par\textit{Tradeoff: A common ML platform improves lineage but may be heavier than a focused tracing tool. Pitfall: Mixing incomparable scorer versions in one trend chart produces false regressions.}
\paragraph{LA-0672 [hard]} Explain MLflow tracing from first principles, then show how you would evaluate and operate it in production.
\noindent\textbf{Answer.} Start with the mechanism: MLflow connects traces, evaluation datasets, scorers, experiments, and production monitoring. Make the tradeoff explicit: A common ML platform improves lineage but may be heavier than a focused tracing tool. Define workload-specific offline metrics and a latency/cost budget, test important slices, instrument the critical path, and create a rollback criterion. The production review must explicitly guard against this failure: Mixing incomparable scorer versions in one trend chart produces false regressions.
\par\textit{A strong answer connects mechanism, measurable tradeoffs, failure isolation, observability, and rollback rather than listing product features.}
\paragraph{MCQ-0673 [medium]} A design requires this behavior: Weave tracks calls, objects, datasets, evaluations, and model application experiments. Which mechanism should the team study?
\begin{enumerate}[label=\Alph*.]
\item Langfuse
\item W\&B Weave
\item Evaluation in production
\item Cost observability
\end{enumerate}
\noindent\textbf{Answer.} W\&B Weave
\par\textit{Weave tracks calls, objects, datasets, evaluations, and model application experiments. Tradeoff: Rich experiment lineage helps collaboration but introduces another governance surface. Watch for: Saving mutable datasets without immutable versions undermines reproducibility.}
\paragraph{MCQ-0674 [hard]} The team proposes W\&B Weave for a real workload. Which finding gives the correct decision boundary?
\begin{enumerate}[label=\Alph*.]
\item Rich experiment lineage helps collaboration but introduces another governance surface.
\item Unit economics become actionable, while provider pricing and caching rules evolve.
\item Self-hosting offers control but shifts scaling, upgrades, and security to the team.
\item Live coverage improves realism but evaluations must not expose or delay users.
\end{enumerate}
\noindent\textbf{Answer.} Rich experiment lineage helps collaboration but introduces another governance surface.
\par\textit{Weave tracks calls, objects, datasets, evaluations, and model application experiments. Tradeoff: Rich experiment lineage helps collaboration but introduces another governance surface. Watch for: Saving mutable datasets without immutable versions undermines reproducibility.}
\paragraph{MCQ-0675 [hard]} Before releasing W\&B Weave, which failure deserves a dedicated regression test?
\begin{enumerate}[label=\Alph*.]
\item Tracking only average cost hides pathological long-running agents.
\item Capturing every token forever creates expensive sensitive retention.
\item Saving mutable datasets without immutable versions undermines reproducibility.
\item Treating implicit clicks as unambiguous quality labels creates bias.
\end{enumerate}
\noindent\textbf{Answer.} Saving mutable datasets without immutable versions undermines reproducibility.
\par\textit{Weave tracks calls, objects, datasets, evaluations, and model application experiments. Tradeoff: Rich experiment lineage helps collaboration but introduces another governance surface. Watch for: Saving mutable datasets without immutable versions undermines reproducibility.}
\paragraph{MCQ-0676 [medium]} A design review reveals this mistake: Saving mutable datasets without immutable versions undermines reproducibility. Which mechanism should be traced first?
\begin{enumerate}[label=\Alph*.]
\item W\&B Weave
\item Langfuse
\item Evaluation in production
\item Cost observability
\end{enumerate}
\noindent\textbf{Answer.} W\&B Weave
\par\textit{Weave tracks calls, objects, datasets, evaluations, and model application experiments. Tradeoff: Rich experiment lineage helps collaboration but introduces another governance surface. Watch for: Saving mutable datasets without immutable versions undermines reproducibility.}
\paragraph{MCQ-0677 [hard]} Which causal explanation of W\&B Weave connects what it does to when it is worth its cost?
\begin{enumerate}[label=\Alph*.]
\item Mechanism: Cost telemetry attributes input, output, cache, retrieval, tool, and infrastructure spend to requests and tenants. Consequence: Unit economics become actionable, while provider pricing and caching rules evolve.
\item Mechanism: Online evaluators combine sampled judges, deterministic checks, human feedback, and business outcomes. Consequence: Live coverage improves realism but evaluations must not expose or delay users.
\item Mechanism: Langfuse is an open-source platform for LLM traces, sessions, prompts, scores, datasets, and cost analytics. Consequence: Self-hosting offers control but shifts scaling, upgrades, and security to the team.
\item Mechanism: Weave tracks calls, objects, datasets, evaluations, and model application experiments. Consequence: Rich experiment lineage helps collaboration but introduces another governance surface.
\end{enumerate}
\noindent\textbf{Answer.} Mechanism: Weave tracks calls, objects, datasets, evaluations, and model application experiments. Consequence: Rich experiment lineage helps collaboration but introduces another governance surface.
\par\textit{Weave tracks calls, objects, datasets, evaluations, and model application experiments. Tradeoff: Rich experiment lineage helps collaboration but introduces another governance surface. Watch for: Saving mutable datasets without immutable versions undermines reproducibility.}
\paragraph{FC-0678 [medium]} Define W\&B Weave in interview-ready language.
\noindent\textbf{Answer.} Weave tracks calls, objects, datasets, evaluations, and model application experiments.
\par\textit{Weave tracks calls, objects, datasets, evaluations, and model application experiments.}
\paragraph{FC-0679 [hard]} State the main tradeoff and production pitfall for W\&B Weave.
\noindent\textbf{Answer.} Tradeoff: Rich experiment lineage helps collaboration but introduces another governance surface. Pitfall: Saving mutable datasets without immutable versions undermines reproducibility.
\par\textit{Tradeoff: Rich experiment lineage helps collaboration but introduces another governance surface. Pitfall: Saving mutable datasets without immutable versions undermines reproducibility.}
\paragraph{LA-0680 [hard]} Explain W\&B Weave from first principles, then show how you would evaluate and operate it in production.
\noindent\textbf{Answer.} Start with the mechanism: Weave tracks calls, objects, datasets, evaluations, and model application experiments. Make the tradeoff explicit: Rich experiment lineage helps collaboration but introduces another governance surface. Define workload-specific offline metrics and a latency/cost budget, test important slices, instrument the critical path, and create a rollback criterion. The production review must explicitly guard against this failure: Saving mutable datasets without immutable versions undermines reproducibility.
\par\textit{A strong answer connects mechanism, measurable tradeoffs, failure isolation, observability, and rollback rather than listing product features.}
\paragraph{MCQ-0681 [medium]} A design requires this behavior: Drift monitoring detects changes in inputs, retrieved evidence, outputs, model versions, and user behavior. Which mechanism should the team study?
\begin{enumerate}[label=\Alph*.]
\item LLM tracing
\item Metrics and SLOs
\item Prompt and model drift
\item LangSmith
\end{enumerate}
\noindent\textbf{Answer.} Prompt and model drift
\par\textit{Drift monitoring detects changes in inputs, retrieved evidence, outputs, model versions, and user behavior. Tradeoff: Alerts catch regressions but thresholds must distinguish seasonality from harm. Watch for: Measuring embedding distribution shift alone does not prove answer quality changed.}
\paragraph{MCQ-0682 [hard]} The team proposes Prompt and model drift for a real workload. Which finding gives the correct decision boundary?
\begin{enumerate}[label=\Alph*.]
\item Tight LangChain integration accelerates work but requires portability planning.
\item Aggregate signals scale well but cannot explain individual semantic failures.
\item Alerts catch regressions but thresholds must distinguish seasonality from harm.
\item Debugging improves, while content capture raises privacy and retention risk.
\end{enumerate}
\noindent\textbf{Answer.} Alerts catch regressions but thresholds must distinguish seasonality from harm.
\par\textit{Drift monitoring detects changes in inputs, retrieved evidence, outputs, model versions, and user behavior. Tradeoff: Alerts catch regressions but thresholds must distinguish seasonality from harm. Watch for: Measuring embedding distribution shift alone does not prove answer quality changed.}
\paragraph{MCQ-0683 [hard]} Before releasing Prompt and model drift, which failure deserves a dedicated regression test?
\begin{enumerate}[label=\Alph*.]
\item Logging secrets or personal data by default creates a second sensitive datastore.
\item Depending only on vendor UI annotations leaves evaluation logic unreproducible.
\item Measuring embedding distribution shift alone does not prove answer quality changed.
\item Averaging latency hides long-tail p95 and p99 user pain.
\end{enumerate}
\noindent\textbf{Answer.} Measuring embedding distribution shift alone does not prove answer quality changed.
\par\textit{Drift monitoring detects changes in inputs, retrieved evidence, outputs, model versions, and user behavior. Tradeoff: Alerts catch regressions but thresholds must distinguish seasonality from harm. Watch for: Measuring embedding distribution shift alone does not prove answer quality changed.}
\paragraph{MCQ-0684 [medium]} A design review reveals this mistake: Measuring embedding distribution shift alone does not prove answer quality changed. Which mechanism should be traced first?
\begin{enumerate}[label=\Alph*.]
\item LangSmith
\item Prompt and model drift
\item LLM tracing
\item Metrics and SLOs
\end{enumerate}
\noindent\textbf{Answer.} Prompt and model drift
\par\textit{Drift monitoring detects changes in inputs, retrieved evidence, outputs, model versions, and user behavior. Tradeoff: Alerts catch regressions but thresholds must distinguish seasonality from harm. Watch for: Measuring embedding distribution shift alone does not prove answer quality changed.}
\paragraph{MCQ-0685 [hard]} Which causal explanation of Prompt and model drift connects what it does to when it is worth its cost?
\begin{enumerate}[label=\Alph*.]
\item Mechanism: LangSmith provides tracing, datasets, evaluation, prompt management, and deployment-oriented tooling for LLM applications. Consequence: Tight LangChain integration accelerates work but requires portability planning.
\item Mechanism: Counters, histograms, gauges, service-level indicators, and objectives quantify reliability and performance. Consequence: Aggregate signals scale well but cannot explain individual semantic failures.
\item Mechanism: Drift monitoring detects changes in inputs, retrieved evidence, outputs, model versions, and user behavior. Consequence: Alerts catch regressions but thresholds must distinguish seasonality from harm.
\item Mechanism: LLM traces capture prompts, outputs, model settings, token usage, latency, tool calls, retrieval, and evaluations. Consequence: Debugging improves, while content capture raises privacy and retention risk.
\end{enumerate}
\noindent\textbf{Answer.} Mechanism: Drift monitoring detects changes in inputs, retrieved evidence, outputs, model versions, and user behavior. Consequence: Alerts catch regressions but thresholds must distinguish seasonality from harm.
\par\textit{Drift monitoring detects changes in inputs, retrieved evidence, outputs, model versions, and user behavior. Tradeoff: Alerts catch regressions but thresholds must distinguish seasonality from harm. Watch for: Measuring embedding distribution shift alone does not prove answer quality changed.}
\paragraph{FC-0686 [medium]} Define Prompt and model drift in interview-ready language.
\noindent\textbf{Answer.} Drift monitoring detects changes in inputs, retrieved evidence, outputs, model versions, and user behavior.
\par\textit{Drift monitoring detects changes in inputs, retrieved evidence, outputs, model versions, and user behavior.}
\paragraph{FC-0687 [hard]} State the main tradeoff and production pitfall for Prompt and model drift.
\noindent\textbf{Answer.} Tradeoff: Alerts catch regressions but thresholds must distinguish seasonality from harm. Pitfall: Measuring embedding distribution shift alone does not prove answer quality changed.
\par\textit{Tradeoff: Alerts catch regressions but thresholds must distinguish seasonality from harm. Pitfall: Measuring embedding distribution shift alone does not prove answer quality changed.}
\paragraph{LA-0688 [hard]} Explain Prompt and model drift from first principles, then show how you would evaluate and operate it in production.
\noindent\textbf{Answer.} Start with the mechanism: Drift monitoring detects changes in inputs, retrieved evidence, outputs, model versions, and user behavior. Make the tradeoff explicit: Alerts catch regressions but thresholds must distinguish seasonality from harm. Define workload-specific offline metrics and a latency/cost budget, test important slices, instrument the critical path, and create a rollback criterion. The production review must explicitly guard against this failure: Measuring embedding distribution shift alone does not prove answer quality changed.
\par\textit{A strong answer connects mechanism, measurable tradeoffs, failure isolation, observability, and rollback rather than listing product features.}
\paragraph{MCQ-0689 [medium]} A design requires this behavior: Cost telemetry attributes input, output, cache, retrieval, tool, and infrastructure spend to requests and tenants. Which mechanism should the team study?
\begin{enumerate}[label=\Alph*.]
\item W\&B Weave
\item LLM tracing
\item Cost observability
\item Evaluation in production
\end{enumerate}
\noindent\textbf{Answer.} Cost observability
\par\textit{Cost telemetry attributes input, output, cache, retrieval, tool, and infrastructure spend to requests and tenants. Tradeoff: Unit economics become actionable, while provider pricing and caching rules evolve. Watch for: Tracking only average cost hides pathological long-running agents.}
\paragraph{MCQ-0690 [hard]} The team proposes Cost observability for a real workload. Which finding gives the correct decision boundary?
\begin{enumerate}[label=\Alph*.]
\item Debugging improves, while content capture raises privacy and retention risk.
\item Unit economics become actionable, while provider pricing and caching rules evolve.
\item Live coverage improves realism but evaluations must not expose or delay users.
\item Rich experiment lineage helps collaboration but introduces another governance surface.
\end{enumerate}
\noindent\textbf{Answer.} Unit economics become actionable, while provider pricing and caching rules evolve.
\par\textit{Cost telemetry attributes input, output, cache, retrieval, tool, and infrastructure spend to requests and tenants. Tradeoff: Unit economics become actionable, while provider pricing and caching rules evolve. Watch for: Tracking only average cost hides pathological long-running agents.}
\paragraph{MCQ-0691 [hard]} Before releasing Cost observability, which failure deserves a dedicated regression test?
\begin{enumerate}[label=\Alph*.]
\item Tracking only average cost hides pathological long-running agents.
\item Logging secrets or personal data by default creates a second sensitive datastore.
\item Saving mutable datasets without immutable versions undermines reproducibility.
\item Treating implicit clicks as unambiguous quality labels creates bias.
\end{enumerate}
\noindent\textbf{Answer.} Tracking only average cost hides pathological long-running agents.
\par\textit{Cost telemetry attributes input, output, cache, retrieval, tool, and infrastructure spend to requests and tenants. Tradeoff: Unit economics become actionable, while provider pricing and caching rules evolve. Watch for: Tracking only average cost hides pathological long-running agents.}
\paragraph{MCQ-0692 [medium]} A design review reveals this mistake: Tracking only average cost hides pathological long-running agents. Which mechanism should be traced first?
\begin{enumerate}[label=\Alph*.]
\item Evaluation in production
\item W\&B Weave
\item Cost observability
\item LLM tracing
\end{enumerate}
\noindent\textbf{Answer.} Cost observability
\par\textit{Cost telemetry attributes input, output, cache, retrieval, tool, and infrastructure spend to requests and tenants. Tradeoff: Unit economics become actionable, while provider pricing and caching rules evolve. Watch for: Tracking only average cost hides pathological long-running agents.}
\paragraph{MCQ-0693 [hard]} Which causal explanation of Cost observability connects what it does to when it is worth its cost?
\begin{enumerate}[label=\Alph*.]
\item Mechanism: LLM traces capture prompts, outputs, model settings, token usage, latency, tool calls, retrieval, and evaluations. Consequence: Debugging improves, while content capture raises privacy and retention risk.
\item Mechanism: Cost telemetry attributes input, output, cache, retrieval, tool, and infrastructure spend to requests and tenants. Consequence: Unit economics become actionable, while provider pricing and caching rules evolve.
\item Mechanism: Online evaluators combine sampled judges, deterministic checks, human feedback, and business outcomes. Consequence: Live coverage improves realism but evaluations must not expose or delay users.
\item Mechanism: Weave tracks calls, objects, datasets, evaluations, and model application experiments. Consequence: Rich experiment lineage helps collaboration but introduces another governance surface.
\end{enumerate}
\noindent\textbf{Answer.} Mechanism: Cost telemetry attributes input, output, cache, retrieval, tool, and infrastructure spend to requests and tenants. Consequence: Unit economics become actionable, while provider pricing and caching rules evolve.
\par\textit{Cost telemetry attributes input, output, cache, retrieval, tool, and infrastructure spend to requests and tenants. Tradeoff: Unit economics become actionable, while provider pricing and caching rules evolve. Watch for: Tracking only average cost hides pathological long-running agents.}
\paragraph{FC-0694 [medium]} Define Cost observability in interview-ready language.
\noindent\textbf{Answer.} Cost telemetry attributes input, output, cache, retrieval, tool, and infrastructure spend to requests and tenants.
\par\textit{Cost telemetry attributes input, output, cache, retrieval, tool, and infrastructure spend to requests and tenants.}
\paragraph{FC-0695 [hard]} State the main tradeoff and production pitfall for Cost observability.
\noindent\textbf{Answer.} Tradeoff: Unit economics become actionable, while provider pricing and caching rules evolve. Pitfall: Tracking only average cost hides pathological long-running agents.
\par\textit{Tradeoff: Unit economics become actionable, while provider pricing and caching rules evolve. Pitfall: Tracking only average cost hides pathological long-running agents.}
\paragraph{LA-0696 [hard]} Explain Cost observability from first principles, then show how you would evaluate and operate it in production.
\noindent\textbf{Answer.} Start with the mechanism: Cost telemetry attributes input, output, cache, retrieval, tool, and infrastructure spend to requests and tenants. Make the tradeoff explicit: Unit economics become actionable, while provider pricing and caching rules evolve. Define workload-specific offline metrics and a latency/cost budget, test important slices, instrument the critical path, and create a rollback criterion. The production review must explicitly guard against this failure: Tracking only average cost hides pathological long-running agents.
\par\textit{A strong answer connects mechanism, measurable tradeoffs, failure isolation, observability, and rollback rather than listing product features.}
\paragraph{MCQ-0697 [medium]} A design requires this behavior: Online evaluators combine sampled judges, deterministic checks, human feedback, and business outcomes. Which mechanism should the team study?
\begin{enumerate}[label=\Alph*.]
\item Evaluation in production
\item Prompt and model drift
\item LangSmith
\item W\&B Weave
\end{enumerate}
\noindent\textbf{Answer.} Evaluation in production
\par\textit{Online evaluators combine sampled judges, deterministic checks, human feedback, and business outcomes. Tradeoff: Live coverage improves realism but evaluations must not expose or delay users. Watch for: Treating implicit clicks as unambiguous quality labels creates bias.}
\paragraph{MCQ-0698 [hard]} The team proposes Evaluation in production for a real workload. Which finding gives the correct decision boundary?
\begin{enumerate}[label=\Alph*.]
\item Live coverage improves realism but evaluations must not expose or delay users.
\item Tight LangChain integration accelerates work but requires portability planning.
\item Alerts catch regressions but thresholds must distinguish seasonality from harm.
\item Rich experiment lineage helps collaboration but introduces another governance surface.
\end{enumerate}
\noindent\textbf{Answer.} Live coverage improves realism but evaluations must not expose or delay users.
\par\textit{Online evaluators combine sampled judges, deterministic checks, human feedback, and business outcomes. Tradeoff: Live coverage improves realism but evaluations must not expose or delay users. Watch for: Treating implicit clicks as unambiguous quality labels creates bias.}
\paragraph{MCQ-0699 [hard]} Before releasing Evaluation in production, which failure deserves a dedicated regression test?
\begin{enumerate}[label=\Alph*.]
\item Saving mutable datasets without immutable versions undermines reproducibility.
\item Measuring embedding distribution shift alone does not prove answer quality changed.
\item Treating implicit clicks as unambiguous quality labels creates bias.
\item Depending only on vendor UI annotations leaves evaluation logic unreproducible.
\end{enumerate}
\noindent\textbf{Answer.} Treating implicit clicks as unambiguous quality labels creates bias.
\par\textit{Online evaluators combine sampled judges, deterministic checks, human feedback, and business outcomes. Tradeoff: Live coverage improves realism but evaluations must not expose or delay users. Watch for: Treating implicit clicks as unambiguous quality labels creates bias.}
\paragraph{MCQ-0700 [medium]} A design review reveals this mistake: Treating implicit clicks as unambiguous quality labels creates bias. Which mechanism should be traced first?
\begin{enumerate}[label=\Alph*.]
\item LangSmith
\item W\&B Weave
\item Evaluation in production
\item Prompt and model drift
\end{enumerate}
\noindent\textbf{Answer.} Evaluation in production
\par\textit{Online evaluators combine sampled judges, deterministic checks, human feedback, and business outcomes. Tradeoff: Live coverage improves realism but evaluations must not expose or delay users. Watch for: Treating implicit clicks as unambiguous quality labels creates bias.}
\paragraph{MCQ-0701 [hard]} Which causal explanation of Evaluation in production connects what it does to when it is worth its cost?
\begin{enumerate}[label=\Alph*.]
\item Mechanism: Weave tracks calls, objects, datasets, evaluations, and model application experiments. Consequence: Rich experiment lineage helps collaboration but introduces another governance surface.
\item Mechanism: Online evaluators combine sampled judges, deterministic checks, human feedback, and business outcomes. Consequence: Live coverage improves realism but evaluations must not expose or delay users.
\item Mechanism: LangSmith provides tracing, datasets, evaluation, prompt management, and deployment-oriented tooling for LLM applications. Consequence: Tight LangChain integration accelerates work but requires portability planning.
\item Mechanism: Drift monitoring detects changes in inputs, retrieved evidence, outputs, model versions, and user behavior. Consequence: Alerts catch regressions but thresholds must distinguish seasonality from harm.
\end{enumerate}
\noindent\textbf{Answer.} Mechanism: Online evaluators combine sampled judges, deterministic checks, human feedback, and business outcomes. Consequence: Live coverage improves realism but evaluations must not expose or delay users.
\par\textit{Online evaluators combine sampled judges, deterministic checks, human feedback, and business outcomes. Tradeoff: Live coverage improves realism but evaluations must not expose or delay users. Watch for: Treating implicit clicks as unambiguous quality labels creates bias.}
\paragraph{FC-0702 [medium]} Define Evaluation in production in interview-ready language.
\noindent\textbf{Answer.} Online evaluators combine sampled judges, deterministic checks, human feedback, and business outcomes.
\par\textit{Online evaluators combine sampled judges, deterministic checks, human feedback, and business outcomes.}
\paragraph{FC-0703 [hard]} State the main tradeoff and production pitfall for Evaluation in production.
\noindent\textbf{Answer.} Tradeoff: Live coverage improves realism but evaluations must not expose or delay users. Pitfall: Treating implicit clicks as unambiguous quality labels creates bias.
\par\textit{Tradeoff: Live coverage improves realism but evaluations must not expose or delay users. Pitfall: Treating implicit clicks as unambiguous quality labels creates bias.}
\paragraph{LA-0704 [hard]} Explain Evaluation in production from first principles, then show how you would evaluate and operate it in production.
\noindent\textbf{Answer.} Start with the mechanism: Online evaluators combine sampled judges, deterministic checks, human feedback, and business outcomes. Make the tradeoff explicit: Live coverage improves realism but evaluations must not expose or delay users. Define workload-specific offline metrics and a latency/cost budget, test important slices, instrument the critical path, and create a rollback criterion. The production review must explicitly guard against this failure: Treating implicit clicks as unambiguous quality labels creates bias.
\par\textit{A strong answer connects mechanism, measurable tradeoffs, failure isolation, observability, and rollback rather than listing product features.}
\subsection{Serving, Deployment, and LLMOps}
\paragraph{MCQ-0705 [medium]} A design requires this behavior: vLLM is a high-throughput inference server using PagedAttention, continuous batching, parallelism, and OpenAI-compatible APIs. Which mechanism should the team study?
\begin{enumerate}[label=\Alph*.]
\item Model registry
\item Text Generation Inference
\item vLLM
\item Kubernetes and GPUs
\end{enumerate}
\noindent\textbf{Answer.} vLLM
\par\textit{vLLM is a high-throughput inference server using PagedAttention, continuous batching, parallelism, and OpenAI-compatible APIs. Tradeoff: Throughput and memory utilization improve, while model and kernel compatibility must be verified. Watch for: Benchmarking only one concurrency level hides latency-throughput collapse points.}
\paragraph{MCQ-0706 [hard]} The team proposes vLLM for a real workload. Which finding gives the correct decision boundary?
\begin{enumerate}[label=\Alph*.]
\item Ecosystem integration is strong, while supported optimization paths vary by model and hardware.
\item Throughput and memory utilization improve, while model and kernel compatibility must be verified.
\item Governance and rollback improve but approval gates can slow iteration.
\item A common control plane helps governance but default scheduling is not token-aware.
\end{enumerate}
\noindent\textbf{Answer.} Throughput and memory utilization improve, while model and kernel compatibility must be verified.
\par\textit{vLLM is a high-throughput inference server using PagedAttention, continuous batching, parallelism, and OpenAI-compatible APIs. Tradeoff: Throughput and memory utilization improve, while model and kernel compatibility must be verified. Watch for: Benchmarking only one concurrency level hides latency-throughput collapse points.}
\paragraph{MCQ-0707 [hard]} Before releasing vLLM, which failure deserves a dedicated regression test?
\begin{enumerate}[label=\Alph*.]
\item Using mutable stage names without immutable digests makes rollbacks uncertain.
\item Requesting whole GPUs for tiny models can strand capacity.
\item Benchmarking only one concurrency level hides latency-throughput collapse points.
\item Enabling quantization without quality and kernel tests can reduce both accuracy and speed.
\end{enumerate}
\noindent\textbf{Answer.} Benchmarking only one concurrency level hides latency-throughput collapse points.
\par\textit{vLLM is a high-throughput inference server using PagedAttention, continuous batching, parallelism, and OpenAI-compatible APIs. Tradeoff: Throughput and memory utilization improve, while model and kernel compatibility must be verified. Watch for: Benchmarking only one concurrency level hides latency-throughput collapse points.}
\paragraph{MCQ-0708 [medium]} A design review reveals this mistake: Benchmarking only one concurrency level hides latency-throughput collapse points. Which mechanism should be traced first?
\begin{enumerate}[label=\Alph*.]
\item Kubernetes and GPUs
\item vLLM
\item Text Generation Inference
\item Model registry
\end{enumerate}
\noindent\textbf{Answer.} vLLM
\par\textit{vLLM is a high-throughput inference server using PagedAttention, continuous batching, parallelism, and OpenAI-compatible APIs. Tradeoff: Throughput and memory utilization improve, while model and kernel compatibility must be verified. Watch for: Benchmarking only one concurrency level hides latency-throughput collapse points.}
\paragraph{MCQ-0709 [hard]} Which causal explanation of vLLM connects what it does to when it is worth its cost?
\begin{enumerate}[label=\Alph*.]
\item Mechanism: Kubernetes schedules containers and device resources while operators add GPU sharing, queues, and autoscaling. Consequence: A common control plane helps governance but default scheduling is not token-aware.
\item Mechanism: TGI is a Hugging Face server for transformer inference with streaming, batching, quantization, and distributed support. Consequence: Ecosystem integration is strong, while supported optimization paths vary by model and hardware.
\item Mechanism: vLLM is a high-throughput inference server using PagedAttention, continuous batching, parallelism, and OpenAI-compatible APIs. Consequence: Throughput and memory utilization improve, while model and kernel compatibility must be verified.
\item Mechanism: A registry versions model artifacts, lineage, signatures, evaluation evidence, aliases, and promotion state. Consequence: Governance and rollback improve but approval gates can slow iteration.
\end{enumerate}
\noindent\textbf{Answer.} Mechanism: vLLM is a high-throughput inference server using PagedAttention, continuous batching, parallelism, and OpenAI-compatible APIs. Consequence: Throughput and memory utilization improve, while model and kernel compatibility must be verified.
\par\textit{vLLM is a high-throughput inference server using PagedAttention, continuous batching, parallelism, and OpenAI-compatible APIs. Tradeoff: Throughput and memory utilization improve, while model and kernel compatibility must be verified. Watch for: Benchmarking only one concurrency level hides latency-throughput collapse points.}
\paragraph{FC-0710 [medium]} Define vLLM in interview-ready language.
\noindent\textbf{Answer.} vLLM is a high-throughput inference server using PagedAttention, continuous batching, parallelism, and OpenAI-compatible APIs.
\par\textit{vLLM is a high-throughput inference server using PagedAttention, continuous batching, parallelism, and OpenAI-compatible APIs.}
\paragraph{FC-0711 [hard]} State the main tradeoff and production pitfall for vLLM.
\noindent\textbf{Answer.} Tradeoff: Throughput and memory utilization improve, while model and kernel compatibility must be verified. Pitfall: Benchmarking only one concurrency level hides latency-throughput collapse points.
\par\textit{Tradeoff: Throughput and memory utilization improve, while model and kernel compatibility must be verified. Pitfall: Benchmarking only one concurrency level hides latency-throughput collapse points.}
\paragraph{LA-0712 [hard]} Explain vLLM from first principles, then show how you would evaluate and operate it in production.
\noindent\textbf{Answer.} Start with the mechanism: vLLM is a high-throughput inference server using PagedAttention, continuous batching, parallelism, and OpenAI-compatible APIs. Make the tradeoff explicit: Throughput and memory utilization improve, while model and kernel compatibility must be verified. Define workload-specific offline metrics and a latency/cost budget, test important slices, instrument the critical path, and create a rollback criterion. The production review must explicitly guard against this failure: Benchmarking only one concurrency level hides latency-throughput collapse points.
\par\textit{A strong answer connects mechanism, measurable tradeoffs, failure isolation, observability, and rollback rather than listing product features.}
\paragraph{MCQ-0713 [medium]} A design requires this behavior: SGLang combines a structured generation runtime with optimized model serving, prefix reuse, batching, and parallelism. Which mechanism should the team study?
\begin{enumerate}[label=\Alph*.]
\item KServe
\item Kubernetes and GPUs
\item SGLang
\item NVIDIA Triton
\end{enumerate}
\noindent\textbf{Answer.} SGLang
\par\textit{SGLang combines a structured generation runtime with optimized model serving, prefix reuse, batching, and parallelism. Tradeoff: Tight runtime-serving integration can accelerate agentic workloads but changes application abstractions. Watch for: Assuming every structured program benefits equally from cache reuse leads to poor capacity models.}
\paragraph{MCQ-0714 [hard]} The team proposes SGLang for a real workload. Which finding gives the correct decision boundary?
\begin{enumerate}[label=\Alph*.]
\item Platform consistency improves, with Kubernetes operational overhead and cold-start concerns.
\item It standardizes diverse serving, though LLM-specific scheduling may require specialized backends.
\item Tight runtime-serving integration can accelerate agentic workloads but changes application abstractions.
\item A common control plane helps governance but default scheduling is not token-aware.
\end{enumerate}
\noindent\textbf{Answer.} Tight runtime-serving integration can accelerate agentic workloads but changes application abstractions.
\par\textit{SGLang combines a structured generation runtime with optimized model serving, prefix reuse, batching, and parallelism. Tradeoff: Tight runtime-serving integration can accelerate agentic workloads but changes application abstractions. Watch for: Assuming every structured program benefits equally from cache reuse leads to poor capacity models.}
\paragraph{MCQ-0715 [hard]} Before releasing SGLang, which failure deserves a dedicated regression test?
\begin{enumerate}[label=\Alph*.]
\item Using dynamic batching without queue-delay limits breaks interactive latency.
\item Requesting whole GPUs for tiny models can strand capacity.
\item Scaling only on CPU misses token queue pressure in GPU workloads.
\item Assuming every structured program benefits equally from cache reuse leads to poor capacity models.
\end{enumerate}
\noindent\textbf{Answer.} Assuming every structured program benefits equally from cache reuse leads to poor capacity models.
\par\textit{SGLang combines a structured generation runtime with optimized model serving, prefix reuse, batching, and parallelism. Tradeoff: Tight runtime-serving integration can accelerate agentic workloads but changes application abstractions. Watch for: Assuming every structured program benefits equally from cache reuse leads to poor capacity models.}
\paragraph{MCQ-0716 [medium]} A design review reveals this mistake: Assuming every structured program benefits equally from cache reuse leads to poor capacity models. Which mechanism should be traced first?
\begin{enumerate}[label=\Alph*.]
\item Kubernetes and GPUs
\item KServe
\item SGLang
\item NVIDIA Triton
\end{enumerate}
\noindent\textbf{Answer.} SGLang
\par\textit{SGLang combines a structured generation runtime with optimized model serving, prefix reuse, batching, and parallelism. Tradeoff: Tight runtime-serving integration can accelerate agentic workloads but changes application abstractions. Watch for: Assuming every structured program benefits equally from cache reuse leads to poor capacity models.}
\paragraph{MCQ-0717 [hard]} Which causal explanation of SGLang connects what it does to when it is worth its cost?
\begin{enumerate}[label=\Alph*.]
\item Mechanism: Kubernetes schedules containers and device resources while operators add GPU sharing, queues, and autoscaling. Consequence: A common control plane helps governance but default scheduling is not token-aware.
\item Mechanism: SGLang combines a structured generation runtime with optimized model serving, prefix reuse, batching, and parallelism. Consequence: Tight runtime-serving integration can accelerate agentic workloads but changes application abstractions.
\item Mechanism: Triton serves multiple model frameworks with dynamic batching, ensembles, metrics, and model repositories. Consequence: It standardizes diverse serving, though LLM-specific scheduling may require specialized backends.
\item Mechanism: KServe provides Kubernetes-native inference services, autoscaling, revisions, traffic splitting, and model runtimes. Consequence: Platform consistency improves, with Kubernetes operational overhead and cold-start concerns.
\end{enumerate}
\noindent\textbf{Answer.} Mechanism: SGLang combines a structured generation runtime with optimized model serving, prefix reuse, batching, and parallelism. Consequence: Tight runtime-serving integration can accelerate agentic workloads but changes application abstractions.
\par\textit{SGLang combines a structured generation runtime with optimized model serving, prefix reuse, batching, and parallelism. Tradeoff: Tight runtime-serving integration can accelerate agentic workloads but changes application abstractions. Watch for: Assuming every structured program benefits equally from cache reuse leads to poor capacity models.}
\paragraph{FC-0718 [medium]} Define SGLang in interview-ready language.
\noindent\textbf{Answer.} SGLang combines a structured generation runtime with optimized model serving, prefix reuse, batching, and parallelism.
\par\textit{SGLang combines a structured generation runtime with optimized model serving, prefix reuse, batching, and parallelism.}
\paragraph{FC-0719 [hard]} State the main tradeoff and production pitfall for SGLang.
\noindent\textbf{Answer.} Tradeoff: Tight runtime-serving integration can accelerate agentic workloads but changes application abstractions. Pitfall: Assuming every structured program benefits equally from cache reuse leads to poor capacity models.
\par\textit{Tradeoff: Tight runtime-serving integration can accelerate agentic workloads but changes application abstractions. Pitfall: Assuming every structured program benefits equally from cache reuse leads to poor capacity models.}
\paragraph{LA-0720 [hard]} Explain SGLang from first principles, then show how you would evaluate and operate it in production.
\noindent\textbf{Answer.} Start with the mechanism: SGLang combines a structured generation runtime with optimized model serving, prefix reuse, batching, and parallelism. Make the tradeoff explicit: Tight runtime-serving integration can accelerate agentic workloads but changes application abstractions. Define workload-specific offline metrics and a latency/cost budget, test important slices, instrument the critical path, and create a rollback criterion. The production review must explicitly guard against this failure: Assuming every structured program benefits equally from cache reuse leads to poor capacity models.
\par\textit{A strong answer connects mechanism, measurable tradeoffs, failure isolation, observability, and rollback rather than listing product features.}
\paragraph{MCQ-0721 [medium]} A design requires this behavior: TGI is a Hugging Face server for transformer inference with streaming, batching, quantization, and distributed support. Which mechanism should the team study?
\begin{enumerate}[label=\Alph*.]
\item Autoscaling for LLMs
\item Text Generation Inference
\item Model registry
\item BentoML
\end{enumerate}
\noindent\textbf{Answer.} Text Generation Inference
\par\textit{TGI is a Hugging Face server for transformer inference with streaming, batching, quantization, and distributed support. Tradeoff: Ecosystem integration is strong, while supported optimization paths vary by model and hardware. Watch for: Enabling quantization without quality and kernel tests can reduce both accuracy and speed.}
\paragraph{MCQ-0722 [hard]} The team proposes Text Generation Inference for a real workload. Which finding gives the correct decision boundary?
\begin{enumerate}[label=\Alph*.]
\item Governance and rollback improve but approval gates can slow iteration.
\item Extra replicas reduce latency but GPUs are costly and slow to warm.
\item Developer workflow is streamlined, while platform integration and runtime tuning still matter.
\item Ecosystem integration is strong, while supported optimization paths vary by model and hardware.
\end{enumerate}
\noindent\textbf{Answer.} Ecosystem integration is strong, while supported optimization paths vary by model and hardware.
\par\textit{TGI is a Hugging Face server for transformer inference with streaming, batching, quantization, and distributed support. Tradeoff: Ecosystem integration is strong, while supported optimization paths vary by model and hardware. Watch for: Enabling quantization without quality and kernel tests can reduce both accuracy and speed.}
\paragraph{MCQ-0723 [hard]} Before releasing Text Generation Inference, which failure deserves a dedicated regression test?
\begin{enumerate}[label=\Alph*.]
\item Using mutable stage names without immutable digests makes rollbacks uncertain.
\item Treating packaging success as load readiness skips concurrency tests.
\item Enabling quantization without quality and kernel tests can reduce both accuracy and speed.
\item CPU-based autoscaling reacts late or incorrectly to GPU saturation.
\end{enumerate}
\noindent\textbf{Answer.} Enabling quantization without quality and kernel tests can reduce both accuracy and speed.
\par\textit{TGI is a Hugging Face server for transformer inference with streaming, batching, quantization, and distributed support. Tradeoff: Ecosystem integration is strong, while supported optimization paths vary by model and hardware. Watch for: Enabling quantization without quality and kernel tests can reduce both accuracy and speed.}
\paragraph{MCQ-0724 [medium]} A design review reveals this mistake: Enabling quantization without quality and kernel tests can reduce both accuracy and speed. Which mechanism should be traced first?
\begin{enumerate}[label=\Alph*.]
\item Model registry
\item Autoscaling for LLMs
\item BentoML
\item Text Generation Inference
\end{enumerate}
\noindent\textbf{Answer.} Text Generation Inference
\par\textit{TGI is a Hugging Face server for transformer inference with streaming, batching, quantization, and distributed support. Tradeoff: Ecosystem integration is strong, while supported optimization paths vary by model and hardware. Watch for: Enabling quantization without quality and kernel tests can reduce both accuracy and speed.}
\paragraph{MCQ-0725 [hard]} Which causal explanation of Text Generation Inference connects what it does to when it is worth its cost?
\begin{enumerate}[label=\Alph*.]
\item Mechanism: BentoML packages model services, dependencies, APIs, and deployment configuration with adaptive batching. Consequence: Developer workflow is streamlined, while platform integration and runtime tuning still matter.
\item Mechanism: LLM autoscaling should consider queue depth, time-to-first-token, tokens per second, memory, and startup time. Consequence: Extra replicas reduce latency but GPUs are costly and slow to warm.
\item Mechanism: TGI is a Hugging Face server for transformer inference with streaming, batching, quantization, and distributed support. Consequence: Ecosystem integration is strong, while supported optimization paths vary by model and hardware.
\item Mechanism: A registry versions model artifacts, lineage, signatures, evaluation evidence, aliases, and promotion state. Consequence: Governance and rollback improve but approval gates can slow iteration.
\end{enumerate}
\noindent\textbf{Answer.} Mechanism: TGI is a Hugging Face server for transformer inference with streaming, batching, quantization, and distributed support. Consequence: Ecosystem integration is strong, while supported optimization paths vary by model and hardware.
\par\textit{TGI is a Hugging Face server for transformer inference with streaming, batching, quantization, and distributed support. Tradeoff: Ecosystem integration is strong, while supported optimization paths vary by model and hardware. Watch for: Enabling quantization without quality and kernel tests can reduce both accuracy and speed.}
\paragraph{FC-0726 [medium]} Define Text Generation Inference in interview-ready language.
\noindent\textbf{Answer.} TGI is a Hugging Face server for transformer inference with streaming, batching, quantization, and distributed support.
\par\textit{TGI is a Hugging Face server for transformer inference with streaming, batching, quantization, and distributed support.}
\paragraph{FC-0727 [hard]} State the main tradeoff and production pitfall for Text Generation Inference.
\noindent\textbf{Answer.} Tradeoff: Ecosystem integration is strong, while supported optimization paths vary by model and hardware. Pitfall: Enabling quantization without quality and kernel tests can reduce both accuracy and speed.
\par\textit{Tradeoff: Ecosystem integration is strong, while supported optimization paths vary by model and hardware. Pitfall: Enabling quantization without quality and kernel tests can reduce both accuracy and speed.}
\paragraph{LA-0728 [hard]} Explain Text Generation Inference from first principles, then show how you would evaluate and operate it in production.
\noindent\textbf{Answer.} Start with the mechanism: TGI is a Hugging Face server for transformer inference with streaming, batching, quantization, and distributed support. Make the tradeoff explicit: Ecosystem integration is strong, while supported optimization paths vary by model and hardware. Define workload-specific offline metrics and a latency/cost budget, test important slices, instrument the critical path, and create a rollback criterion. The production review must explicitly guard against this failure: Enabling quantization without quality and kernel tests can reduce both accuracy and speed.
\par\textit{A strong answer connects mechanism, measurable tradeoffs, failure isolation, observability, and rollback rather than listing product features.}
\paragraph{MCQ-0729 [medium]} A design requires this behavior: TensorRT-LLM compiles and serves NVIDIA-optimized LLM graphs and kernels across GPUs. Which mechanism should the team study?
\begin{enumerate}[label=\Alph*.]
\item NVIDIA Triton
\item Ray Serve
\item TensorRT-LLM
\item Canary and shadow releases
\end{enumerate}
\noindent\textbf{Answer.} TensorRT-LLM
\par\textit{TensorRT-LLM compiles and serves NVIDIA-optimized LLM graphs and kernels across GPUs. Tradeoff: Maximum NVIDIA performance comes with build complexity and hardware coupling. Watch for: Rebuilding engines for every shape without a deployment cache slows releases.}
\paragraph{MCQ-0730 [hard]} The team proposes TensorRT-LLM for a real workload. Which finding gives the correct decision boundary?
\begin{enumerate}[label=\Alph*.]
\item It standardizes diverse serving, though LLM-specific scheduling may require specialized backends.
\item Maximum NVIDIA performance comes with build complexity and hardware coupling.
\item Flexible distributed Python pipelines fit multi-model systems but require cluster and object-store discipline.
\item Risk falls, while cost, comparison logic, and privacy obligations grow.
\end{enumerate}
\noindent\textbf{Answer.} Maximum NVIDIA performance comes with build complexity and hardware coupling.
\par\textit{TensorRT-LLM compiles and serves NVIDIA-optimized LLM graphs and kernels across GPUs. Tradeoff: Maximum NVIDIA performance comes with build complexity and hardware coupling. Watch for: Rebuilding engines for every shape without a deployment cache slows releases.}
\paragraph{MCQ-0731 [hard]} Before releasing TensorRT-LLM, which failure deserves a dedicated regression test?
\begin{enumerate}[label=\Alph*.]
\item Comparing canaries with different traffic mixes confounds the result.
\item Rebuilding engines for every shape without a deployment cache slows releases.
\item Passing large tensors repeatedly through serialization paths wastes memory and latency.
\item Using dynamic batching without queue-delay limits breaks interactive latency.
\end{enumerate}
\noindent\textbf{Answer.} Rebuilding engines for every shape without a deployment cache slows releases.
\par\textit{TensorRT-LLM compiles and serves NVIDIA-optimized LLM graphs and kernels across GPUs. Tradeoff: Maximum NVIDIA performance comes with build complexity and hardware coupling. Watch for: Rebuilding engines for every shape without a deployment cache slows releases.}
\paragraph{MCQ-0732 [medium]} A design review reveals this mistake: Rebuilding engines for every shape without a deployment cache slows releases. Which mechanism should be traced first?
\begin{enumerate}[label=\Alph*.]
\item NVIDIA Triton
\item TensorRT-LLM
\item Canary and shadow releases
\item Ray Serve
\end{enumerate}
\noindent\textbf{Answer.} TensorRT-LLM
\par\textit{TensorRT-LLM compiles and serves NVIDIA-optimized LLM graphs and kernels across GPUs. Tradeoff: Maximum NVIDIA performance comes with build complexity and hardware coupling. Watch for: Rebuilding engines for every shape without a deployment cache slows releases.}
\paragraph{MCQ-0733 [hard]} Which causal explanation of TensorRT-LLM connects what it does to when it is worth its cost?
\begin{enumerate}[label=\Alph*.]
\item Mechanism: Triton serves multiple model frameworks with dynamic batching, ensembles, metrics, and model repositories. Consequence: It standardizes diverse serving, though LLM-specific scheduling may require specialized backends.
\item Mechanism: Canary traffic exposes a new version to a small cohort; shadowing duplicates requests without serving the new output. Consequence: Risk falls, while cost, comparison logic, and privacy obligations grow.
\item Mechanism: TensorRT-LLM compiles and serves NVIDIA-optimized LLM graphs and kernels across GPUs. Consequence: Maximum NVIDIA performance comes with build complexity and hardware coupling.
\item Mechanism: Ray Serve composes Python model deployments and scales replicas across a Ray cluster. Consequence: Flexible distributed Python pipelines fit multi-model systems but require cluster and object-store discipline.
\end{enumerate}
\noindent\textbf{Answer.} Mechanism: TensorRT-LLM compiles and serves NVIDIA-optimized LLM graphs and kernels across GPUs. Consequence: Maximum NVIDIA performance comes with build complexity and hardware coupling.
\par\textit{TensorRT-LLM compiles and serves NVIDIA-optimized LLM graphs and kernels across GPUs. Tradeoff: Maximum NVIDIA performance comes with build complexity and hardware coupling. Watch for: Rebuilding engines for every shape without a deployment cache slows releases.}
\paragraph{FC-0734 [medium]} Define TensorRT-LLM in interview-ready language.
\noindent\textbf{Answer.} TensorRT-LLM compiles and serves NVIDIA-optimized LLM graphs and kernels across GPUs.
\par\textit{TensorRT-LLM compiles and serves NVIDIA-optimized LLM graphs and kernels across GPUs.}
\paragraph{FC-0735 [hard]} State the main tradeoff and production pitfall for TensorRT-LLM.
\noindent\textbf{Answer.} Tradeoff: Maximum NVIDIA performance comes with build complexity and hardware coupling. Pitfall: Rebuilding engines for every shape without a deployment cache slows releases.
\par\textit{Tradeoff: Maximum NVIDIA performance comes with build complexity and hardware coupling. Pitfall: Rebuilding engines for every shape without a deployment cache slows releases.}
\paragraph{LA-0736 [hard]} Explain TensorRT-LLM from first principles, then show how you would evaluate and operate it in production.
\noindent\textbf{Answer.} Start with the mechanism: TensorRT-LLM compiles and serves NVIDIA-optimized LLM graphs and kernels across GPUs. Make the tradeoff explicit: Maximum NVIDIA performance comes with build complexity and hardware coupling. Define workload-specific offline metrics and a latency/cost budget, test important slices, instrument the critical path, and create a rollback criterion. The production review must explicitly guard against this failure: Rebuilding engines for every shape without a deployment cache slows releases.
\par\textit{A strong answer connects mechanism, measurable tradeoffs, failure isolation, observability, and rollback rather than listing product features.}
\paragraph{MCQ-0737 [medium]} A design requires this behavior: Triton serves multiple model frameworks with dynamic batching, ensembles, metrics, and model repositories. Which mechanism should the team study?
\begin{enumerate}[label=\Alph*.]
\item Model registry
\item Text Generation Inference
\item Canary and shadow releases
\item NVIDIA Triton
\end{enumerate}
\noindent\textbf{Answer.} NVIDIA Triton
\par\textit{Triton serves multiple model frameworks with dynamic batching, ensembles, metrics, and model repositories. Tradeoff: It standardizes diverse serving, though LLM-specific scheduling may require specialized backends. Watch for: Using dynamic batching without queue-delay limits breaks interactive latency.}
\paragraph{MCQ-0738 [hard]} The team proposes NVIDIA Triton for a real workload. Which finding gives the correct decision boundary?
\begin{enumerate}[label=\Alph*.]
\item Ecosystem integration is strong, while supported optimization paths vary by model and hardware.
\item Governance and rollback improve but approval gates can slow iteration.
\item Risk falls, while cost, comparison logic, and privacy obligations grow.
\item It standardizes diverse serving, though LLM-specific scheduling may require specialized backends.
\end{enumerate}
\noindent\textbf{Answer.} It standardizes diverse serving, though LLM-specific scheduling may require specialized backends.
\par\textit{Triton serves multiple model frameworks with dynamic batching, ensembles, metrics, and model repositories. Tradeoff: It standardizes diverse serving, though LLM-specific scheduling may require specialized backends. Watch for: Using dynamic batching without queue-delay limits breaks interactive latency.}
\paragraph{MCQ-0739 [hard]} Before releasing NVIDIA Triton, which failure deserves a dedicated regression test?
\begin{enumerate}[label=\Alph*.]
\item Comparing canaries with different traffic mixes confounds the result.
\item Using dynamic batching without queue-delay limits breaks interactive latency.
\item Enabling quantization without quality and kernel tests can reduce both accuracy and speed.
\item Using mutable stage names without immutable digests makes rollbacks uncertain.
\end{enumerate}
\noindent\textbf{Answer.} Using dynamic batching without queue-delay limits breaks interactive latency.
\par\textit{Triton serves multiple model frameworks with dynamic batching, ensembles, metrics, and model repositories. Tradeoff: It standardizes diverse serving, though LLM-specific scheduling may require specialized backends. Watch for: Using dynamic batching without queue-delay limits breaks interactive latency.}
\paragraph{MCQ-0740 [medium]} A design review reveals this mistake: Using dynamic batching without queue-delay limits breaks interactive latency. Which mechanism should be traced first?
\begin{enumerate}[label=\Alph*.]
\item NVIDIA Triton
\item Canary and shadow releases
\item Model registry
\item Text Generation Inference
\end{enumerate}
\noindent\textbf{Answer.} NVIDIA Triton
\par\textit{Triton serves multiple model frameworks with dynamic batching, ensembles, metrics, and model repositories. Tradeoff: It standardizes diverse serving, though LLM-specific scheduling may require specialized backends. Watch for: Using dynamic batching without queue-delay limits breaks interactive latency.}
\paragraph{MCQ-0741 [hard]} Which causal explanation of NVIDIA Triton connects what it does to when it is worth its cost?
\begin{enumerate}[label=\Alph*.]
\item Mechanism: Canary traffic exposes a new version to a small cohort; shadowing duplicates requests without serving the new output. Consequence: Risk falls, while cost, comparison logic, and privacy obligations grow.
\item Mechanism: TGI is a Hugging Face server for transformer inference with streaming, batching, quantization, and distributed support. Consequence: Ecosystem integration is strong, while supported optimization paths vary by model and hardware.
\item Mechanism: A registry versions model artifacts, lineage, signatures, evaluation evidence, aliases, and promotion state. Consequence: Governance and rollback improve but approval gates can slow iteration.
\item Mechanism: Triton serves multiple model frameworks with dynamic batching, ensembles, metrics, and model repositories. Consequence: It standardizes diverse serving, though LLM-specific scheduling may require specialized backends.
\end{enumerate}
\noindent\textbf{Answer.} Mechanism: Triton serves multiple model frameworks with dynamic batching, ensembles, metrics, and model repositories. Consequence: It standardizes diverse serving, though LLM-specific scheduling may require specialized backends.
\par\textit{Triton serves multiple model frameworks with dynamic batching, ensembles, metrics, and model repositories. Tradeoff: It standardizes diverse serving, though LLM-specific scheduling may require specialized backends. Watch for: Using dynamic batching without queue-delay limits breaks interactive latency.}
\paragraph{FC-0742 [medium]} Define NVIDIA Triton in interview-ready language.
\noindent\textbf{Answer.} Triton serves multiple model frameworks with dynamic batching, ensembles, metrics, and model repositories.
\par\textit{Triton serves multiple model frameworks with dynamic batching, ensembles, metrics, and model repositories.}
\paragraph{FC-0743 [hard]} State the main tradeoff and production pitfall for NVIDIA Triton.
\noindent\textbf{Answer.} Tradeoff: It standardizes diverse serving, though LLM-specific scheduling may require specialized backends. Pitfall: Using dynamic batching without queue-delay limits breaks interactive latency.
\par\textit{Tradeoff: It standardizes diverse serving, though LLM-specific scheduling may require specialized backends. Pitfall: Using dynamic batching without queue-delay limits breaks interactive latency.}
\paragraph{LA-0744 [hard]} Explain NVIDIA Triton from first principles, then show how you would evaluate and operate it in production.
\noindent\textbf{Answer.} Start with the mechanism: Triton serves multiple model frameworks with dynamic batching, ensembles, metrics, and model repositories. Make the tradeoff explicit: It standardizes diverse serving, though LLM-specific scheduling may require specialized backends. Define workload-specific offline metrics and a latency/cost budget, test important slices, instrument the critical path, and create a rollback criterion. The production review must explicitly guard against this failure: Using dynamic batching without queue-delay limits breaks interactive latency.
\par\textit{A strong answer connects mechanism, measurable tradeoffs, failure isolation, observability, and rollback rather than listing product features.}
\paragraph{MCQ-0745 [medium]} A design requires this behavior: llama.cpp runs quantized models efficiently on CPUs and varied local accelerators through a portable C++ stack. Which mechanism should the team study?
\begin{enumerate}[label=\Alph*.]
\item Canary and shadow releases
\item KServe
\item TensorRT-LLM
\item llama.cpp
\end{enumerate}
\noindent\textbf{Answer.} llama.cpp
\par\textit{llama.cpp runs quantized models efficiently on CPUs and varied local accelerators through a portable C++ stack. Tradeoff: Edge privacy and low dependency count trade off against datacenter throughput and model size. Watch for: Selecting a quantization by file size alone ignores accuracy and hardware kernels.}
\paragraph{MCQ-0746 [hard]} The team proposes llama.cpp for a real workload. Which finding gives the correct decision boundary?
\begin{enumerate}[label=\Alph*.]
\item Maximum NVIDIA performance comes with build complexity and hardware coupling.
\item Risk falls, while cost, comparison logic, and privacy obligations grow.
\item Edge privacy and low dependency count trade off against datacenter throughput and model size.
\item Platform consistency improves, with Kubernetes operational overhead and cold-start concerns.
\end{enumerate}
\noindent\textbf{Answer.} Edge privacy and low dependency count trade off against datacenter throughput and model size.
\par\textit{llama.cpp runs quantized models efficiently on CPUs and varied local accelerators through a portable C++ stack. Tradeoff: Edge privacy and low dependency count trade off against datacenter throughput and model size. Watch for: Selecting a quantization by file size alone ignores accuracy and hardware kernels.}
\paragraph{MCQ-0747 [hard]} Before releasing llama.cpp, which failure deserves a dedicated regression test?
\begin{enumerate}[label=\Alph*.]
\item Scaling only on CPU misses token queue pressure in GPU workloads.
\item Selecting a quantization by file size alone ignores accuracy and hardware kernels.
\item Comparing canaries with different traffic mixes confounds the result.
\item Rebuilding engines for every shape without a deployment cache slows releases.
\end{enumerate}
\noindent\textbf{Answer.} Selecting a quantization by file size alone ignores accuracy and hardware kernels.
\par\textit{llama.cpp runs quantized models efficiently on CPUs and varied local accelerators through a portable C++ stack. Tradeoff: Edge privacy and low dependency count trade off against datacenter throughput and model size. Watch for: Selecting a quantization by file size alone ignores accuracy and hardware kernels.}
\paragraph{MCQ-0748 [medium]} A design review reveals this mistake: Selecting a quantization by file size alone ignores accuracy and hardware kernels. Which mechanism should be traced first?
\begin{enumerate}[label=\Alph*.]
\item llama.cpp
\item Canary and shadow releases
\item TensorRT-LLM
\item KServe
\end{enumerate}
\noindent\textbf{Answer.} llama.cpp
\par\textit{llama.cpp runs quantized models efficiently on CPUs and varied local accelerators through a portable C++ stack. Tradeoff: Edge privacy and low dependency count trade off against datacenter throughput and model size. Watch for: Selecting a quantization by file size alone ignores accuracy and hardware kernels.}
\paragraph{MCQ-0749 [hard]} Which causal explanation of llama.cpp connects what it does to when it is worth its cost?
\begin{enumerate}[label=\Alph*.]
\item Mechanism: KServe provides Kubernetes-native inference services, autoscaling, revisions, traffic splitting, and model runtimes. Consequence: Platform consistency improves, with Kubernetes operational overhead and cold-start concerns.
\item Mechanism: TensorRT-LLM compiles and serves NVIDIA-optimized LLM graphs and kernels across GPUs. Consequence: Maximum NVIDIA performance comes with build complexity and hardware coupling.
\item Mechanism: llama.cpp runs quantized models efficiently on CPUs and varied local accelerators through a portable C++ stack. Consequence: Edge privacy and low dependency count trade off against datacenter throughput and model size.
\item Mechanism: Canary traffic exposes a new version to a small cohort; shadowing duplicates requests without serving the new output. Consequence: Risk falls, while cost, comparison logic, and privacy obligations grow.
\end{enumerate}
\noindent\textbf{Answer.} Mechanism: llama.cpp runs quantized models efficiently on CPUs and varied local accelerators through a portable C++ stack. Consequence: Edge privacy and low dependency count trade off against datacenter throughput and model size.
\par\textit{llama.cpp runs quantized models efficiently on CPUs and varied local accelerators through a portable C++ stack. Tradeoff: Edge privacy and low dependency count trade off against datacenter throughput and model size. Watch for: Selecting a quantization by file size alone ignores accuracy and hardware kernels.}
\paragraph{FC-0750 [medium]} Define llama.cpp in interview-ready language.
\noindent\textbf{Answer.} llama.cpp runs quantized models efficiently on CPUs and varied local accelerators through a portable C++ stack.
\par\textit{llama.cpp runs quantized models efficiently on CPUs and varied local accelerators through a portable C++ stack.}
\paragraph{FC-0751 [hard]} State the main tradeoff and production pitfall for llama.cpp.
\noindent\textbf{Answer.} Tradeoff: Edge privacy and low dependency count trade off against datacenter throughput and model size. Pitfall: Selecting a quantization by file size alone ignores accuracy and hardware kernels.
\par\textit{Tradeoff: Edge privacy and low dependency count trade off against datacenter throughput and model size. Pitfall: Selecting a quantization by file size alone ignores accuracy and hardware kernels.}
\paragraph{LA-0752 [hard]} Explain llama.cpp from first principles, then show how you would evaluate and operate it in production.
\noindent\textbf{Answer.} Start with the mechanism: llama.cpp runs quantized models efficiently on CPUs and varied local accelerators through a portable C++ stack. Make the tradeoff explicit: Edge privacy and low dependency count trade off against datacenter throughput and model size. Define workload-specific offline metrics and a latency/cost budget, test important slices, instrument the critical path, and create a rollback criterion. The production review must explicitly guard against this failure: Selecting a quantization by file size alone ignores accuracy and hardware kernels.
\par\textit{A strong answer connects mechanism, measurable tradeoffs, failure isolation, observability, and rollback rather than listing product features.}
\paragraph{MCQ-0753 [medium]} A design requires this behavior: KServe provides Kubernetes-native inference services, autoscaling, revisions, traffic splitting, and model runtimes. Which mechanism should the team study?
\begin{enumerate}[label=\Alph*.]
\item Model registry
\item vLLM
\item KServe
\item BentoML
\end{enumerate}
\noindent\textbf{Answer.} KServe
\par\textit{KServe provides Kubernetes-native inference services, autoscaling, revisions, traffic splitting, and model runtimes. Tradeoff: Platform consistency improves, with Kubernetes operational overhead and cold-start concerns. Watch for: Scaling only on CPU misses token queue pressure in GPU workloads.}
\paragraph{MCQ-0754 [hard]} The team proposes KServe for a real workload. Which finding gives the correct decision boundary?
\begin{enumerate}[label=\Alph*.]
\item Developer workflow is streamlined, while platform integration and runtime tuning still matter.
\item Throughput and memory utilization improve, while model and kernel compatibility must be verified.
\item Platform consistency improves, with Kubernetes operational overhead and cold-start concerns.
\item Governance and rollback improve but approval gates can slow iteration.
\end{enumerate}
\noindent\textbf{Answer.} Platform consistency improves, with Kubernetes operational overhead and cold-start concerns.
\par\textit{KServe provides Kubernetes-native inference services, autoscaling, revisions, traffic splitting, and model runtimes. Tradeoff: Platform consistency improves, with Kubernetes operational overhead and cold-start concerns. Watch for: Scaling only on CPU misses token queue pressure in GPU workloads.}
\paragraph{MCQ-0755 [hard]} Before releasing KServe, which failure deserves a dedicated regression test?
\begin{enumerate}[label=\Alph*.]
\item Scaling only on CPU misses token queue pressure in GPU workloads.
\item Using mutable stage names without immutable digests makes rollbacks uncertain.
\item Benchmarking only one concurrency level hides latency-throughput collapse points.
\item Treating packaging success as load readiness skips concurrency tests.
\end{enumerate}
\noindent\textbf{Answer.} Scaling only on CPU misses token queue pressure in GPU workloads.
\par\textit{KServe provides Kubernetes-native inference services, autoscaling, revisions, traffic splitting, and model runtimes. Tradeoff: Platform consistency improves, with Kubernetes operational overhead and cold-start concerns. Watch for: Scaling only on CPU misses token queue pressure in GPU workloads.}
\paragraph{MCQ-0756 [medium]} A design review reveals this mistake: Scaling only on CPU misses token queue pressure in GPU workloads. Which mechanism should be traced first?
\begin{enumerate}[label=\Alph*.]
\item BentoML
\item KServe
\item vLLM
\item Model registry
\end{enumerate}
\noindent\textbf{Answer.} KServe
\par\textit{KServe provides Kubernetes-native inference services, autoscaling, revisions, traffic splitting, and model runtimes. Tradeoff: Platform consistency improves, with Kubernetes operational overhead and cold-start concerns. Watch for: Scaling only on CPU misses token queue pressure in GPU workloads.}
\paragraph{MCQ-0757 [hard]} Which causal explanation of KServe connects what it does to when it is worth its cost?
\begin{enumerate}[label=\Alph*.]
\item Mechanism: KServe provides Kubernetes-native inference services, autoscaling, revisions, traffic splitting, and model runtimes. Consequence: Platform consistency improves, with Kubernetes operational overhead and cold-start concerns.
\item Mechanism: BentoML packages model services, dependencies, APIs, and deployment configuration with adaptive batching. Consequence: Developer workflow is streamlined, while platform integration and runtime tuning still matter.
\item Mechanism: A registry versions model artifacts, lineage, signatures, evaluation evidence, aliases, and promotion state. Consequence: Governance and rollback improve but approval gates can slow iteration.
\item Mechanism: vLLM is a high-throughput inference server using PagedAttention, continuous batching, parallelism, and OpenAI-compatible APIs. Consequence: Throughput and memory utilization improve, while model and kernel compatibility must be verified.
\end{enumerate}
\noindent\textbf{Answer.} Mechanism: KServe provides Kubernetes-native inference services, autoscaling, revisions, traffic splitting, and model runtimes. Consequence: Platform consistency improves, with Kubernetes operational overhead and cold-start concerns.
\par\textit{KServe provides Kubernetes-native inference services, autoscaling, revisions, traffic splitting, and model runtimes. Tradeoff: Platform consistency improves, with Kubernetes operational overhead and cold-start concerns. Watch for: Scaling only on CPU misses token queue pressure in GPU workloads.}
\paragraph{FC-0758 [medium]} Define KServe in interview-ready language.
\noindent\textbf{Answer.} KServe provides Kubernetes-native inference services, autoscaling, revisions, traffic splitting, and model runtimes.
\par\textit{KServe provides Kubernetes-native inference services, autoscaling, revisions, traffic splitting, and model runtimes.}
\paragraph{FC-0759 [hard]} State the main tradeoff and production pitfall for KServe.
\noindent\textbf{Answer.} Tradeoff: Platform consistency improves, with Kubernetes operational overhead and cold-start concerns. Pitfall: Scaling only on CPU misses token queue pressure in GPU workloads.
\par\textit{Tradeoff: Platform consistency improves, with Kubernetes operational overhead and cold-start concerns. Pitfall: Scaling only on CPU misses token queue pressure in GPU workloads.}
\paragraph{LA-0760 [hard]} Explain KServe from first principles, then show how you would evaluate and operate it in production.
\noindent\textbf{Answer.} Start with the mechanism: KServe provides Kubernetes-native inference services, autoscaling, revisions, traffic splitting, and model runtimes. Make the tradeoff explicit: Platform consistency improves, with Kubernetes operational overhead and cold-start concerns. Define workload-specific offline metrics and a latency/cost budget, test important slices, instrument the critical path, and create a rollback criterion. The production review must explicitly guard against this failure: Scaling only on CPU misses token queue pressure in GPU workloads.
\par\textit{A strong answer connects mechanism, measurable tradeoffs, failure isolation, observability, and rollback rather than listing product features.}
\paragraph{MCQ-0761 [medium]} A design requires this behavior: Ray Serve composes Python model deployments and scales replicas across a Ray cluster. Which mechanism should the team study?
\begin{enumerate}[label=\Alph*.]
\item Canary and shadow releases
\item Ray Serve
\item Autoscaling for LLMs
\item vLLM
\end{enumerate}
\noindent\textbf{Answer.} Ray Serve
\par\textit{Ray Serve composes Python model deployments and scales replicas across a Ray cluster. Tradeoff: Flexible distributed Python pipelines fit multi-model systems but require cluster and object-store discipline. Watch for: Passing large tensors repeatedly through serialization paths wastes memory and latency.}
\paragraph{MCQ-0762 [hard]} The team proposes Ray Serve for a real workload. Which finding gives the correct decision boundary?
\begin{enumerate}[label=\Alph*.]
\item Flexible distributed Python pipelines fit multi-model systems but require cluster and object-store discipline.
\item Risk falls, while cost, comparison logic, and privacy obligations grow.
\item Extra replicas reduce latency but GPUs are costly and slow to warm.
\item Throughput and memory utilization improve, while model and kernel compatibility must be verified.
\end{enumerate}
\noindent\textbf{Answer.} Flexible distributed Python pipelines fit multi-model systems but require cluster and object-store discipline.
\par\textit{Ray Serve composes Python model deployments and scales replicas across a Ray cluster. Tradeoff: Flexible distributed Python pipelines fit multi-model systems but require cluster and object-store discipline. Watch for: Passing large tensors repeatedly through serialization paths wastes memory and latency.}
\paragraph{MCQ-0763 [hard]} Before releasing Ray Serve, which failure deserves a dedicated regression test?
\begin{enumerate}[label=\Alph*.]
\item Passing large tensors repeatedly through serialization paths wastes memory and latency.
\item Benchmarking only one concurrency level hides latency-throughput collapse points.
\item CPU-based autoscaling reacts late or incorrectly to GPU saturation.
\item Comparing canaries with different traffic mixes confounds the result.
\end{enumerate}
\noindent\textbf{Answer.} Passing large tensors repeatedly through serialization paths wastes memory and latency.
\par\textit{Ray Serve composes Python model deployments and scales replicas across a Ray cluster. Tradeoff: Flexible distributed Python pipelines fit multi-model systems but require cluster and object-store discipline. Watch for: Passing large tensors repeatedly through serialization paths wastes memory and latency.}
\paragraph{MCQ-0764 [medium]} A design review reveals this mistake: Passing large tensors repeatedly through serialization paths wastes memory and latency. Which mechanism should be traced first?
\begin{enumerate}[label=\Alph*.]
\item Ray Serve
\item vLLM
\item Autoscaling for LLMs
\item Canary and shadow releases
\end{enumerate}
\noindent\textbf{Answer.} Ray Serve
\par\textit{Ray Serve composes Python model deployments and scales replicas across a Ray cluster. Tradeoff: Flexible distributed Python pipelines fit multi-model systems but require cluster and object-store discipline. Watch for: Passing large tensors repeatedly through serialization paths wastes memory and latency.}
\paragraph{MCQ-0765 [hard]} Which causal explanation of Ray Serve connects what it does to when it is worth its cost?
\begin{enumerate}[label=\Alph*.]
\item Mechanism: Ray Serve composes Python model deployments and scales replicas across a Ray cluster. Consequence: Flexible distributed Python pipelines fit multi-model systems but require cluster and object-store discipline.
\item Mechanism: Canary traffic exposes a new version to a small cohort; shadowing duplicates requests without serving the new output. Consequence: Risk falls, while cost, comparison logic, and privacy obligations grow.
\item Mechanism: LLM autoscaling should consider queue depth, time-to-first-token, tokens per second, memory, and startup time. Consequence: Extra replicas reduce latency but GPUs are costly and slow to warm.
\item Mechanism: vLLM is a high-throughput inference server using PagedAttention, continuous batching, parallelism, and OpenAI-compatible APIs. Consequence: Throughput and memory utilization improve, while model and kernel compatibility must be verified.
\end{enumerate}
\noindent\textbf{Answer.} Mechanism: Ray Serve composes Python model deployments and scales replicas across a Ray cluster. Consequence: Flexible distributed Python pipelines fit multi-model systems but require cluster and object-store discipline.
\par\textit{Ray Serve composes Python model deployments and scales replicas across a Ray cluster. Tradeoff: Flexible distributed Python pipelines fit multi-model systems but require cluster and object-store discipline. Watch for: Passing large tensors repeatedly through serialization paths wastes memory and latency.}
\paragraph{FC-0766 [medium]} Define Ray Serve in interview-ready language.
\noindent\textbf{Answer.} Ray Serve composes Python model deployments and scales replicas across a Ray cluster.
\par\textit{Ray Serve composes Python model deployments and scales replicas across a Ray cluster.}
\paragraph{FC-0767 [hard]} State the main tradeoff and production pitfall for Ray Serve.
\noindent\textbf{Answer.} Tradeoff: Flexible distributed Python pipelines fit multi-model systems but require cluster and object-store discipline. Pitfall: Passing large tensors repeatedly through serialization paths wastes memory and latency.
\par\textit{Tradeoff: Flexible distributed Python pipelines fit multi-model systems but require cluster and object-store discipline. Pitfall: Passing large tensors repeatedly through serialization paths wastes memory and latency.}
\paragraph{LA-0768 [hard]} Explain Ray Serve from first principles, then show how you would evaluate and operate it in production.
\noindent\textbf{Answer.} Start with the mechanism: Ray Serve composes Python model deployments and scales replicas across a Ray cluster. Make the tradeoff explicit: Flexible distributed Python pipelines fit multi-model systems but require cluster and object-store discipline. Define workload-specific offline metrics and a latency/cost budget, test important slices, instrument the critical path, and create a rollback criterion. The production review must explicitly guard against this failure: Passing large tensors repeatedly through serialization paths wastes memory and latency.
\par\textit{A strong answer connects mechanism, measurable tradeoffs, failure isolation, observability, and rollback rather than listing product features.}
\paragraph{MCQ-0769 [medium]} A design requires this behavior: BentoML packages model services, dependencies, APIs, and deployment configuration with adaptive batching. Which mechanism should the team study?
\begin{enumerate}[label=\Alph*.]
\item SGLang
\item KServe
\item Autoscaling for LLMs
\item BentoML
\end{enumerate}
\noindent\textbf{Answer.} BentoML
\par\textit{BentoML packages model services, dependencies, APIs, and deployment configuration with adaptive batching. Tradeoff: Developer workflow is streamlined, while platform integration and runtime tuning still matter. Watch for: Treating packaging success as load readiness skips concurrency tests.}
\paragraph{MCQ-0770 [hard]} The team proposes BentoML for a real workload. Which finding gives the correct decision boundary?
\begin{enumerate}[label=\Alph*.]
\item Platform consistency improves, with Kubernetes operational overhead and cold-start concerns.
\item Developer workflow is streamlined, while platform integration and runtime tuning still matter.
\item Extra replicas reduce latency but GPUs are costly and slow to warm.
\item Tight runtime-serving integration can accelerate agentic workloads but changes application abstractions.
\end{enumerate}
\noindent\textbf{Answer.} Developer workflow is streamlined, while platform integration and runtime tuning still matter.
\par\textit{BentoML packages model services, dependencies, APIs, and deployment configuration with adaptive batching. Tradeoff: Developer workflow is streamlined, while platform integration and runtime tuning still matter. Watch for: Treating packaging success as load readiness skips concurrency tests.}
\paragraph{MCQ-0771 [hard]} Before releasing BentoML, which failure deserves a dedicated regression test?
\begin{enumerate}[label=\Alph*.]
\item Scaling only on CPU misses token queue pressure in GPU workloads.
\item Assuming every structured program benefits equally from cache reuse leads to poor capacity models.
\item CPU-based autoscaling reacts late or incorrectly to GPU saturation.
\item Treating packaging success as load readiness skips concurrency tests.
\end{enumerate}
\noindent\textbf{Answer.} Treating packaging success as load readiness skips concurrency tests.
\par\textit{BentoML packages model services, dependencies, APIs, and deployment configuration with adaptive batching. Tradeoff: Developer workflow is streamlined, while platform integration and runtime tuning still matter. Watch for: Treating packaging success as load readiness skips concurrency tests.}
\paragraph{MCQ-0772 [medium]} A design review reveals this mistake: Treating packaging success as load readiness skips concurrency tests. Which mechanism should be traced first?
\begin{enumerate}[label=\Alph*.]
\item BentoML
\item Autoscaling for LLMs
\item KServe
\item SGLang
\end{enumerate}
\noindent\textbf{Answer.} BentoML
\par\textit{BentoML packages model services, dependencies, APIs, and deployment configuration with adaptive batching. Tradeoff: Developer workflow is streamlined, while platform integration and runtime tuning still matter. Watch for: Treating packaging success as load readiness skips concurrency tests.}
\paragraph{MCQ-0773 [hard]} Which causal explanation of BentoML connects what it does to when it is worth its cost?
\begin{enumerate}[label=\Alph*.]
\item Mechanism: BentoML packages model services, dependencies, APIs, and deployment configuration with adaptive batching. Consequence: Developer workflow is streamlined, while platform integration and runtime tuning still matter.
\item Mechanism: LLM autoscaling should consider queue depth, time-to-first-token, tokens per second, memory, and startup time. Consequence: Extra replicas reduce latency but GPUs are costly and slow to warm.
\item Mechanism: SGLang combines a structured generation runtime with optimized model serving, prefix reuse, batching, and parallelism. Consequence: Tight runtime-serving integration can accelerate agentic workloads but changes application abstractions.
\item Mechanism: KServe provides Kubernetes-native inference services, autoscaling, revisions, traffic splitting, and model runtimes. Consequence: Platform consistency improves, with Kubernetes operational overhead and cold-start concerns.
\end{enumerate}
\noindent\textbf{Answer.} Mechanism: BentoML packages model services, dependencies, APIs, and deployment configuration with adaptive batching. Consequence: Developer workflow is streamlined, while platform integration and runtime tuning still matter.
\par\textit{BentoML packages model services, dependencies, APIs, and deployment configuration with adaptive batching. Tradeoff: Developer workflow is streamlined, while platform integration and runtime tuning still matter. Watch for: Treating packaging success as load readiness skips concurrency tests.}
\paragraph{FC-0774 [medium]} Define BentoML in interview-ready language.
\noindent\textbf{Answer.} BentoML packages model services, dependencies, APIs, and deployment configuration with adaptive batching.
\par\textit{BentoML packages model services, dependencies, APIs, and deployment configuration with adaptive batching.}
\paragraph{FC-0775 [hard]} State the main tradeoff and production pitfall for BentoML.
\noindent\textbf{Answer.} Tradeoff: Developer workflow is streamlined, while platform integration and runtime tuning still matter. Pitfall: Treating packaging success as load readiness skips concurrency tests.
\par\textit{Tradeoff: Developer workflow is streamlined, while platform integration and runtime tuning still matter. Pitfall: Treating packaging success as load readiness skips concurrency tests.}
\paragraph{LA-0776 [hard]} Explain BentoML from first principles, then show how you would evaluate and operate it in production.
\noindent\textbf{Answer.} Start with the mechanism: BentoML packages model services, dependencies, APIs, and deployment configuration with adaptive batching. Make the tradeoff explicit: Developer workflow is streamlined, while platform integration and runtime tuning still matter. Define workload-specific offline metrics and a latency/cost budget, test important slices, instrument the critical path, and create a rollback criterion. The production review must explicitly guard against this failure: Treating packaging success as load readiness skips concurrency tests.
\par\textit{A strong answer connects mechanism, measurable tradeoffs, failure isolation, observability, and rollback rather than listing product features.}
\paragraph{MCQ-0777 [medium]} A design requires this behavior: Kubernetes schedules containers and device resources while operators add GPU sharing, queues, and autoscaling. Which mechanism should the team study?
\begin{enumerate}[label=\Alph*.]
\item Kubernetes and GPUs
\item llama.cpp
\item vLLM
\item Ray Serve
\end{enumerate}
\noindent\textbf{Answer.} Kubernetes and GPUs
\par\textit{Kubernetes schedules containers and device resources while operators add GPU sharing, queues, and autoscaling. Tradeoff: A common control plane helps governance but default scheduling is not token-aware. Watch for: Requesting whole GPUs for tiny models can strand capacity.}
\paragraph{MCQ-0778 [hard]} The team proposes Kubernetes and GPUs for a real workload. Which finding gives the correct decision boundary?
\begin{enumerate}[label=\Alph*.]
\item Edge privacy and low dependency count trade off against datacenter throughput and model size.
\item A common control plane helps governance but default scheduling is not token-aware.
\item Throughput and memory utilization improve, while model and kernel compatibility must be verified.
\item Flexible distributed Python pipelines fit multi-model systems but require cluster and object-store discipline.
\end{enumerate}
\noindent\textbf{Answer.} A common control plane helps governance but default scheduling is not token-aware.
\par\textit{Kubernetes schedules containers and device resources while operators add GPU sharing, queues, and autoscaling. Tradeoff: A common control plane helps governance but default scheduling is not token-aware. Watch for: Requesting whole GPUs for tiny models can strand capacity.}
\paragraph{MCQ-0779 [hard]} Before releasing Kubernetes and GPUs, which failure deserves a dedicated regression test?
\begin{enumerate}[label=\Alph*.]
\item Benchmarking only one concurrency level hides latency-throughput collapse points.
\item Selecting a quantization by file size alone ignores accuracy and hardware kernels.
\item Passing large tensors repeatedly through serialization paths wastes memory and latency.
\item Requesting whole GPUs for tiny models can strand capacity.
\end{enumerate}
\noindent\textbf{Answer.} Requesting whole GPUs for tiny models can strand capacity.
\par\textit{Kubernetes schedules containers and device resources while operators add GPU sharing, queues, and autoscaling. Tradeoff: A common control plane helps governance but default scheduling is not token-aware. Watch for: Requesting whole GPUs for tiny models can strand capacity.}
\paragraph{MCQ-0780 [medium]} A design review reveals this mistake: Requesting whole GPUs for tiny models can strand capacity. Which mechanism should be traced first?
\begin{enumerate}[label=\Alph*.]
\item vLLM
\item Kubernetes and GPUs
\item llama.cpp
\item Ray Serve
\end{enumerate}
\noindent\textbf{Answer.} Kubernetes and GPUs
\par\textit{Kubernetes schedules containers and device resources while operators add GPU sharing, queues, and autoscaling. Tradeoff: A common control plane helps governance but default scheduling is not token-aware. Watch for: Requesting whole GPUs for tiny models can strand capacity.}
\paragraph{MCQ-0781 [hard]} Which causal explanation of Kubernetes and GPUs connects what it does to when it is worth its cost?
\begin{enumerate}[label=\Alph*.]
\item Mechanism: Ray Serve composes Python model deployments and scales replicas across a Ray cluster. Consequence: Flexible distributed Python pipelines fit multi-model systems but require cluster and object-store discipline.
\item Mechanism: vLLM is a high-throughput inference server using PagedAttention, continuous batching, parallelism, and OpenAI-compatible APIs. Consequence: Throughput and memory utilization improve, while model and kernel compatibility must be verified.
\item Mechanism: llama.cpp runs quantized models efficiently on CPUs and varied local accelerators through a portable C++ stack. Consequence: Edge privacy and low dependency count trade off against datacenter throughput and model size.
\item Mechanism: Kubernetes schedules containers and device resources while operators add GPU sharing, queues, and autoscaling. Consequence: A common control plane helps governance but default scheduling is not token-aware.
\end{enumerate}
\noindent\textbf{Answer.} Mechanism: Kubernetes schedules containers and device resources while operators add GPU sharing, queues, and autoscaling. Consequence: A common control plane helps governance but default scheduling is not token-aware.
\par\textit{Kubernetes schedules containers and device resources while operators add GPU sharing, queues, and autoscaling. Tradeoff: A common control plane helps governance but default scheduling is not token-aware. Watch for: Requesting whole GPUs for tiny models can strand capacity.}
\paragraph{FC-0782 [medium]} Define Kubernetes and GPUs in interview-ready language.
\noindent\textbf{Answer.} Kubernetes schedules containers and device resources while operators add GPU sharing, queues, and autoscaling.
\par\textit{Kubernetes schedules containers and device resources while operators add GPU sharing, queues, and autoscaling.}
\paragraph{FC-0783 [hard]} State the main tradeoff and production pitfall for Kubernetes and GPUs.
\noindent\textbf{Answer.} Tradeoff: A common control plane helps governance but default scheduling is not token-aware. Pitfall: Requesting whole GPUs for tiny models can strand capacity.
\par\textit{Tradeoff: A common control plane helps governance but default scheduling is not token-aware. Pitfall: Requesting whole GPUs for tiny models can strand capacity.}
\paragraph{LA-0784 [hard]} Explain Kubernetes and GPUs from first principles, then show how you would evaluate and operate it in production.
\noindent\textbf{Answer.} Start with the mechanism: Kubernetes schedules containers and device resources while operators add GPU sharing, queues, and autoscaling. Make the tradeoff explicit: A common control plane helps governance but default scheduling is not token-aware. Define workload-specific offline metrics and a latency/cost budget, test important slices, instrument the critical path, and create a rollback criterion. The production review must explicitly guard against this failure: Requesting whole GPUs for tiny models can strand capacity.
\par\textit{A strong answer connects mechanism, measurable tradeoffs, failure isolation, observability, and rollback rather than listing product features.}
\paragraph{MCQ-0785 [medium]} A design requires this behavior: Gateways normalize provider APIs and add routing, budgets, rate limits, fallbacks, caching, and audit policy. Which mechanism should the team study?
\begin{enumerate}[label=\Alph*.]
\item LLM gateways
\item SGLang
\item Text Generation Inference
\item vLLM
\end{enumerate}
\noindent\textbf{Answer.} LLM gateways
\par\textit{Gateways normalize provider APIs and add routing, budgets, rate limits, fallbacks, caching, and audit policy. Tradeoff: Central control improves governance but becomes a high-impact availability and security dependency. Watch for: Automatic fallback across non-equivalent models can silently change behavior.}
\paragraph{MCQ-0786 [hard]} The team proposes LLM gateways for a real workload. Which finding gives the correct decision boundary?
\begin{enumerate}[label=\Alph*.]
\item Tight runtime-serving integration can accelerate agentic workloads but changes application abstractions.
\item Central control improves governance but becomes a high-impact availability and security dependency.
\item Throughput and memory utilization improve, while model and kernel compatibility must be verified.
\item Ecosystem integration is strong, while supported optimization paths vary by model and hardware.
\end{enumerate}
\noindent\textbf{Answer.} Central control improves governance but becomes a high-impact availability and security dependency.
\par\textit{Gateways normalize provider APIs and add routing, budgets, rate limits, fallbacks, caching, and audit policy. Tradeoff: Central control improves governance but becomes a high-impact availability and security dependency. Watch for: Automatic fallback across non-equivalent models can silently change behavior.}
\paragraph{MCQ-0787 [hard]} Before releasing LLM gateways, which failure deserves a dedicated regression test?
\begin{enumerate}[label=\Alph*.]
\item Automatic fallback across non-equivalent models can silently change behavior.
\item Assuming every structured program benefits equally from cache reuse leads to poor capacity models.
\item Enabling quantization without quality and kernel tests can reduce both accuracy and speed.
\item Benchmarking only one concurrency level hides latency-throughput collapse points.
\end{enumerate}
\noindent\textbf{Answer.} Automatic fallback across non-equivalent models can silently change behavior.
\par\textit{Gateways normalize provider APIs and add routing, budgets, rate limits, fallbacks, caching, and audit policy. Tradeoff: Central control improves governance but becomes a high-impact availability and security dependency. Watch for: Automatic fallback across non-equivalent models can silently change behavior.}
\paragraph{MCQ-0788 [medium]} A design review reveals this mistake: Automatic fallback across non-equivalent models can silently change behavior. Which mechanism should be traced first?
\begin{enumerate}[label=\Alph*.]
\item LLM gateways
\item vLLM
\item SGLang
\item Text Generation Inference
\end{enumerate}
\noindent\textbf{Answer.} LLM gateways
\par\textit{Gateways normalize provider APIs and add routing, budgets, rate limits, fallbacks, caching, and audit policy. Tradeoff: Central control improves governance but becomes a high-impact availability and security dependency. Watch for: Automatic fallback across non-equivalent models can silently change behavior.}
\paragraph{MCQ-0789 [hard]} Which causal explanation of LLM gateways connects what it does to when it is worth its cost?
\begin{enumerate}[label=\Alph*.]
\item Mechanism: SGLang combines a structured generation runtime with optimized model serving, prefix reuse, batching, and parallelism. Consequence: Tight runtime-serving integration can accelerate agentic workloads but changes application abstractions.
\item Mechanism: vLLM is a high-throughput inference server using PagedAttention, continuous batching, parallelism, and OpenAI-compatible APIs. Consequence: Throughput and memory utilization improve, while model and kernel compatibility must be verified.
\item Mechanism: TGI is a Hugging Face server for transformer inference with streaming, batching, quantization, and distributed support. Consequence: Ecosystem integration is strong, while supported optimization paths vary by model and hardware.
\item Mechanism: Gateways normalize provider APIs and add routing, budgets, rate limits, fallbacks, caching, and audit policy. Consequence: Central control improves governance but becomes a high-impact availability and security dependency.
\end{enumerate}
\noindent\textbf{Answer.} Mechanism: Gateways normalize provider APIs and add routing, budgets, rate limits, fallbacks, caching, and audit policy. Consequence: Central control improves governance but becomes a high-impact availability and security dependency.
\par\textit{Gateways normalize provider APIs and add routing, budgets, rate limits, fallbacks, caching, and audit policy. Tradeoff: Central control improves governance but becomes a high-impact availability and security dependency. Watch for: Automatic fallback across non-equivalent models can silently change behavior.}
\paragraph{FC-0790 [medium]} Define LLM gateways in interview-ready language.
\noindent\textbf{Answer.} Gateways normalize provider APIs and add routing, budgets, rate limits, fallbacks, caching, and audit policy.
\par\textit{Gateways normalize provider APIs and add routing, budgets, rate limits, fallbacks, caching, and audit policy.}
\paragraph{FC-0791 [hard]} State the main tradeoff and production pitfall for LLM gateways.
\noindent\textbf{Answer.} Tradeoff: Central control improves governance but becomes a high-impact availability and security dependency. Pitfall: Automatic fallback across non-equivalent models can silently change behavior.
\par\textit{Tradeoff: Central control improves governance but becomes a high-impact availability and security dependency. Pitfall: Automatic fallback across non-equivalent models can silently change behavior.}
\paragraph{LA-0792 [hard]} Explain LLM gateways from first principles, then show how you would evaluate and operate it in production.
\noindent\textbf{Answer.} Start with the mechanism: Gateways normalize provider APIs and add routing, budgets, rate limits, fallbacks, caching, and audit policy. Make the tradeoff explicit: Central control improves governance but becomes a high-impact availability and security dependency. Define workload-specific offline metrics and a latency/cost budget, test important slices, instrument the critical path, and create a rollback criterion. The production review must explicitly guard against this failure: Automatic fallback across non-equivalent models can silently change behavior.
\par\textit{A strong answer connects mechanism, measurable tradeoffs, failure isolation, observability, and rollback rather than listing product features.}
\paragraph{MCQ-0793 [medium]} A design requires this behavior: Canary traffic exposes a new version to a small cohort; shadowing duplicates requests without serving the new output. Which mechanism should the team study?
\begin{enumerate}[label=\Alph*.]
\item Canary and shadow releases
\item BentoML
\item KServe
\item llama.cpp
\end{enumerate}
\noindent\textbf{Answer.} Canary and shadow releases
\par\textit{Canary traffic exposes a new version to a small cohort; shadowing duplicates requests without serving the new output. Tradeoff: Risk falls, while cost, comparison logic, and privacy obligations grow. Watch for: Comparing canaries with different traffic mixes confounds the result.}
\paragraph{MCQ-0794 [hard]} The team proposes Canary and shadow releases for a real workload. Which finding gives the correct decision boundary?
\begin{enumerate}[label=\Alph*.]
\item Platform consistency improves, with Kubernetes operational overhead and cold-start concerns.
\item Edge privacy and low dependency count trade off against datacenter throughput and model size.
\item Risk falls, while cost, comparison logic, and privacy obligations grow.
\item Developer workflow is streamlined, while platform integration and runtime tuning still matter.
\end{enumerate}
\noindent\textbf{Answer.} Risk falls, while cost, comparison logic, and privacy obligations grow.
\par\textit{Canary traffic exposes a new version to a small cohort; shadowing duplicates requests without serving the new output. Tradeoff: Risk falls, while cost, comparison logic, and privacy obligations grow. Watch for: Comparing canaries with different traffic mixes confounds the result.}
\paragraph{MCQ-0795 [hard]} Before releasing Canary and shadow releases, which failure deserves a dedicated regression test?
\begin{enumerate}[label=\Alph*.]
\item Comparing canaries with different traffic mixes confounds the result.
\item Selecting a quantization by file size alone ignores accuracy and hardware kernels.
\item Scaling only on CPU misses token queue pressure in GPU workloads.
\item Treating packaging success as load readiness skips concurrency tests.
\end{enumerate}
\noindent\textbf{Answer.} Comparing canaries with different traffic mixes confounds the result.
\par\textit{Canary traffic exposes a new version to a small cohort; shadowing duplicates requests without serving the new output. Tradeoff: Risk falls, while cost, comparison logic, and privacy obligations grow. Watch for: Comparing canaries with different traffic mixes confounds the result.}
\paragraph{MCQ-0796 [medium]} A design review reveals this mistake: Comparing canaries with different traffic mixes confounds the result. Which mechanism should be traced first?
\begin{enumerate}[label=\Alph*.]
\item BentoML
\item llama.cpp
\item KServe
\item Canary and shadow releases
\end{enumerate}
\noindent\textbf{Answer.} Canary and shadow releases
\par\textit{Canary traffic exposes a new version to a small cohort; shadowing duplicates requests without serving the new output. Tradeoff: Risk falls, while cost, comparison logic, and privacy obligations grow. Watch for: Comparing canaries with different traffic mixes confounds the result.}
\paragraph{MCQ-0797 [hard]} Which causal explanation of Canary and shadow releases connects what it does to when it is worth its cost?
\begin{enumerate}[label=\Alph*.]
\item Mechanism: llama.cpp runs quantized models efficiently on CPUs and varied local accelerators through a portable C++ stack. Consequence: Edge privacy and low dependency count trade off against datacenter throughput and model size.
\item Mechanism: KServe provides Kubernetes-native inference services, autoscaling, revisions, traffic splitting, and model runtimes. Consequence: Platform consistency improves, with Kubernetes operational overhead and cold-start concerns.
\item Mechanism: Canary traffic exposes a new version to a small cohort; shadowing duplicates requests without serving the new output. Consequence: Risk falls, while cost, comparison logic, and privacy obligations grow.
\item Mechanism: BentoML packages model services, dependencies, APIs, and deployment configuration with adaptive batching. Consequence: Developer workflow is streamlined, while platform integration and runtime tuning still matter.
\end{enumerate}
\noindent\textbf{Answer.} Mechanism: Canary traffic exposes a new version to a small cohort; shadowing duplicates requests without serving the new output. Consequence: Risk falls, while cost, comparison logic, and privacy obligations grow.
\par\textit{Canary traffic exposes a new version to a small cohort; shadowing duplicates requests without serving the new output. Tradeoff: Risk falls, while cost, comparison logic, and privacy obligations grow. Watch for: Comparing canaries with different traffic mixes confounds the result.}
\paragraph{FC-0798 [medium]} Define Canary and shadow releases in interview-ready language.
\noindent\textbf{Answer.} Canary traffic exposes a new version to a small cohort; shadowing duplicates requests without serving the new output.
\par\textit{Canary traffic exposes a new version to a small cohort; shadowing duplicates requests without serving the new output.}
\paragraph{FC-0799 [hard]} State the main tradeoff and production pitfall for Canary and shadow releases.
\noindent\textbf{Answer.} Tradeoff: Risk falls, while cost, comparison logic, and privacy obligations grow. Pitfall: Comparing canaries with different traffic mixes confounds the result.
\par\textit{Tradeoff: Risk falls, while cost, comparison logic, and privacy obligations grow. Pitfall: Comparing canaries with different traffic mixes confounds the result.}
\paragraph{LA-0800 [hard]} Explain Canary and shadow releases from first principles, then show how you would evaluate and operate it in production.
\noindent\textbf{Answer.} Start with the mechanism: Canary traffic exposes a new version to a small cohort; shadowing duplicates requests without serving the new output. Make the tradeoff explicit: Risk falls, while cost, comparison logic, and privacy obligations grow. Define workload-specific offline metrics and a latency/cost budget, test important slices, instrument the critical path, and create a rollback criterion. The production review must explicitly guard against this failure: Comparing canaries with different traffic mixes confounds the result.
\par\textit{A strong answer connects mechanism, measurable tradeoffs, failure isolation, observability, and rollback rather than listing product features.}
\paragraph{MCQ-0801 [medium]} A design requires this behavior: A registry versions model artifacts, lineage, signatures, evaluation evidence, aliases, and promotion state. Which mechanism should the team study?
\begin{enumerate}[label=\Alph*.]
\item Canary and shadow releases
\item Ray Serve
\item TensorRT-LLM
\item Model registry
\end{enumerate}
\noindent\textbf{Answer.} Model registry
\par\textit{A registry versions model artifacts, lineage, signatures, evaluation evidence, aliases, and promotion state. Tradeoff: Governance and rollback improve but approval gates can slow iteration. Watch for: Using mutable stage names without immutable digests makes rollbacks uncertain.}
\paragraph{MCQ-0802 [hard]} The team proposes Model registry for a real workload. Which finding gives the correct decision boundary?
\begin{enumerate}[label=\Alph*.]
\item Flexible distributed Python pipelines fit multi-model systems but require cluster and object-store discipline.
\item Risk falls, while cost, comparison logic, and privacy obligations grow.
\item Governance and rollback improve but approval gates can slow iteration.
\item Maximum NVIDIA performance comes with build complexity and hardware coupling.
\end{enumerate}
\noindent\textbf{Answer.} Governance and rollback improve but approval gates can slow iteration.
\par\textit{A registry versions model artifacts, lineage, signatures, evaluation evidence, aliases, and promotion state. Tradeoff: Governance and rollback improve but approval gates can slow iteration. Watch for: Using mutable stage names without immutable digests makes rollbacks uncertain.}
\paragraph{MCQ-0803 [hard]} Before releasing Model registry, which failure deserves a dedicated regression test?
\begin{enumerate}[label=\Alph*.]
\item Rebuilding engines for every shape without a deployment cache slows releases.
\item Comparing canaries with different traffic mixes confounds the result.
\item Passing large tensors repeatedly through serialization paths wastes memory and latency.
\item Using mutable stage names without immutable digests makes rollbacks uncertain.
\end{enumerate}
\noindent\textbf{Answer.} Using mutable stage names without immutable digests makes rollbacks uncertain.
\par\textit{A registry versions model artifacts, lineage, signatures, evaluation evidence, aliases, and promotion state. Tradeoff: Governance and rollback improve but approval gates can slow iteration. Watch for: Using mutable stage names without immutable digests makes rollbacks uncertain.}
\paragraph{MCQ-0804 [medium]} A design review reveals this mistake: Using mutable stage names without immutable digests makes rollbacks uncertain. Which mechanism should be traced first?
\begin{enumerate}[label=\Alph*.]
\item TensorRT-LLM
\item Model registry
\item Canary and shadow releases
\item Ray Serve
\end{enumerate}
\noindent\textbf{Answer.} Model registry
\par\textit{A registry versions model artifacts, lineage, signatures, evaluation evidence, aliases, and promotion state. Tradeoff: Governance and rollback improve but approval gates can slow iteration. Watch for: Using mutable stage names without immutable digests makes rollbacks uncertain.}
\paragraph{MCQ-0805 [hard]} Which causal explanation of Model registry connects what it does to when it is worth its cost?
\begin{enumerate}[label=\Alph*.]
\item Mechanism: Canary traffic exposes a new version to a small cohort; shadowing duplicates requests without serving the new output. Consequence: Risk falls, while cost, comparison logic, and privacy obligations grow.
\item Mechanism: A registry versions model artifacts, lineage, signatures, evaluation evidence, aliases, and promotion state. Consequence: Governance and rollback improve but approval gates can slow iteration.
\item Mechanism: TensorRT-LLM compiles and serves NVIDIA-optimized LLM graphs and kernels across GPUs. Consequence: Maximum NVIDIA performance comes with build complexity and hardware coupling.
\item Mechanism: Ray Serve composes Python model deployments and scales replicas across a Ray cluster. Consequence: Flexible distributed Python pipelines fit multi-model systems but require cluster and object-store discipline.
\end{enumerate}
\noindent\textbf{Answer.} Mechanism: A registry versions model artifacts, lineage, signatures, evaluation evidence, aliases, and promotion state. Consequence: Governance and rollback improve but approval gates can slow iteration.
\par\textit{A registry versions model artifacts, lineage, signatures, evaluation evidence, aliases, and promotion state. Tradeoff: Governance and rollback improve but approval gates can slow iteration. Watch for: Using mutable stage names without immutable digests makes rollbacks uncertain.}
\paragraph{FC-0806 [medium]} Define Model registry in interview-ready language.
\noindent\textbf{Answer.} A registry versions model artifacts, lineage, signatures, evaluation evidence, aliases, and promotion state.
\par\textit{A registry versions model artifacts, lineage, signatures, evaluation evidence, aliases, and promotion state.}
\paragraph{FC-0807 [hard]} State the main tradeoff and production pitfall for Model registry.
\noindent\textbf{Answer.} Tradeoff: Governance and rollback improve but approval gates can slow iteration. Pitfall: Using mutable stage names without immutable digests makes rollbacks uncertain.
\par\textit{Tradeoff: Governance and rollback improve but approval gates can slow iteration. Pitfall: Using mutable stage names without immutable digests makes rollbacks uncertain.}
\paragraph{LA-0808 [hard]} Explain Model registry from first principles, then show how you would evaluate and operate it in production.
\noindent\textbf{Answer.} Start with the mechanism: A registry versions model artifacts, lineage, signatures, evaluation evidence, aliases, and promotion state. Make the tradeoff explicit: Governance and rollback improve but approval gates can slow iteration. Define workload-specific offline metrics and a latency/cost budget, test important slices, instrument the critical path, and create a rollback criterion. The production review must explicitly guard against this failure: Using mutable stage names without immutable digests makes rollbacks uncertain.
\par\textit{A strong answer connects mechanism, measurable tradeoffs, failure isolation, observability, and rollback rather than listing product features.}
\paragraph{MCQ-0809 [medium]} A design requires this behavior: LLM autoscaling should consider queue depth, time-to-first-token, tokens per second, memory, and startup time. Which mechanism should the team study?
\begin{enumerate}[label=\Alph*.]
\item Model registry
\item Autoscaling for LLMs
\item LLM gateways
\item Canary and shadow releases
\end{enumerate}
\noindent\textbf{Answer.} Autoscaling for LLMs
\par\textit{LLM autoscaling should consider queue depth, time-to-first-token, tokens per second, memory, and startup time. Tradeoff: Extra replicas reduce latency but GPUs are costly and slow to warm. Watch for: CPU-based autoscaling reacts late or incorrectly to GPU saturation.}
\paragraph{MCQ-0810 [hard]} The team proposes Autoscaling for LLMs for a real workload. Which finding gives the correct decision boundary?
\begin{enumerate}[label=\Alph*.]
\item Risk falls, while cost, comparison logic, and privacy obligations grow.
\item Governance and rollback improve but approval gates can slow iteration.
\item Central control improves governance but becomes a high-impact availability and security dependency.
\item Extra replicas reduce latency but GPUs are costly and slow to warm.
\end{enumerate}
\noindent\textbf{Answer.} Extra replicas reduce latency but GPUs are costly and slow to warm.
\par\textit{LLM autoscaling should consider queue depth, time-to-first-token, tokens per second, memory, and startup time. Tradeoff: Extra replicas reduce latency but GPUs are costly and slow to warm. Watch for: CPU-based autoscaling reacts late or incorrectly to GPU saturation.}
\paragraph{MCQ-0811 [hard]} Before releasing Autoscaling for LLMs, which failure deserves a dedicated regression test?
\begin{enumerate}[label=\Alph*.]
\item Using mutable stage names without immutable digests makes rollbacks uncertain.
\item CPU-based autoscaling reacts late or incorrectly to GPU saturation.
\item Automatic fallback across non-equivalent models can silently change behavior.
\item Comparing canaries with different traffic mixes confounds the result.
\end{enumerate}
\noindent\textbf{Answer.} CPU-based autoscaling reacts late or incorrectly to GPU saturation.
\par\textit{LLM autoscaling should consider queue depth, time-to-first-token, tokens per second, memory, and startup time. Tradeoff: Extra replicas reduce latency but GPUs are costly and slow to warm. Watch for: CPU-based autoscaling reacts late or incorrectly to GPU saturation.}
\paragraph{MCQ-0812 [medium]} A design review reveals this mistake: CPU-based autoscaling reacts late or incorrectly to GPU saturation. Which mechanism should be traced first?
\begin{enumerate}[label=\Alph*.]
\item Autoscaling for LLMs
\item LLM gateways
\item Canary and shadow releases
\item Model registry
\end{enumerate}
\noindent\textbf{Answer.} Autoscaling for LLMs
\par\textit{LLM autoscaling should consider queue depth, time-to-first-token, tokens per second, memory, and startup time. Tradeoff: Extra replicas reduce latency but GPUs are costly and slow to warm. Watch for: CPU-based autoscaling reacts late or incorrectly to GPU saturation.}
\paragraph{MCQ-0813 [hard]} Which causal explanation of Autoscaling for LLMs connects what it does to when it is worth its cost?
\begin{enumerate}[label=\Alph*.]
\item Mechanism: Gateways normalize provider APIs and add routing, budgets, rate limits, fallbacks, caching, and audit policy. Consequence: Central control improves governance but becomes a high-impact availability and security dependency.
\item Mechanism: LLM autoscaling should consider queue depth, time-to-first-token, tokens per second, memory, and startup time. Consequence: Extra replicas reduce latency but GPUs are costly and slow to warm.
\item Mechanism: A registry versions model artifacts, lineage, signatures, evaluation evidence, aliases, and promotion state. Consequence: Governance and rollback improve but approval gates can slow iteration.
\item Mechanism: Canary traffic exposes a new version to a small cohort; shadowing duplicates requests without serving the new output. Consequence: Risk falls, while cost, comparison logic, and privacy obligations grow.
\end{enumerate}
\noindent\textbf{Answer.} Mechanism: LLM autoscaling should consider queue depth, time-to-first-token, tokens per second, memory, and startup time. Consequence: Extra replicas reduce latency but GPUs are costly and slow to warm.
\par\textit{LLM autoscaling should consider queue depth, time-to-first-token, tokens per second, memory, and startup time. Tradeoff: Extra replicas reduce latency but GPUs are costly and slow to warm. Watch for: CPU-based autoscaling reacts late or incorrectly to GPU saturation.}
\paragraph{FC-0814 [medium]} Define Autoscaling for LLMs in interview-ready language.
\noindent\textbf{Answer.} LLM autoscaling should consider queue depth, time-to-first-token, tokens per second, memory, and startup time.
\par\textit{LLM autoscaling should consider queue depth, time-to-first-token, tokens per second, memory, and startup time.}
\paragraph{FC-0815 [hard]} State the main tradeoff and production pitfall for Autoscaling for LLMs.
\noindent\textbf{Answer.} Tradeoff: Extra replicas reduce latency but GPUs are costly and slow to warm. Pitfall: CPU-based autoscaling reacts late or incorrectly to GPU saturation.
\par\textit{Tradeoff: Extra replicas reduce latency but GPUs are costly and slow to warm. Pitfall: CPU-based autoscaling reacts late or incorrectly to GPU saturation.}
\paragraph{LA-0816 [hard]} Explain Autoscaling for LLMs from first principles, then show how you would evaluate and operate it in production.
\noindent\textbf{Answer.} Start with the mechanism: LLM autoscaling should consider queue depth, time-to-first-token, tokens per second, memory, and startup time. Make the tradeoff explicit: Extra replicas reduce latency but GPUs are costly and slow to warm. Define workload-specific offline metrics and a latency/cost budget, test important slices, instrument the critical path, and create a rollback criterion. The production review must explicitly guard against this failure: CPU-based autoscaling reacts late or incorrectly to GPU saturation.
\par\textit{A strong answer connects mechanism, measurable tradeoffs, failure isolation, observability, and rollback rather than listing product features.}
\subsection{Data and ML Lifecycle}
\paragraph{MCQ-0817 [medium]} A design requires this behavior: DVC versions large data and model pointers with Git and defines reproducible pipeline stages and experiments. Which mechanism should the team study?
\begin{enumerate}[label=\Alph*.]
\item DVC
\item Delta Lake and Iceberg
\item Reproducible environments
\item Feature stores
\end{enumerate}
\noindent\textbf{Answer.} DVC
\par\textit{DVC versions large data and model pointers with Git and defines reproducible pipeline stages and experiments. Tradeoff: Familiar Git workflows improve reproducibility, while remote storage and cache discipline are required. Watch for: Committing only code while silently replacing remote data breaks lineage.}
\paragraph{MCQ-0818 [hard]} The team proposes DVC for a real workload. Which finding gives the correct decision boundary?
\begin{enumerate}[label=\Alph*.]
\item Reproducibility improves at the expense of build maintenance and sometimes performance.
\item Familiar Git workflows improve reproducibility, while remote storage and cache discipline are required.
\item Reliable analytics improve, while compaction and catalog operations become essential.
\item Training-serving consistency improves but adds infrastructure and ownership requirements.
\end{enumerate}
\noindent\textbf{Answer.} Familiar Git workflows improve reproducibility, while remote storage and cache discipline are required.
\par\textit{DVC versions large data and model pointers with Git and defines reproducible pipeline stages and experiments. Tradeoff: Familiar Git workflows improve reproducibility, while remote storage and cache discipline are required. Watch for: Committing only code while silently replacing remote data breaks lineage.}
\paragraph{MCQ-0819 [hard]} Before releasing DVC, which failure deserves a dedicated regression test?
\begin{enumerate}[label=\Alph*.]
\item A random seed alone cannot guarantee deterministic GPU execution.
\item Committing only code while silently replacing remote data breaks lineage.
\item Recomputing features differently online creates skew.
\item Excess small files destroy scan performance despite correct table semantics.
\end{enumerate}
\noindent\textbf{Answer.} Committing only code while silently replacing remote data breaks lineage.
\par\textit{DVC versions large data and model pointers with Git and defines reproducible pipeline stages and experiments. Tradeoff: Familiar Git workflows improve reproducibility, while remote storage and cache discipline are required. Watch for: Committing only code while silently replacing remote data breaks lineage.}
\paragraph{MCQ-0820 [medium]} A design review reveals this mistake: Committing only code while silently replacing remote data breaks lineage. Which mechanism should be traced first?
\begin{enumerate}[label=\Alph*.]
\item DVC
\item Reproducible environments
\item Feature stores
\item Delta Lake and Iceberg
\end{enumerate}
\noindent\textbf{Answer.} DVC
\par\textit{DVC versions large data and model pointers with Git and defines reproducible pipeline stages and experiments. Tradeoff: Familiar Git workflows improve reproducibility, while remote storage and cache discipline are required. Watch for: Committing only code while silently replacing remote data breaks lineage.}
\paragraph{MCQ-0821 [hard]} Which causal explanation of DVC connects what it does to when it is worth its cost?
\begin{enumerate}[label=\Alph*.]
\item Mechanism: Open table formats add schemas, snapshots, partition evolution, and transactional metadata over object storage. Consequence: Reliable analytics improve, while compaction and catalog operations become essential.
\item Mechanism: DVC versions large data and model pointers with Git and defines reproducible pipeline stages and experiments. Consequence: Familiar Git workflows improve reproducibility, while remote storage and cache discipline are required.
\item Mechanism: Lockfiles, containers, hardware metadata, seeds, and deterministic settings constrain run variability. Consequence: Reproducibility improves at the expense of build maintenance and sometimes performance.
\item Mechanism: Feature stores define reusable features and connect historical offline data with low-latency online serving. Consequence: Training-serving consistency improves but adds infrastructure and ownership requirements.
\end{enumerate}
\noindent\textbf{Answer.} Mechanism: DVC versions large data and model pointers with Git and defines reproducible pipeline stages and experiments. Consequence: Familiar Git workflows improve reproducibility, while remote storage and cache discipline are required.
\par\textit{DVC versions large data and model pointers with Git and defines reproducible pipeline stages and experiments. Tradeoff: Familiar Git workflows improve reproducibility, while remote storage and cache discipline are required. Watch for: Committing only code while silently replacing remote data breaks lineage.}
\paragraph{FC-0822 [medium]} Define DVC in interview-ready language.
\noindent\textbf{Answer.} DVC versions large data and model pointers with Git and defines reproducible pipeline stages and experiments.
\par\textit{DVC versions large data and model pointers with Git and defines reproducible pipeline stages and experiments.}
\paragraph{FC-0823 [hard]} State the main tradeoff and production pitfall for DVC.
\noindent\textbf{Answer.} Tradeoff: Familiar Git workflows improve reproducibility, while remote storage and cache discipline are required. Pitfall: Committing only code while silently replacing remote data breaks lineage.
\par\textit{Tradeoff: Familiar Git workflows improve reproducibility, while remote storage and cache discipline are required. Pitfall: Committing only code while silently replacing remote data breaks lineage.}
\paragraph{LA-0824 [hard]} Explain DVC from first principles, then show how you would evaluate and operate it in production.
\noindent\textbf{Answer.} Start with the mechanism: DVC versions large data and model pointers with Git and defines reproducible pipeline stages and experiments. Make the tradeoff explicit: Familiar Git workflows improve reproducibility, while remote storage and cache discipline are required. Define workload-specific offline metrics and a latency/cost budget, test important slices, instrument the critical path, and create a rollback criterion. The production review must explicitly guard against this failure: Committing only code while silently replacing remote data breaks lineage.
\par\textit{A strong answer connects mechanism, measurable tradeoffs, failure isolation, observability, and rollback rather than listing product features.}
\paragraph{MCQ-0825 [medium]} A design requires this behavior: lakeFS adds Git-like branches, commits, merges, and hooks over object-storage data lakes. Which mechanism should the team study?
\begin{enumerate}[label=\Alph*.]
\item Feature stores
\item Data quality testing
\item Experiment tracking
\item lakeFS
\end{enumerate}
\noindent\textbf{Answer.} lakeFS
\par\textit{lakeFS adds Git-like branches, commits, merges, and hooks over object-storage data lakes. Tradeoff: Atomic data-lake changes enable testing, with server and metadata operations to manage. Watch for: Treating object versioning alone as branch-level data semantics is insufficient.}
\paragraph{MCQ-0826 [hard]} The team proposes lakeFS for a real workload. Which finding gives the correct decision boundary?
\begin{enumerate}[label=\Alph*.]
\item Training-serving consistency improves but adds infrastructure and ownership requirements.
\item Comparison and reproducibility improve but inconsistent naming creates a junk drawer.
\item Atomic data-lake changes enable testing, with server and metadata operations to manage.
\item Early detection prevents downstream damage, while brittle thresholds create alert fatigue.
\end{enumerate}
\noindent\textbf{Answer.} Atomic data-lake changes enable testing, with server and metadata operations to manage.
\par\textit{lakeFS adds Git-like branches, commits, merges, and hooks over object-storage data lakes. Tradeoff: Atomic data-lake changes enable testing, with server and metadata operations to manage. Watch for: Treating object versioning alone as branch-level data semantics is insufficient.}
\paragraph{MCQ-0827 [hard]} Before releasing lakeFS, which failure deserves a dedicated regression test?
\begin{enumerate}[label=\Alph*.]
\item Recomputing features differently online creates skew.
\item Treating object versioning alone as branch-level data semantics is insufficient.
\item Logging a metric without dataset and code versions makes it uninterpretable.
\item Checking only aggregate distributions misses failures in important slices.
\end{enumerate}
\noindent\textbf{Answer.} Treating object versioning alone as branch-level data semantics is insufficient.
\par\textit{lakeFS adds Git-like branches, commits, merges, and hooks over object-storage data lakes. Tradeoff: Atomic data-lake changes enable testing, with server and metadata operations to manage. Watch for: Treating object versioning alone as branch-level data semantics is insufficient.}
\paragraph{MCQ-0828 [medium]} A design review reveals this mistake: Treating object versioning alone as branch-level data semantics is insufficient. Which mechanism should be traced first?
\begin{enumerate}[label=\Alph*.]
\item lakeFS
\item Experiment tracking
\item Feature stores
\item Data quality testing
\end{enumerate}
\noindent\textbf{Answer.} lakeFS
\par\textit{lakeFS adds Git-like branches, commits, merges, and hooks over object-storage data lakes. Tradeoff: Atomic data-lake changes enable testing, with server and metadata operations to manage. Watch for: Treating object versioning alone as branch-level data semantics is insufficient.}
\paragraph{MCQ-0829 [hard]} Which causal explanation of lakeFS connects what it does to when it is worth its cost?
\begin{enumerate}[label=\Alph*.]
\item Mechanism: Tests cover schema, nulls, ranges, uniqueness, distribution, freshness, and cross-table invariants. Consequence: Early detection prevents downstream damage, while brittle thresholds create alert fatigue.
\item Mechanism: Feature stores define reusable features and connect historical offline data with low-latency online serving. Consequence: Training-serving consistency improves but adds infrastructure and ownership requirements.
\item Mechanism: Runs capture code, parameters, data, metrics, artifacts, environment, and parent relationships. Consequence: Comparison and reproducibility improve but inconsistent naming creates a junk drawer.
\item Mechanism: lakeFS adds Git-like branches, commits, merges, and hooks over object-storage data lakes. Consequence: Atomic data-lake changes enable testing, with server and metadata operations to manage.
\end{enumerate}
\noindent\textbf{Answer.} Mechanism: lakeFS adds Git-like branches, commits, merges, and hooks over object-storage data lakes. Consequence: Atomic data-lake changes enable testing, with server and metadata operations to manage.
\par\textit{lakeFS adds Git-like branches, commits, merges, and hooks over object-storage data lakes. Tradeoff: Atomic data-lake changes enable testing, with server and metadata operations to manage. Watch for: Treating object versioning alone as branch-level data semantics is insufficient.}
\paragraph{FC-0830 [medium]} Define lakeFS in interview-ready language.
\noindent\textbf{Answer.} lakeFS adds Git-like branches, commits, merges, and hooks over object-storage data lakes.
\par\textit{lakeFS adds Git-like branches, commits, merges, and hooks over object-storage data lakes.}
\paragraph{FC-0831 [hard]} State the main tradeoff and production pitfall for lakeFS.
\noindent\textbf{Answer.} Tradeoff: Atomic data-lake changes enable testing, with server and metadata operations to manage. Pitfall: Treating object versioning alone as branch-level data semantics is insufficient.
\par\textit{Tradeoff: Atomic data-lake changes enable testing, with server and metadata operations to manage. Pitfall: Treating object versioning alone as branch-level data semantics is insufficient.}
\paragraph{LA-0832 [hard]} Explain lakeFS from first principles, then show how you would evaluate and operate it in production.
\noindent\textbf{Answer.} Start with the mechanism: lakeFS adds Git-like branches, commits, merges, and hooks over object-storage data lakes. Make the tradeoff explicit: Atomic data-lake changes enable testing, with server and metadata operations to manage. Define workload-specific offline metrics and a latency/cost budget, test important slices, instrument the critical path, and create a rollback criterion. The production review must explicitly guard against this failure: Treating object versioning alone as branch-level data semantics is insufficient.
\par\textit{A strong answer connects mechanism, measurable tradeoffs, failure isolation, observability, and rollback rather than listing product features.}
\paragraph{MCQ-0833 [medium]} A design requires this behavior: Open table formats add schemas, snapshots, partition evolution, and transactional metadata over object storage. Which mechanism should the team study?
\begin{enumerate}[label=\Alph*.]
\item Reproducible environments
\item Delta Lake and Iceberg
\item lakeFS
\item Experiment tracking
\end{enumerate}
\noindent\textbf{Answer.} Delta Lake and Iceberg
\par\textit{Open table formats add schemas, snapshots, partition evolution, and transactional metadata over object storage. Tradeoff: Reliable analytics improve, while compaction and catalog operations become essential. Watch for: Excess small files destroy scan performance despite correct table semantics.}
\paragraph{MCQ-0834 [hard]} The team proposes Delta Lake and Iceberg for a real workload. Which finding gives the correct decision boundary?
\begin{enumerate}[label=\Alph*.]
\item Atomic data-lake changes enable testing, with server and metadata operations to manage.
\item Reliable analytics improve, while compaction and catalog operations become essential.
\item Comparison and reproducibility improve but inconsistent naming creates a junk drawer.
\item Reproducibility improves at the expense of build maintenance and sometimes performance.
\end{enumerate}
\noindent\textbf{Answer.} Reliable analytics improve, while compaction and catalog operations become essential.
\par\textit{Open table formats add schemas, snapshots, partition evolution, and transactional metadata over object storage. Tradeoff: Reliable analytics improve, while compaction and catalog operations become essential. Watch for: Excess small files destroy scan performance despite correct table semantics.}
\paragraph{MCQ-0835 [hard]} Before releasing Delta Lake and Iceberg, which failure deserves a dedicated regression test?
\begin{enumerate}[label=\Alph*.]
\item Excess small files destroy scan performance despite correct table semantics.
\item Treating object versioning alone as branch-level data semantics is insufficient.
\item Logging a metric without dataset and code versions makes it uninterpretable.
\item A random seed alone cannot guarantee deterministic GPU execution.
\end{enumerate}
\noindent\textbf{Answer.} Excess small files destroy scan performance despite correct table semantics.
\par\textit{Open table formats add schemas, snapshots, partition evolution, and transactional metadata over object storage. Tradeoff: Reliable analytics improve, while compaction and catalog operations become essential. Watch for: Excess small files destroy scan performance despite correct table semantics.}
\paragraph{MCQ-0836 [medium]} A design review reveals this mistake: Excess small files destroy scan performance despite correct table semantics. Which mechanism should be traced first?
\begin{enumerate}[label=\Alph*.]
\item lakeFS
\item Experiment tracking
\item Reproducible environments
\item Delta Lake and Iceberg
\end{enumerate}
\noindent\textbf{Answer.} Delta Lake and Iceberg
\par\textit{Open table formats add schemas, snapshots, partition evolution, and transactional metadata over object storage. Tradeoff: Reliable analytics improve, while compaction and catalog operations become essential. Watch for: Excess small files destroy scan performance despite correct table semantics.}
\paragraph{MCQ-0837 [hard]} Which causal explanation of Delta Lake and Iceberg connects what it does to when it is worth its cost?
\begin{enumerate}[label=\Alph*.]
\item Mechanism: Lockfiles, containers, hardware metadata, seeds, and deterministic settings constrain run variability. Consequence: Reproducibility improves at the expense of build maintenance and sometimes performance.
\item Mechanism: lakeFS adds Git-like branches, commits, merges, and hooks over object-storage data lakes. Consequence: Atomic data-lake changes enable testing, with server and metadata operations to manage.
\item Mechanism: Runs capture code, parameters, data, metrics, artifacts, environment, and parent relationships. Consequence: Comparison and reproducibility improve but inconsistent naming creates a junk drawer.
\item Mechanism: Open table formats add schemas, snapshots, partition evolution, and transactional metadata over object storage. Consequence: Reliable analytics improve, while compaction and catalog operations become essential.
\end{enumerate}
\noindent\textbf{Answer.} Mechanism: Open table formats add schemas, snapshots, partition evolution, and transactional metadata over object storage. Consequence: Reliable analytics improve, while compaction and catalog operations become essential.
\par\textit{Open table formats add schemas, snapshots, partition evolution, and transactional metadata over object storage. Tradeoff: Reliable analytics improve, while compaction and catalog operations become essential. Watch for: Excess small files destroy scan performance despite correct table semantics.}
\paragraph{FC-0838 [medium]} Define Delta Lake and Iceberg in interview-ready language.
\noindent\textbf{Answer.} Open table formats add schemas, snapshots, partition evolution, and transactional metadata over object storage.
\par\textit{Open table formats add schemas, snapshots, partition evolution, and transactional metadata over object storage.}
\paragraph{FC-0839 [hard]} State the main tradeoff and production pitfall for Delta Lake and Iceberg.
\noindent\textbf{Answer.} Tradeoff: Reliable analytics improve, while compaction and catalog operations become essential. Pitfall: Excess small files destroy scan performance despite correct table semantics.
\par\textit{Tradeoff: Reliable analytics improve, while compaction and catalog operations become essential. Pitfall: Excess small files destroy scan performance despite correct table semantics.}
\paragraph{LA-0840 [hard]} Explain Delta Lake and Iceberg from first principles, then show how you would evaluate and operate it in production.
\noindent\textbf{Answer.} Start with the mechanism: Open table formats add schemas, snapshots, partition evolution, and transactional metadata over object storage. Make the tradeoff explicit: Reliable analytics improve, while compaction and catalog operations become essential. Define workload-specific offline metrics and a latency/cost budget, test important slices, instrument the critical path, and create a rollback criterion. The production review must explicitly guard against this failure: Excess small files destroy scan performance despite correct table semantics.
\par\textit{A strong answer connects mechanism, measurable tradeoffs, failure isolation, observability, and rollback rather than listing product features.}
\paragraph{MCQ-0841 [medium]} A design requires this behavior: Feature stores define reusable features and connect historical offline data with low-latency online serving. Which mechanism should the team study?
\begin{enumerate}[label=\Alph*.]
\item Experiment tracking
\item Data quality testing
\item Point-in-time correctness
\item Feature stores
\end{enumerate}
\noindent\textbf{Answer.} Feature stores
\par\textit{Feature stores define reusable features and connect historical offline data with low-latency online serving. Tradeoff: Training-serving consistency improves but adds infrastructure and ownership requirements. Watch for: Recomputing features differently online creates skew.}
\paragraph{MCQ-0842 [hard]} The team proposes Feature stores for a real workload. Which finding gives the correct decision boundary?
\begin{enumerate}[label=\Alph*.]
\item Comparison and reproducibility improve but inconsistent naming creates a junk drawer.
\item Training-serving consistency improves but adds infrastructure and ownership requirements.
\item Leakage is prevented at the cost of temporal keys and more expensive joins.
\item Early detection prevents downstream damage, while brittle thresholds create alert fatigue.
\end{enumerate}
\noindent\textbf{Answer.} Training-serving consistency improves but adds infrastructure and ownership requirements.
\par\textit{Feature stores define reusable features and connect historical offline data with low-latency online serving. Tradeoff: Training-serving consistency improves but adds infrastructure and ownership requirements. Watch for: Recomputing features differently online creates skew.}
\paragraph{MCQ-0843 [hard]} Before releasing Feature stores, which failure deserves a dedicated regression test?
\begin{enumerate}[label=\Alph*.]
\item Recomputing features differently online creates skew.
\item Checking only aggregate distributions misses failures in important slices.
\item Joining the latest feature value into old examples creates unrealistically strong training results.
\item Logging a metric without dataset and code versions makes it uninterpretable.
\end{enumerate}
\noindent\textbf{Answer.} Recomputing features differently online creates skew.
\par\textit{Feature stores define reusable features and connect historical offline data with low-latency online serving. Tradeoff: Training-serving consistency improves but adds infrastructure and ownership requirements. Watch for: Recomputing features differently online creates skew.}
\paragraph{MCQ-0844 [medium]} A design review reveals this mistake: Recomputing features differently online creates skew. Which mechanism should be traced first?
\begin{enumerate}[label=\Alph*.]
\item Experiment tracking
\item Point-in-time correctness
\item Data quality testing
\item Feature stores
\end{enumerate}
\noindent\textbf{Answer.} Feature stores
\par\textit{Feature stores define reusable features and connect historical offline data with low-latency online serving. Tradeoff: Training-serving consistency improves but adds infrastructure and ownership requirements. Watch for: Recomputing features differently online creates skew.}
\paragraph{MCQ-0845 [hard]} Which causal explanation of Feature stores connects what it does to when it is worth its cost?
\begin{enumerate}[label=\Alph*.]
\item Mechanism: Feature stores define reusable features and connect historical offline data with low-latency online serving. Consequence: Training-serving consistency improves but adds infrastructure and ownership requirements.
\item Mechanism: Tests cover schema, nulls, ranges, uniqueness, distribution, freshness, and cross-table invariants. Consequence: Early detection prevents downstream damage, while brittle thresholds create alert fatigue.
\item Mechanism: Historical joins select only feature values known at each label timestamp. Consequence: Leakage is prevented at the cost of temporal keys and more expensive joins.
\item Mechanism: Runs capture code, parameters, data, metrics, artifacts, environment, and parent relationships. Consequence: Comparison and reproducibility improve but inconsistent naming creates a junk drawer.
\end{enumerate}
\noindent\textbf{Answer.} Mechanism: Feature stores define reusable features and connect historical offline data with low-latency online serving. Consequence: Training-serving consistency improves but adds infrastructure and ownership requirements.
\par\textit{Feature stores define reusable features and connect historical offline data with low-latency online serving. Tradeoff: Training-serving consistency improves but adds infrastructure and ownership requirements. Watch for: Recomputing features differently online creates skew.}
\paragraph{FC-0846 [medium]} Define Feature stores in interview-ready language.
\noindent\textbf{Answer.} Feature stores define reusable features and connect historical offline data with low-latency online serving.
\par\textit{Feature stores define reusable features and connect historical offline data with low-latency online serving.}
\paragraph{FC-0847 [hard]} State the main tradeoff and production pitfall for Feature stores.
\noindent\textbf{Answer.} Tradeoff: Training-serving consistency improves but adds infrastructure and ownership requirements. Pitfall: Recomputing features differently online creates skew.
\par\textit{Tradeoff: Training-serving consistency improves but adds infrastructure and ownership requirements. Pitfall: Recomputing features differently online creates skew.}
\paragraph{LA-0848 [hard]} Explain Feature stores from first principles, then show how you would evaluate and operate it in production.
\noindent\textbf{Answer.} Start with the mechanism: Feature stores define reusable features and connect historical offline data with low-latency online serving. Make the tradeoff explicit: Training-serving consistency improves but adds infrastructure and ownership requirements. Define workload-specific offline metrics and a latency/cost budget, test important slices, instrument the critical path, and create a rollback criterion. The production review must explicitly guard against this failure: Recomputing features differently online creates skew.
\par\textit{A strong answer connects mechanism, measurable tradeoffs, failure isolation, observability, and rollback rather than listing product features.}
\paragraph{MCQ-0849 [medium]} A design requires this behavior: Historical joins select only feature values known at each label timestamp. Which mechanism should the team study?
\begin{enumerate}[label=\Alph*.]
\item Data lineage
\item Point-in-time correctness
\item lakeFS
\item Reproducible environments
\end{enumerate}
\noindent\textbf{Answer.} Point-in-time correctness
\par\textit{Historical joins select only feature values known at each label timestamp. Tradeoff: Leakage is prevented at the cost of temporal keys and more expensive joins. Watch for: Joining the latest feature value into old examples creates unrealistically strong training results.}
\paragraph{MCQ-0850 [hard]} The team proposes Point-in-time correctness for a real workload. Which finding gives the correct decision boundary?
\begin{enumerate}[label=\Alph*.]
\item Leakage is prevented at the cost of temporal keys and more expensive joins.
\item Impact analysis and auditability improve, while automated lineage can miss dynamic behavior.
\item Atomic data-lake changes enable testing, with server and metadata operations to manage.
\item Reproducibility improves at the expense of build maintenance and sometimes performance.
\end{enumerate}
\noindent\textbf{Answer.} Leakage is prevented at the cost of temporal keys and more expensive joins.
\par\textit{Historical joins select only feature values known at each label timestamp. Tradeoff: Leakage is prevented at the cost of temporal keys and more expensive joins. Watch for: Joining the latest feature value into old examples creates unrealistically strong training results.}
\paragraph{MCQ-0851 [hard]} Before releasing Point-in-time correctness, which failure deserves a dedicated regression test?
\begin{enumerate}[label=\Alph*.]
\item Treating object versioning alone as branch-level data semantics is insufficient.
\item Joining the latest feature value into old examples creates unrealistically strong training results.
\item Lineage without immutable identifiers gives a false sense of reproducibility.
\item A random seed alone cannot guarantee deterministic GPU execution.
\end{enumerate}
\noindent\textbf{Answer.} Joining the latest feature value into old examples creates unrealistically strong training results.
\par\textit{Historical joins select only feature values known at each label timestamp. Tradeoff: Leakage is prevented at the cost of temporal keys and more expensive joins. Watch for: Joining the latest feature value into old examples creates unrealistically strong training results.}
\paragraph{MCQ-0852 [medium]} A design review reveals this mistake: Joining the latest feature value into old examples creates unrealistically strong training results. Which mechanism should be traced first?
\begin{enumerate}[label=\Alph*.]
\item lakeFS
\item Point-in-time correctness
\item Data lineage
\item Reproducible environments
\end{enumerate}
\noindent\textbf{Answer.} Point-in-time correctness
\par\textit{Historical joins select only feature values known at each label timestamp. Tradeoff: Leakage is prevented at the cost of temporal keys and more expensive joins. Watch for: Joining the latest feature value into old examples creates unrealistically strong training results.}
\paragraph{MCQ-0853 [hard]} Which causal explanation of Point-in-time correctness connects what it does to when it is worth its cost?
\begin{enumerate}[label=\Alph*.]
\item Mechanism: Lineage links sources, transformations, datasets, features, prompts, models, evaluations, and deployments. Consequence: Impact analysis and auditability improve, while automated lineage can miss dynamic behavior.
\item Mechanism: lakeFS adds Git-like branches, commits, merges, and hooks over object-storage data lakes. Consequence: Atomic data-lake changes enable testing, with server and metadata operations to manage.
\item Mechanism: Historical joins select only feature values known at each label timestamp. Consequence: Leakage is prevented at the cost of temporal keys and more expensive joins.
\item Mechanism: Lockfiles, containers, hardware metadata, seeds, and deterministic settings constrain run variability. Consequence: Reproducibility improves at the expense of build maintenance and sometimes performance.
\end{enumerate}
\noindent\textbf{Answer.} Mechanism: Historical joins select only feature values known at each label timestamp. Consequence: Leakage is prevented at the cost of temporal keys and more expensive joins.
\par\textit{Historical joins select only feature values known at each label timestamp. Tradeoff: Leakage is prevented at the cost of temporal keys and more expensive joins. Watch for: Joining the latest feature value into old examples creates unrealistically strong training results.}
\paragraph{FC-0854 [medium]} Define Point-in-time correctness in interview-ready language.
\noindent\textbf{Answer.} Historical joins select only feature values known at each label timestamp.
\par\textit{Historical joins select only feature values known at each label timestamp.}
\paragraph{FC-0855 [hard]} State the main tradeoff and production pitfall for Point-in-time correctness.
\noindent\textbf{Answer.} Tradeoff: Leakage is prevented at the cost of temporal keys and more expensive joins. Pitfall: Joining the latest feature value into old examples creates unrealistically strong training results.
\par\textit{Tradeoff: Leakage is prevented at the cost of temporal keys and more expensive joins. Pitfall: Joining the latest feature value into old examples creates unrealistically strong training results.}
\paragraph{LA-0856 [hard]} Explain Point-in-time correctness from first principles, then show how you would evaluate and operate it in production.
\noindent\textbf{Answer.} Start with the mechanism: Historical joins select only feature values known at each label timestamp. Make the tradeoff explicit: Leakage is prevented at the cost of temporal keys and more expensive joins. Define workload-specific offline metrics and a latency/cost budget, test important slices, instrument the critical path, and create a rollback criterion. The production review must explicitly guard against this failure: Joining the latest feature value into old examples creates unrealistically strong training results.
\par\textit{A strong answer connects mechanism, measurable tradeoffs, failure isolation, observability, and rollback rather than listing product features.}
\paragraph{MCQ-0857 [medium]} A design requires this behavior: Runs capture code, parameters, data, metrics, artifacts, environment, and parent relationships. Which mechanism should the team study?
\begin{enumerate}[label=\Alph*.]
\item Feature stores
\item Point-in-time correctness
\item Experiment tracking
\item lakeFS
\end{enumerate}
\noindent\textbf{Answer.} Experiment tracking
\par\textit{Runs capture code, parameters, data, metrics, artifacts, environment, and parent relationships. Tradeoff: Comparison and reproducibility improve but inconsistent naming creates a junk drawer. Watch for: Logging a metric without dataset and code versions makes it uninterpretable.}
\paragraph{MCQ-0858 [hard]} The team proposes Experiment tracking for a real workload. Which finding gives the correct decision boundary?
\begin{enumerate}[label=\Alph*.]
\item Training-serving consistency improves but adds infrastructure and ownership requirements.
\item Leakage is prevented at the cost of temporal keys and more expensive joins.
\item Atomic data-lake changes enable testing, with server and metadata operations to manage.
\item Comparison and reproducibility improve but inconsistent naming creates a junk drawer.
\end{enumerate}
\noindent\textbf{Answer.} Comparison and reproducibility improve but inconsistent naming creates a junk drawer.
\par\textit{Runs capture code, parameters, data, metrics, artifacts, environment, and parent relationships. Tradeoff: Comparison and reproducibility improve but inconsistent naming creates a junk drawer. Watch for: Logging a metric without dataset and code versions makes it uninterpretable.}
\paragraph{MCQ-0859 [hard]} Before releasing Experiment tracking, which failure deserves a dedicated regression test?
\begin{enumerate}[label=\Alph*.]
\item Treating object versioning alone as branch-level data semantics is insufficient.
\item Joining the latest feature value into old examples creates unrealistically strong training results.
\item Logging a metric without dataset and code versions makes it uninterpretable.
\item Recomputing features differently online creates skew.
\end{enumerate}
\noindent\textbf{Answer.} Logging a metric without dataset and code versions makes it uninterpretable.
\par\textit{Runs capture code, parameters, data, metrics, artifacts, environment, and parent relationships. Tradeoff: Comparison and reproducibility improve but inconsistent naming creates a junk drawer. Watch for: Logging a metric without dataset and code versions makes it uninterpretable.}
\paragraph{MCQ-0860 [medium]} A design review reveals this mistake: Logging a metric without dataset and code versions makes it uninterpretable. Which mechanism should be traced first?
\begin{enumerate}[label=\Alph*.]
\item Experiment tracking
\item Feature stores
\item lakeFS
\item Point-in-time correctness
\end{enumerate}
\noindent\textbf{Answer.} Experiment tracking
\par\textit{Runs capture code, parameters, data, metrics, artifacts, environment, and parent relationships. Tradeoff: Comparison and reproducibility improve but inconsistent naming creates a junk drawer. Watch for: Logging a metric without dataset and code versions makes it uninterpretable.}
\paragraph{MCQ-0861 [hard]} Which causal explanation of Experiment tracking connects what it does to when it is worth its cost?
\begin{enumerate}[label=\Alph*.]
\item Mechanism: Historical joins select only feature values known at each label timestamp. Consequence: Leakage is prevented at the cost of temporal keys and more expensive joins.
\item Mechanism: lakeFS adds Git-like branches, commits, merges, and hooks over object-storage data lakes. Consequence: Atomic data-lake changes enable testing, with server and metadata operations to manage.
\item Mechanism: Runs capture code, parameters, data, metrics, artifacts, environment, and parent relationships. Consequence: Comparison and reproducibility improve but inconsistent naming creates a junk drawer.
\item Mechanism: Feature stores define reusable features and connect historical offline data with low-latency online serving. Consequence: Training-serving consistency improves but adds infrastructure and ownership requirements.
\end{enumerate}
\noindent\textbf{Answer.} Mechanism: Runs capture code, parameters, data, metrics, artifacts, environment, and parent relationships. Consequence: Comparison and reproducibility improve but inconsistent naming creates a junk drawer.
\par\textit{Runs capture code, parameters, data, metrics, artifacts, environment, and parent relationships. Tradeoff: Comparison and reproducibility improve but inconsistent naming creates a junk drawer. Watch for: Logging a metric without dataset and code versions makes it uninterpretable.}
\paragraph{FC-0862 [medium]} Define Experiment tracking in interview-ready language.
\noindent\textbf{Answer.} Runs capture code, parameters, data, metrics, artifacts, environment, and parent relationships.
\par\textit{Runs capture code, parameters, data, metrics, artifacts, environment, and parent relationships.}
\paragraph{FC-0863 [hard]} State the main tradeoff and production pitfall for Experiment tracking.
\noindent\textbf{Answer.} Tradeoff: Comparison and reproducibility improve but inconsistent naming creates a junk drawer. Pitfall: Logging a metric without dataset and code versions makes it uninterpretable.
\par\textit{Tradeoff: Comparison and reproducibility improve but inconsistent naming creates a junk drawer. Pitfall: Logging a metric without dataset and code versions makes it uninterpretable.}
\paragraph{LA-0864 [hard]} Explain Experiment tracking from first principles, then show how you would evaluate and operate it in production.
\noindent\textbf{Answer.} Start with the mechanism: Runs capture code, parameters, data, metrics, artifacts, environment, and parent relationships. Make the tradeoff explicit: Comparison and reproducibility improve but inconsistent naming creates a junk drawer. Define workload-specific offline metrics and a latency/cost budget, test important slices, instrument the critical path, and create a rollback criterion. The production review must explicitly guard against this failure: Logging a metric without dataset and code versions makes it uninterpretable.
\par\textit{A strong answer connects mechanism, measurable tradeoffs, failure isolation, observability, and rollback rather than listing product features.}
\paragraph{MCQ-0865 [medium]} A design requires this behavior: Contracts specify schema, semantics, freshness, ownership, quality, and compatibility expectations between producers and consumers. Which mechanism should the team study?
\begin{enumerate}[label=\Alph*.]
\item Feature stores
\item Data contracts
\item Reproducible environments
\item Data lineage
\end{enumerate}
\noindent\textbf{Answer.} Data contracts
\par\textit{Contracts specify schema, semantics, freshness, ownership, quality, and compatibility expectations between producers and consumers. Tradeoff: Clear boundaries reduce silent breakage but require governance and enforcement. Watch for: Validating shape without semantic definitions misses unit and meaning changes.}
\paragraph{MCQ-0866 [hard]} The team proposes Data contracts for a real workload. Which finding gives the correct decision boundary?
\begin{enumerate}[label=\Alph*.]
\item Training-serving consistency improves but adds infrastructure and ownership requirements.
\item Impact analysis and auditability improve, while automated lineage can miss dynamic behavior.
\item Reproducibility improves at the expense of build maintenance and sometimes performance.
\item Clear boundaries reduce silent breakage but require governance and enforcement.
\end{enumerate}
\noindent\textbf{Answer.} Clear boundaries reduce silent breakage but require governance and enforcement.
\par\textit{Contracts specify schema, semantics, freshness, ownership, quality, and compatibility expectations between producers and consumers. Tradeoff: Clear boundaries reduce silent breakage but require governance and enforcement. Watch for: Validating shape without semantic definitions misses unit and meaning changes.}
\paragraph{MCQ-0867 [hard]} Before releasing Data contracts, which failure deserves a dedicated regression test?
\begin{enumerate}[label=\Alph*.]
\item Lineage without immutable identifiers gives a false sense of reproducibility.
\item Recomputing features differently online creates skew.
\item Validating shape without semantic definitions misses unit and meaning changes.
\item A random seed alone cannot guarantee deterministic GPU execution.
\end{enumerate}
\noindent\textbf{Answer.} Validating shape without semantic definitions misses unit and meaning changes.
\par\textit{Contracts specify schema, semantics, freshness, ownership, quality, and compatibility expectations between producers and consumers. Tradeoff: Clear boundaries reduce silent breakage but require governance and enforcement. Watch for: Validating shape without semantic definitions misses unit and meaning changes.}
\paragraph{MCQ-0868 [medium]} A design review reveals this mistake: Validating shape without semantic definitions misses unit and meaning changes. Which mechanism should be traced first?
\begin{enumerate}[label=\Alph*.]
\item Data lineage
\item Feature stores
\item Reproducible environments
\item Data contracts
\end{enumerate}
\noindent\textbf{Answer.} Data contracts
\par\textit{Contracts specify schema, semantics, freshness, ownership, quality, and compatibility expectations between producers and consumers. Tradeoff: Clear boundaries reduce silent breakage but require governance and enforcement. Watch for: Validating shape without semantic definitions misses unit and meaning changes.}
\paragraph{MCQ-0869 [hard]} Which causal explanation of Data contracts connects what it does to when it is worth its cost?
\begin{enumerate}[label=\Alph*.]
\item Mechanism: Contracts specify schema, semantics, freshness, ownership, quality, and compatibility expectations between producers and consumers. Consequence: Clear boundaries reduce silent breakage but require governance and enforcement.
\item Mechanism: Feature stores define reusable features and connect historical offline data with low-latency online serving. Consequence: Training-serving consistency improves but adds infrastructure and ownership requirements.
\item Mechanism: Lineage links sources, transformations, datasets, features, prompts, models, evaluations, and deployments. Consequence: Impact analysis and auditability improve, while automated lineage can miss dynamic behavior.
\item Mechanism: Lockfiles, containers, hardware metadata, seeds, and deterministic settings constrain run variability. Consequence: Reproducibility improves at the expense of build maintenance and sometimes performance.
\end{enumerate}
\noindent\textbf{Answer.} Mechanism: Contracts specify schema, semantics, freshness, ownership, quality, and compatibility expectations between producers and consumers. Consequence: Clear boundaries reduce silent breakage but require governance and enforcement.
\par\textit{Contracts specify schema, semantics, freshness, ownership, quality, and compatibility expectations between producers and consumers. Tradeoff: Clear boundaries reduce silent breakage but require governance and enforcement. Watch for: Validating shape without semantic definitions misses unit and meaning changes.}
\paragraph{FC-0870 [medium]} Define Data contracts in interview-ready language.
\noindent\textbf{Answer.} Contracts specify schema, semantics, freshness, ownership, quality, and compatibility expectations between producers and consumers.
\par\textit{Contracts specify schema, semantics, freshness, ownership, quality, and compatibility expectations between producers and consumers.}
\paragraph{FC-0871 [hard]} State the main tradeoff and production pitfall for Data contracts.
\noindent\textbf{Answer.} Tradeoff: Clear boundaries reduce silent breakage but require governance and enforcement. Pitfall: Validating shape without semantic definitions misses unit and meaning changes.
\par\textit{Tradeoff: Clear boundaries reduce silent breakage but require governance and enforcement. Pitfall: Validating shape without semantic definitions misses unit and meaning changes.}
\paragraph{LA-0872 [hard]} Explain Data contracts from first principles, then show how you would evaluate and operate it in production.
\noindent\textbf{Answer.} Start with the mechanism: Contracts specify schema, semantics, freshness, ownership, quality, and compatibility expectations between producers and consumers. Make the tradeoff explicit: Clear boundaries reduce silent breakage but require governance and enforcement. Define workload-specific offline metrics and a latency/cost budget, test important slices, instrument the critical path, and create a rollback criterion. The production review must explicitly guard against this failure: Validating shape without semantic definitions misses unit and meaning changes.
\par\textit{A strong answer connects mechanism, measurable tradeoffs, failure isolation, observability, and rollback rather than listing product features.}
\paragraph{MCQ-0873 [medium]} A design requires this behavior: Tests cover schema, nulls, ranges, uniqueness, distribution, freshness, and cross-table invariants. Which mechanism should the team study?
\begin{enumerate}[label=\Alph*.]
\item DVC
\item Data contracts
\item Point-in-time correctness
\item Data quality testing
\end{enumerate}
\noindent\textbf{Answer.} Data quality testing
\par\textit{Tests cover schema, nulls, ranges, uniqueness, distribution, freshness, and cross-table invariants. Tradeoff: Early detection prevents downstream damage, while brittle thresholds create alert fatigue. Watch for: Checking only aggregate distributions misses failures in important slices.}
\paragraph{MCQ-0874 [hard]} The team proposes Data quality testing for a real workload. Which finding gives the correct decision boundary?
\begin{enumerate}[label=\Alph*.]
\item Early detection prevents downstream damage, while brittle thresholds create alert fatigue.
\item Leakage is prevented at the cost of temporal keys and more expensive joins.
\item Clear boundaries reduce silent breakage but require governance and enforcement.
\item Familiar Git workflows improve reproducibility, while remote storage and cache discipline are required.
\end{enumerate}
\noindent\textbf{Answer.} Early detection prevents downstream damage, while brittle thresholds create alert fatigue.
\par\textit{Tests cover schema, nulls, ranges, uniqueness, distribution, freshness, and cross-table invariants. Tradeoff: Early detection prevents downstream damage, while brittle thresholds create alert fatigue. Watch for: Checking only aggregate distributions misses failures in important slices.}
\paragraph{MCQ-0875 [hard]} Before releasing Data quality testing, which failure deserves a dedicated regression test?
\begin{enumerate}[label=\Alph*.]
\item Committing only code while silently replacing remote data breaks lineage.
\item Checking only aggregate distributions misses failures in important slices.
\item Joining the latest feature value into old examples creates unrealistically strong training results.
\item Validating shape without semantic definitions misses unit and meaning changes.
\end{enumerate}
\noindent\textbf{Answer.} Checking only aggregate distributions misses failures in important slices.
\par\textit{Tests cover schema, nulls, ranges, uniqueness, distribution, freshness, and cross-table invariants. Tradeoff: Early detection prevents downstream damage, while brittle thresholds create alert fatigue. Watch for: Checking only aggregate distributions misses failures in important slices.}
\paragraph{MCQ-0876 [medium]} A design review reveals this mistake: Checking only aggregate distributions misses failures in important slices. Which mechanism should be traced first?
\begin{enumerate}[label=\Alph*.]
\item Data contracts
\item DVC
\item Data quality testing
\item Point-in-time correctness
\end{enumerate}
\noindent\textbf{Answer.} Data quality testing
\par\textit{Tests cover schema, nulls, ranges, uniqueness, distribution, freshness, and cross-table invariants. Tradeoff: Early detection prevents downstream damage, while brittle thresholds create alert fatigue. Watch for: Checking only aggregate distributions misses failures in important slices.}
\paragraph{MCQ-0877 [hard]} Which causal explanation of Data quality testing connects what it does to when it is worth its cost?
\begin{enumerate}[label=\Alph*.]
\item Mechanism: Tests cover schema, nulls, ranges, uniqueness, distribution, freshness, and cross-table invariants. Consequence: Early detection prevents downstream damage, while brittle thresholds create alert fatigue.
\item Mechanism: Contracts specify schema, semantics, freshness, ownership, quality, and compatibility expectations between producers and consumers. Consequence: Clear boundaries reduce silent breakage but require governance and enforcement.
\item Mechanism: Historical joins select only feature values known at each label timestamp. Consequence: Leakage is prevented at the cost of temporal keys and more expensive joins.
\item Mechanism: DVC versions large data and model pointers with Git and defines reproducible pipeline stages and experiments. Consequence: Familiar Git workflows improve reproducibility, while remote storage and cache discipline are required.
\end{enumerate}
\noindent\textbf{Answer.} Mechanism: Tests cover schema, nulls, ranges, uniqueness, distribution, freshness, and cross-table invariants. Consequence: Early detection prevents downstream damage, while brittle thresholds create alert fatigue.
\par\textit{Tests cover schema, nulls, ranges, uniqueness, distribution, freshness, and cross-table invariants. Tradeoff: Early detection prevents downstream damage, while brittle thresholds create alert fatigue. Watch for: Checking only aggregate distributions misses failures in important slices.}
\paragraph{FC-0878 [medium]} Define Data quality testing in interview-ready language.
\noindent\textbf{Answer.} Tests cover schema, nulls, ranges, uniqueness, distribution, freshness, and cross-table invariants.
\par\textit{Tests cover schema, nulls, ranges, uniqueness, distribution, freshness, and cross-table invariants.}
\paragraph{FC-0879 [hard]} State the main tradeoff and production pitfall for Data quality testing.
\noindent\textbf{Answer.} Tradeoff: Early detection prevents downstream damage, while brittle thresholds create alert fatigue. Pitfall: Checking only aggregate distributions misses failures in important slices.
\par\textit{Tradeoff: Early detection prevents downstream damage, while brittle thresholds create alert fatigue. Pitfall: Checking only aggregate distributions misses failures in important slices.}
\paragraph{LA-0880 [hard]} Explain Data quality testing from first principles, then show how you would evaluate and operate it in production.
\noindent\textbf{Answer.} Start with the mechanism: Tests cover schema, nulls, ranges, uniqueness, distribution, freshness, and cross-table invariants. Make the tradeoff explicit: Early detection prevents downstream damage, while brittle thresholds create alert fatigue. Define workload-specific offline metrics and a latency/cost budget, test important slices, instrument the critical path, and create a rollback criterion. The production review must explicitly guard against this failure: Checking only aggregate distributions misses failures in important slices.
\par\textit{A strong answer connects mechanism, measurable tradeoffs, failure isolation, observability, and rollback rather than listing product features.}
\paragraph{MCQ-0881 [medium]} A design requires this behavior: Lineage links sources, transformations, datasets, features, prompts, models, evaluations, and deployments. Which mechanism should the team study?
\begin{enumerate}[label=\Alph*.]
\item Data lineage
\item lakeFS
\item Data contracts
\item Point-in-time correctness
\end{enumerate}
\noindent\textbf{Answer.} Data lineage
\par\textit{Lineage links sources, transformations, datasets, features, prompts, models, evaluations, and deployments. Tradeoff: Impact analysis and auditability improve, while automated lineage can miss dynamic behavior. Watch for: Lineage without immutable identifiers gives a false sense of reproducibility.}
\paragraph{MCQ-0882 [hard]} The team proposes Data lineage for a real workload. Which finding gives the correct decision boundary?
\begin{enumerate}[label=\Alph*.]
\item Atomic data-lake changes enable testing, with server and metadata operations to manage.
\item Clear boundaries reduce silent breakage but require governance and enforcement.
\item Leakage is prevented at the cost of temporal keys and more expensive joins.
\item Impact analysis and auditability improve, while automated lineage can miss dynamic behavior.
\end{enumerate}
\noindent\textbf{Answer.} Impact analysis and auditability improve, while automated lineage can miss dynamic behavior.
\par\textit{Lineage links sources, transformations, datasets, features, prompts, models, evaluations, and deployments. Tradeoff: Impact analysis and auditability improve, while automated lineage can miss dynamic behavior. Watch for: Lineage without immutable identifiers gives a false sense of reproducibility.}
\paragraph{MCQ-0883 [hard]} Before releasing Data lineage, which failure deserves a dedicated regression test?
\begin{enumerate}[label=\Alph*.]
\item Lineage without immutable identifiers gives a false sense of reproducibility.
\item Validating shape without semantic definitions misses unit and meaning changes.
\item Treating object versioning alone as branch-level data semantics is insufficient.
\item Joining the latest feature value into old examples creates unrealistically strong training results.
\end{enumerate}
\noindent\textbf{Answer.} Lineage without immutable identifiers gives a false sense of reproducibility.
\par\textit{Lineage links sources, transformations, datasets, features, prompts, models, evaluations, and deployments. Tradeoff: Impact analysis and auditability improve, while automated lineage can miss dynamic behavior. Watch for: Lineage without immutable identifiers gives a false sense of reproducibility.}
\paragraph{MCQ-0884 [medium]} A design review reveals this mistake: Lineage without immutable identifiers gives a false sense of reproducibility. Which mechanism should be traced first?
\begin{enumerate}[label=\Alph*.]
\item Point-in-time correctness
\item lakeFS
\item Data lineage
\item Data contracts
\end{enumerate}
\noindent\textbf{Answer.} Data lineage
\par\textit{Lineage links sources, transformations, datasets, features, prompts, models, evaluations, and deployments. Tradeoff: Impact analysis and auditability improve, while automated lineage can miss dynamic behavior. Watch for: Lineage without immutable identifiers gives a false sense of reproducibility.}
\paragraph{MCQ-0885 [hard]} Which causal explanation of Data lineage connects what it does to when it is worth its cost?
\begin{enumerate}[label=\Alph*.]
\item Mechanism: Historical joins select only feature values known at each label timestamp. Consequence: Leakage is prevented at the cost of temporal keys and more expensive joins.
\item Mechanism: Lineage links sources, transformations, datasets, features, prompts, models, evaluations, and deployments. Consequence: Impact analysis and auditability improve, while automated lineage can miss dynamic behavior.
\item Mechanism: Contracts specify schema, semantics, freshness, ownership, quality, and compatibility expectations between producers and consumers. Consequence: Clear boundaries reduce silent breakage but require governance and enforcement.
\item Mechanism: lakeFS adds Git-like branches, commits, merges, and hooks over object-storage data lakes. Consequence: Atomic data-lake changes enable testing, with server and metadata operations to manage.
\end{enumerate}
\noindent\textbf{Answer.} Mechanism: Lineage links sources, transformations, datasets, features, prompts, models, evaluations, and deployments. Consequence: Impact analysis and auditability improve, while automated lineage can miss dynamic behavior.
\par\textit{Lineage links sources, transformations, datasets, features, prompts, models, evaluations, and deployments. Tradeoff: Impact analysis and auditability improve, while automated lineage can miss dynamic behavior. Watch for: Lineage without immutable identifiers gives a false sense of reproducibility.}
\paragraph{FC-0886 [medium]} Define Data lineage in interview-ready language.
\noindent\textbf{Answer.} Lineage links sources, transformations, datasets, features, prompts, models, evaluations, and deployments.
\par\textit{Lineage links sources, transformations, datasets, features, prompts, models, evaluations, and deployments.}
\paragraph{FC-0887 [hard]} State the main tradeoff and production pitfall for Data lineage.
\noindent\textbf{Answer.} Tradeoff: Impact analysis and auditability improve, while automated lineage can miss dynamic behavior. Pitfall: Lineage without immutable identifiers gives a false sense of reproducibility.
\par\textit{Tradeoff: Impact analysis and auditability improve, while automated lineage can miss dynamic behavior. Pitfall: Lineage without immutable identifiers gives a false sense of reproducibility.}
\paragraph{LA-0888 [hard]} Explain Data lineage from first principles, then show how you would evaluate and operate it in production.
\noindent\textbf{Answer.} Start with the mechanism: Lineage links sources, transformations, datasets, features, prompts, models, evaluations, and deployments. Make the tradeoff explicit: Impact analysis and auditability improve, while automated lineage can miss dynamic behavior. Define workload-specific offline metrics and a latency/cost budget, test important slices, instrument the critical path, and create a rollback criterion. The production review must explicitly guard against this failure: Lineage without immutable identifiers gives a false sense of reproducibility.
\par\textit{A strong answer connects mechanism, measurable tradeoffs, failure isolation, observability, and rollback rather than listing product features.}
\paragraph{MCQ-0889 [medium]} A design requires this behavior: Lockfiles, containers, hardware metadata, seeds, and deterministic settings constrain run variability. Which mechanism should the team study?
\begin{enumerate}[label=\Alph*.]
\item Reproducible environments
\item lakeFS
\item DVC
\item Data lineage
\end{enumerate}
\noindent\textbf{Answer.} Reproducible environments
\par\textit{Lockfiles, containers, hardware metadata, seeds, and deterministic settings constrain run variability. Tradeoff: Reproducibility improves at the expense of build maintenance and sometimes performance. Watch for: A random seed alone cannot guarantee deterministic GPU execution.}
\paragraph{MCQ-0890 [hard]} The team proposes Reproducible environments for a real workload. Which finding gives the correct decision boundary?
\begin{enumerate}[label=\Alph*.]
\item Atomic data-lake changes enable testing, with server and metadata operations to manage.
\item Reproducibility improves at the expense of build maintenance and sometimes performance.
\item Impact analysis and auditability improve, while automated lineage can miss dynamic behavior.
\item Familiar Git workflows improve reproducibility, while remote storage and cache discipline are required.
\end{enumerate}
\noindent\textbf{Answer.} Reproducibility improves at the expense of build maintenance and sometimes performance.
\par\textit{Lockfiles, containers, hardware metadata, seeds, and deterministic settings constrain run variability. Tradeoff: Reproducibility improves at the expense of build maintenance and sometimes performance. Watch for: A random seed alone cannot guarantee deterministic GPU execution.}
\paragraph{MCQ-0891 [hard]} Before releasing Reproducible environments, which failure deserves a dedicated regression test?
\begin{enumerate}[label=\Alph*.]
\item A random seed alone cannot guarantee deterministic GPU execution.
\item Committing only code while silently replacing remote data breaks lineage.
\item Treating object versioning alone as branch-level data semantics is insufficient.
\item Lineage without immutable identifiers gives a false sense of reproducibility.
\end{enumerate}
\noindent\textbf{Answer.} A random seed alone cannot guarantee deterministic GPU execution.
\par\textit{Lockfiles, containers, hardware metadata, seeds, and deterministic settings constrain run variability. Tradeoff: Reproducibility improves at the expense of build maintenance and sometimes performance. Watch for: A random seed alone cannot guarantee deterministic GPU execution.}
\paragraph{MCQ-0892 [medium]} A design review reveals this mistake: A random seed alone cannot guarantee deterministic GPU execution. Which mechanism should be traced first?
\begin{enumerate}[label=\Alph*.]
\item Reproducible environments
\item Data lineage
\item DVC
\item lakeFS
\end{enumerate}
\noindent\textbf{Answer.} Reproducible environments
\par\textit{Lockfiles, containers, hardware metadata, seeds, and deterministic settings constrain run variability. Tradeoff: Reproducibility improves at the expense of build maintenance and sometimes performance. Watch for: A random seed alone cannot guarantee deterministic GPU execution.}
\paragraph{MCQ-0893 [hard]} Which causal explanation of Reproducible environments connects what it does to when it is worth its cost?
\begin{enumerate}[label=\Alph*.]
\item Mechanism: Lockfiles, containers, hardware metadata, seeds, and deterministic settings constrain run variability. Consequence: Reproducibility improves at the expense of build maintenance and sometimes performance.
\item Mechanism: lakeFS adds Git-like branches, commits, merges, and hooks over object-storage data lakes. Consequence: Atomic data-lake changes enable testing, with server and metadata operations to manage.
\item Mechanism: DVC versions large data and model pointers with Git and defines reproducible pipeline stages and experiments. Consequence: Familiar Git workflows improve reproducibility, while remote storage and cache discipline are required.
\item Mechanism: Lineage links sources, transformations, datasets, features, prompts, models, evaluations, and deployments. Consequence: Impact analysis and auditability improve, while automated lineage can miss dynamic behavior.
\end{enumerate}
\noindent\textbf{Answer.} Mechanism: Lockfiles, containers, hardware metadata, seeds, and deterministic settings constrain run variability. Consequence: Reproducibility improves at the expense of build maintenance and sometimes performance.
\par\textit{Lockfiles, containers, hardware metadata, seeds, and deterministic settings constrain run variability. Tradeoff: Reproducibility improves at the expense of build maintenance and sometimes performance. Watch for: A random seed alone cannot guarantee deterministic GPU execution.}
\paragraph{FC-0894 [medium]} Define Reproducible environments in interview-ready language.
\noindent\textbf{Answer.} Lockfiles, containers, hardware metadata, seeds, and deterministic settings constrain run variability.
\par\textit{Lockfiles, containers, hardware metadata, seeds, and deterministic settings constrain run variability.}
\paragraph{FC-0895 [hard]} State the main tradeoff and production pitfall for Reproducible environments.
\noindent\textbf{Answer.} Tradeoff: Reproducibility improves at the expense of build maintenance and sometimes performance. Pitfall: A random seed alone cannot guarantee deterministic GPU execution.
\par\textit{Tradeoff: Reproducibility improves at the expense of build maintenance and sometimes performance. Pitfall: A random seed alone cannot guarantee deterministic GPU execution.}
\paragraph{LA-0896 [hard]} Explain Reproducible environments from first principles, then show how you would evaluate and operate it in production.
\noindent\textbf{Answer.} Start with the mechanism: Lockfiles, containers, hardware metadata, seeds, and deterministic settings constrain run variability. Make the tradeoff explicit: Reproducibility improves at the expense of build maintenance and sometimes performance. Define workload-specific offline metrics and a latency/cost budget, test important slices, instrument the critical path, and create a rollback criterion. The production review must explicitly guard against this failure: A random seed alone cannot guarantee deterministic GPU execution.
\par\textit{A strong answer connects mechanism, measurable tradeoffs, failure isolation, observability, and rollback rather than listing product features.}
\subsection{Transformer Theory and Training}
\paragraph{MCQ-0897 [medium]} A design requires this behavior: Tokenizers map text or bytes into model vocabulary IDs using algorithms such as BPE, WordPiece, or unigram models. Which mechanism should the team study?
\begin{enumerate}[label=\Alph*.]
\item DPO
\item QLoRA
\item Causal language modeling
\item Tokenization
\end{enumerate}
\noindent\textbf{Answer.} Tokenization
\par\textit{Tokenizers map text or bytes into model vocabulary IDs using algorithms such as BPE, WordPiece, or unigram models. Tradeoff: Larger vocabularies shorten sequences but increase embedding parameters and sparse coverage. Watch for: Comparing context lengths across tokenizers as if tokens were equal misstates capacity and cost.}
\paragraph{MCQ-0898 [hard]} The team proposes Tokenization for a real workload. Which finding gives the correct decision boundary?
\begin{enumerate}[label=\Alph*.]
\item It simplifies preference training but still depends on pair quality and distribution coverage.
\item Fine-tuning large models becomes cheaper, while kernels and quantization settings affect stability.
\item It provides a simple scalable objective but does not directly optimize instruction usefulness or truth.
\item Larger vocabularies shorten sequences but increase embedding parameters and sparse coverage.
\end{enumerate}
\noindent\textbf{Answer.} Larger vocabularies shorten sequences but increase embedding parameters and sparse coverage.
\par\textit{Tokenizers map text or bytes into model vocabulary IDs using algorithms such as BPE, WordPiece, or unigram models. Tradeoff: Larger vocabularies shorten sequences but increase embedding parameters and sparse coverage. Watch for: Comparing context lengths across tokenizers as if tokens were equal misstates capacity and cost.}
\paragraph{MCQ-0899 [hard]} Before releasing Tokenization, which failure deserves a dedicated regression test?
\begin{enumerate}[label=\Alph*.]
\item Treating implicit user choices as clean preferences introduces confounding.
\item Interpreting low perplexity as safe task performance is invalid.
\item Confusing training quantization with final serving quantization produces wrong expectations.
\item Comparing context lengths across tokenizers as if tokens were equal misstates capacity and cost.
\end{enumerate}
\noindent\textbf{Answer.} Comparing context lengths across tokenizers as if tokens were equal misstates capacity and cost.
\par\textit{Tokenizers map text or bytes into model vocabulary IDs using algorithms such as BPE, WordPiece, or unigram models. Tradeoff: Larger vocabularies shorten sequences but increase embedding parameters and sparse coverage. Watch for: Comparing context lengths across tokenizers as if tokens were equal misstates capacity and cost.}
\paragraph{MCQ-0900 [medium]} A design review reveals this mistake: Comparing context lengths across tokenizers as if tokens were equal misstates capacity and cost. Which mechanism should be traced first?
\begin{enumerate}[label=\Alph*.]
\item DPO
\item Tokenization
\item Causal language modeling
\item QLoRA
\end{enumerate}
\noindent\textbf{Answer.} Tokenization
\par\textit{Tokenizers map text or bytes into model vocabulary IDs using algorithms such as BPE, WordPiece, or unigram models. Tradeoff: Larger vocabularies shorten sequences but increase embedding parameters and sparse coverage. Watch for: Comparing context lengths across tokenizers as if tokens were equal misstates capacity and cost.}
\paragraph{MCQ-0901 [hard]} Which causal explanation of Tokenization connects what it does to when it is worth its cost?
\begin{enumerate}[label=\Alph*.]
\item Mechanism: Next-token training factorizes sequence likelihood from left to right under a causal mask. Consequence: It provides a simple scalable objective but does not directly optimize instruction usefulness or truth.
\item Mechanism: Direct Preference Optimization fits preferred over rejected responses through a classification-style objective relative to a reference policy. Consequence: It simplifies preference training but still depends on pair quality and distribution coverage.
\item Mechanism: QLoRA backpropagates through a frozen quantized base model into LoRA adapters using memory-saving quantization techniques. Consequence: Fine-tuning large models becomes cheaper, while kernels and quantization settings affect stability.
\item Mechanism: Tokenizers map text or bytes into model vocabulary IDs using algorithms such as BPE, WordPiece, or unigram models. Consequence: Larger vocabularies shorten sequences but increase embedding parameters and sparse coverage.
\end{enumerate}
\noindent\textbf{Answer.} Mechanism: Tokenizers map text or bytes into model vocabulary IDs using algorithms such as BPE, WordPiece, or unigram models. Consequence: Larger vocabularies shorten sequences but increase embedding parameters and sparse coverage.
\par\textit{Tokenizers map text or bytes into model vocabulary IDs using algorithms such as BPE, WordPiece, or unigram models. Tradeoff: Larger vocabularies shorten sequences but increase embedding parameters and sparse coverage. Watch for: Comparing context lengths across tokenizers as if tokens were equal misstates capacity and cost.}
\paragraph{FC-0902 [medium]} Define Tokenization in interview-ready language.
\noindent\textbf{Answer.} Tokenizers map text or bytes into model vocabulary IDs using algorithms such as BPE, WordPiece, or unigram models.
\par\textit{Tokenizers map text or bytes into model vocabulary IDs using algorithms such as BPE, WordPiece, or unigram models.}
\paragraph{FC-0903 [hard]} State the main tradeoff and production pitfall for Tokenization.
\noindent\textbf{Answer.} Tradeoff: Larger vocabularies shorten sequences but increase embedding parameters and sparse coverage. Pitfall: Comparing context lengths across tokenizers as if tokens were equal misstates capacity and cost.
\par\textit{Tradeoff: Larger vocabularies shorten sequences but increase embedding parameters and sparse coverage. Pitfall: Comparing context lengths across tokenizers as if tokens were equal misstates capacity and cost.}
\paragraph{LA-0904 [hard]} Explain Tokenization from first principles, then show how you would evaluate and operate it in production.
\noindent\textbf{Answer.} Start with the mechanism: Tokenizers map text or bytes into model vocabulary IDs using algorithms such as BPE, WordPiece, or unigram models. Make the tradeoff explicit: Larger vocabularies shorten sequences but increase embedding parameters and sparse coverage. Define workload-specific offline metrics and a latency/cost budget, test important slices, instrument the critical path, and create a rollback criterion. The production review must explicitly guard against this failure: Comparing context lengths across tokenizers as if tokens were equal misstates capacity and cost.
\par\textit{A strong answer connects mechanism, measurable tradeoffs, failure isolation, observability, and rollback rather than listing product features.}
\paragraph{MCQ-0905 [medium]} A design requires this behavior: Learned token vectors map discrete IDs into continuous model space and are often tied to output weights. Which mechanism should the team study?
\begin{enumerate}[label=\Alph*.]
\item Multi-head attention
\item RLHF and PPO
\item Embedding layer
\item Feed-forward networks
\end{enumerate}
\noindent\textbf{Answer.} Embedding layer
\par\textit{Learned token vectors map discrete IDs into continuous model space and are often tied to output weights. Tradeoff: Weight tying saves parameters but constrains input-output geometry. Watch for: Forgetting vocabulary changes invalidates checkpoint embedding shapes.}
\paragraph{MCQ-0906 [hard]} The team proposes Embedding layer for a real workload. Which finding gives the correct decision boundary?
\begin{enumerate}[label=\Alph*.]
\item It can optimize human preferences but adds reward hacking and training complexity.
\item Diverse interactions improve capacity but add projections and KV memory.
\item Weight tying saves parameters but constrains input-output geometry.
\item They hold substantial parameters and compute while offering parallel token processing.
\end{enumerate}
\noindent\textbf{Answer.} Weight tying saves parameters but constrains input-output geometry.
\par\textit{Learned token vectors map discrete IDs into continuous model space and are often tied to output weights. Tradeoff: Weight tying saves parameters but constrains input-output geometry. Watch for: Forgetting vocabulary changes invalidates checkpoint embedding shapes.}
\paragraph{MCQ-0907 [hard]} Before releasing Embedding layer, which failure deserves a dedicated regression test?
\begin{enumerate}[label=\Alph*.]
\item Assuming heads are independently interpretable can overstate mechanistic conclusions.
\item Forgetting vocabulary changes invalidates checkpoint embedding shapes.
\item Ignoring activation memory underestimates training capacity.
\item A biased reward model can produce confidently optimized failure modes.
\end{enumerate}
\noindent\textbf{Answer.} Forgetting vocabulary changes invalidates checkpoint embedding shapes.
\par\textit{Learned token vectors map discrete IDs into continuous model space and are often tied to output weights. Tradeoff: Weight tying saves parameters but constrains input-output geometry. Watch for: Forgetting vocabulary changes invalidates checkpoint embedding shapes.}
\paragraph{MCQ-0908 [medium]} A design review reveals this mistake: Forgetting vocabulary changes invalidates checkpoint embedding shapes. Which mechanism should be traced first?
\begin{enumerate}[label=\Alph*.]
\item Multi-head attention
\item Embedding layer
\item RLHF and PPO
\item Feed-forward networks
\end{enumerate}
\noindent\textbf{Answer.} Embedding layer
\par\textit{Learned token vectors map discrete IDs into continuous model space and are often tied to output weights. Tradeoff: Weight tying saves parameters but constrains input-output geometry. Watch for: Forgetting vocabulary changes invalidates checkpoint embedding shapes.}
\paragraph{MCQ-0909 [hard]} Which causal explanation of Embedding layer connects what it does to when it is worth its cost?
\begin{enumerate}[label=\Alph*.]
\item Mechanism: Learned token vectors map discrete IDs into continuous model space and are often tied to output weights. Consequence: Weight tying saves parameters but constrains input-output geometry.
\item Mechanism: Multiple heads project attention into separate subspaces before concatenation and output projection. Consequence: Diverse interactions improve capacity but add projections and KV memory.
\item Mechanism: A preference model supplies rewards while PPO constrains policy updates relative to a reference model. Consequence: It can optimize human preferences but adds reward hacking and training complexity.
\item Mechanism: Transformer MLP blocks expand hidden dimensions, apply nonlinear gates, and project back per token. Consequence: They hold substantial parameters and compute while offering parallel token processing.
\end{enumerate}
\noindent\textbf{Answer.} Mechanism: Learned token vectors map discrete IDs into continuous model space and are often tied to output weights. Consequence: Weight tying saves parameters but constrains input-output geometry.
\par\textit{Learned token vectors map discrete IDs into continuous model space and are often tied to output weights. Tradeoff: Weight tying saves parameters but constrains input-output geometry. Watch for: Forgetting vocabulary changes invalidates checkpoint embedding shapes.}
\paragraph{FC-0910 [medium]} Define Embedding layer in interview-ready language.
\noindent\textbf{Answer.} Learned token vectors map discrete IDs into continuous model space and are often tied to output weights.
\par\textit{Learned token vectors map discrete IDs into continuous model space and are often tied to output weights.}
\paragraph{FC-0911 [hard]} State the main tradeoff and production pitfall for Embedding layer.
\noindent\textbf{Answer.} Tradeoff: Weight tying saves parameters but constrains input-output geometry. Pitfall: Forgetting vocabulary changes invalidates checkpoint embedding shapes.
\par\textit{Tradeoff: Weight tying saves parameters but constrains input-output geometry. Pitfall: Forgetting vocabulary changes invalidates checkpoint embedding shapes.}
\paragraph{LA-0912 [hard]} Explain Embedding layer from first principles, then show how you would evaluate and operate it in production.
\noindent\textbf{Answer.} Start with the mechanism: Learned token vectors map discrete IDs into continuous model space and are often tied to output weights. Make the tradeoff explicit: Weight tying saves parameters but constrains input-output geometry. Define workload-specific offline metrics and a latency/cost budget, test important slices, instrument the critical path, and create a rollback criterion. The production review must explicitly guard against this failure: Forgetting vocabulary changes invalidates checkpoint embedding shapes.
\par\textit{A strong answer connects mechanism, measurable tradeoffs, failure isolation, observability, and rollback rather than listing product features.}
\paragraph{MCQ-0913 [medium]} A design requires this behavior: Attention computes scaled query-key scores, applies masking and softmax, then mixes value vectors. Which mechanism should the team study?
\begin{enumerate}[label=\Alph*.]
\item Self-attention math
\item Mixture of Experts
\item RoPE
\item Embedding layer
\end{enumerate}
\noindent\textbf{Answer.} Self-attention math
\par\textit{Attention computes scaled query-key scores, applies masking and softmax, then mixes value vectors. Tradeoff: Every token can interact directly, while standard sequence cost grows quadratically. Watch for: Omitting the square-root head-dimension scale saturates softmax gradients.}
\paragraph{MCQ-0914 [hard]} The team proposes Self-attention math for a real workload. Which finding gives the correct decision boundary?
\begin{enumerate}[label=\Alph*.]
\item Every token can interact directly, while standard sequence cost grows quadratically.
\item Weight tying saves parameters but constrains input-output geometry.
\item RoPE integrates cleanly with attention but long-context scaling needs careful frequency treatment.
\item Capacity scales efficiently but routing balance, communication, and memory are hard.
\end{enumerate}
\noindent\textbf{Answer.} Every token can interact directly, while standard sequence cost grows quadratically.
\par\textit{Attention computes scaled query-key scores, applies masking and softmax, then mixes value vectors. Tradeoff: Every token can interact directly, while standard sequence cost grows quadratically. Watch for: Omitting the square-root head-dimension scale saturates softmax gradients.}
\paragraph{MCQ-0915 [hard]} Before releasing Self-attention math, which failure deserves a dedicated regression test?
\begin{enumerate}[label=\Alph*.]
\item Expert collapse and uneven routing create stragglers and weak specialization.
\item Forgetting vocabulary changes invalidates checkpoint embedding shapes.
\item Omitting the square-root head-dimension scale saturates softmax gradients.
\item Naively increasing the maximum length does not preserve trained frequency behavior.
\end{enumerate}
\noindent\textbf{Answer.} Omitting the square-root head-dimension scale saturates softmax gradients.
\par\textit{Attention computes scaled query-key scores, applies masking and softmax, then mixes value vectors. Tradeoff: Every token can interact directly, while standard sequence cost grows quadratically. Watch for: Omitting the square-root head-dimension scale saturates softmax gradients.}
\paragraph{MCQ-0916 [medium]} A design review reveals this mistake: Omitting the square-root head-dimension scale saturates softmax gradients. Which mechanism should be traced first?
\begin{enumerate}[label=\Alph*.]
\item Self-attention math
\item Mixture of Experts
\item Embedding layer
\item RoPE
\end{enumerate}
\noindent\textbf{Answer.} Self-attention math
\par\textit{Attention computes scaled query-key scores, applies masking and softmax, then mixes value vectors. Tradeoff: Every token can interact directly, while standard sequence cost grows quadratically. Watch for: Omitting the square-root head-dimension scale saturates softmax gradients.}
\paragraph{MCQ-0917 [hard]} Which causal explanation of Self-attention math connects what it does to when it is worth its cost?
\begin{enumerate}[label=\Alph*.]
\item Mechanism: Attention computes scaled query-key scores, applies masking and softmax, then mixes value vectors. Consequence: Every token can interact directly, while standard sequence cost grows quadratically.
\item Mechanism: MoE routes each token to a subset of expert networks, increasing parameters without proportional per-token compute. Consequence: Capacity scales efficiently but routing balance, communication, and memory are hard.
\item Mechanism: Learned token vectors map discrete IDs into continuous model space and are often tied to output weights. Consequence: Weight tying saves parameters but constrains input-output geometry.
\item Mechanism: Rotary position embeddings rotate query and key coordinates so dot products encode relative offsets. Consequence: RoPE integrates cleanly with attention but long-context scaling needs careful frequency treatment.
\end{enumerate}
\noindent\textbf{Answer.} Mechanism: Attention computes scaled query-key scores, applies masking and softmax, then mixes value vectors. Consequence: Every token can interact directly, while standard sequence cost grows quadratically.
\par\textit{Attention computes scaled query-key scores, applies masking and softmax, then mixes value vectors. Tradeoff: Every token can interact directly, while standard sequence cost grows quadratically. Watch for: Omitting the square-root head-dimension scale saturates softmax gradients.}
\paragraph{FC-0918 [medium]} Define Self-attention math in interview-ready language.
\noindent\textbf{Answer.} Attention computes scaled query-key scores, applies masking and softmax, then mixes value vectors.
\par\textit{Attention computes scaled query-key scores, applies masking and softmax, then mixes value vectors.}
\paragraph{FC-0919 [hard]} State the main tradeoff and production pitfall for Self-attention math.
\noindent\textbf{Answer.} Tradeoff: Every token can interact directly, while standard sequence cost grows quadratically. Pitfall: Omitting the square-root head-dimension scale saturates softmax gradients.
\par\textit{Tradeoff: Every token can interact directly, while standard sequence cost grows quadratically. Pitfall: Omitting the square-root head-dimension scale saturates softmax gradients.}
\paragraph{LA-0920 [hard]} Explain Self-attention math from first principles, then show how you would evaluate and operate it in production.
\noindent\textbf{Answer.} Start with the mechanism: Attention computes scaled query-key scores, applies masking and softmax, then mixes value vectors. Make the tradeoff explicit: Every token can interact directly, while standard sequence cost grows quadratically. Define workload-specific offline metrics and a latency/cost budget, test important slices, instrument the critical path, and create a rollback criterion. The production review must explicitly guard against this failure: Omitting the square-root head-dimension scale saturates softmax gradients.
\par\textit{A strong answer connects mechanism, measurable tradeoffs, failure isolation, observability, and rollback rather than listing product features.}
\paragraph{MCQ-0921 [medium]} A design requires this behavior: Multiple heads project attention into separate subspaces before concatenation and output projection. Which mechanism should the team study?
\begin{enumerate}[label=\Alph*.]
\item RoPE
\item Layer normalization and RMSNorm
\item Multi-head attention
\item Residual connections
\end{enumerate}
\noindent\textbf{Answer.} Multi-head attention
\par\textit{Multiple heads project attention into separate subspaces before concatenation and output projection. Tradeoff: Diverse interactions improve capacity but add projections and KV memory. Watch for: Assuming heads are independently interpretable can overstate mechanistic conclusions.}
\paragraph{MCQ-0922 [hard]} The team proposes Multi-head attention for a real workload. Which finding gives the correct decision boundary?
\begin{enumerate}[label=\Alph*.]
\item Simpler normalization can be faster but architecture placement changes optimization dynamics.
\item Diverse interactions improve capacity but add projections and KV memory.
\item RoPE integrates cleanly with attention but long-context scaling needs careful frequency treatment.
\item Deep optimization becomes feasible, while residual scale can grow and require controls.
\end{enumerate}
\noindent\textbf{Answer.} Diverse interactions improve capacity but add projections and KV memory.
\par\textit{Multiple heads project attention into separate subspaces before concatenation and output projection. Tradeoff: Diverse interactions improve capacity but add projections and KV memory. Watch for: Assuming heads are independently interpretable can overstate mechanistic conclusions.}
\paragraph{MCQ-0923 [hard]} Before releasing Multi-head attention, which failure deserves a dedicated regression test?
\begin{enumerate}[label=\Alph*.]
\item Assuming heads are independently interpretable can overstate mechanistic conclusions.
\item In-place mutation of residual tensors can break autograd or checkpointing.
\item Moving pre-norm to post-norm without retuning can destabilize deep training.
\item Naively increasing the maximum length does not preserve trained frequency behavior.
\end{enumerate}
\noindent\textbf{Answer.} Assuming heads are independently interpretable can overstate mechanistic conclusions.
\par\textit{Multiple heads project attention into separate subspaces before concatenation and output projection. Tradeoff: Diverse interactions improve capacity but add projections and KV memory. Watch for: Assuming heads are independently interpretable can overstate mechanistic conclusions.}
\paragraph{MCQ-0924 [medium]} A design review reveals this mistake: Assuming heads are independently interpretable can overstate mechanistic conclusions. Which mechanism should be traced first?
\begin{enumerate}[label=\Alph*.]
\item Multi-head attention
\item RoPE
\item Layer normalization and RMSNorm
\item Residual connections
\end{enumerate}
\noindent\textbf{Answer.} Multi-head attention
\par\textit{Multiple heads project attention into separate subspaces before concatenation and output projection. Tradeoff: Diverse interactions improve capacity but add projections and KV memory. Watch for: Assuming heads are independently interpretable can overstate mechanistic conclusions.}
\paragraph{MCQ-0925 [hard]} Which causal explanation of Multi-head attention connects what it does to when it is worth its cost?
\begin{enumerate}[label=\Alph*.]
\item Mechanism: Normalization stabilizes activations; RMSNorm removes mean centering and normalizes by root-mean-square magnitude. Consequence: Simpler normalization can be faster but architecture placement changes optimization dynamics.
\item Mechanism: Rotary position embeddings rotate query and key coordinates so dot products encode relative offsets. Consequence: RoPE integrates cleanly with attention but long-context scaling needs careful frequency treatment.
\item Mechanism: Multiple heads project attention into separate subspaces before concatenation and output projection. Consequence: Diverse interactions improve capacity but add projections and KV memory.
\item Mechanism: Residual paths add transformed activations back to the stream to improve gradient flow and feature reuse. Consequence: Deep optimization becomes feasible, while residual scale can grow and require controls.
\end{enumerate}
\noindent\textbf{Answer.} Mechanism: Multiple heads project attention into separate subspaces before concatenation and output projection. Consequence: Diverse interactions improve capacity but add projections and KV memory.
\par\textit{Multiple heads project attention into separate subspaces before concatenation and output projection. Tradeoff: Diverse interactions improve capacity but add projections and KV memory. Watch for: Assuming heads are independently interpretable can overstate mechanistic conclusions.}
\paragraph{FC-0926 [medium]} Define Multi-head attention in interview-ready language.
\noindent\textbf{Answer.} Multiple heads project attention into separate subspaces before concatenation and output projection.
\par\textit{Multiple heads project attention into separate subspaces before concatenation and output projection.}
\paragraph{FC-0927 [hard]} State the main tradeoff and production pitfall for Multi-head attention.
\noindent\textbf{Answer.} Tradeoff: Diverse interactions improve capacity but add projections and KV memory. Pitfall: Assuming heads are independently interpretable can overstate mechanistic conclusions.
\par\textit{Tradeoff: Diverse interactions improve capacity but add projections and KV memory. Pitfall: Assuming heads are independently interpretable can overstate mechanistic conclusions.}
\paragraph{LA-0928 [hard]} Explain Multi-head attention from first principles, then show how you would evaluate and operate it in production.
\noindent\textbf{Answer.} Start with the mechanism: Multiple heads project attention into separate subspaces before concatenation and output projection. Make the tradeoff explicit: Diverse interactions improve capacity but add projections and KV memory. Define workload-specific offline metrics and a latency/cost budget, test important slices, instrument the critical path, and create a rollback criterion. The production review must explicitly guard against this failure: Assuming heads are independently interpretable can overstate mechanistic conclusions.
\par\textit{A strong answer connects mechanism, measurable tradeoffs, failure isolation, observability, and rollback rather than listing product features.}
\paragraph{MCQ-0929 [medium]} A design requires this behavior: Position signals break permutation invariance using absolute, relative, rotary, or bias-based schemes. Which mechanism should the team study?
\begin{enumerate}[label=\Alph*.]
\item SFT
\item Positional encoding
\item Knowledge distillation
\item RLHF and PPO
\end{enumerate}
\noindent\textbf{Answer.} Positional encoding
\par\textit{Position signals break permutation invariance using absolute, relative, rotary, or bias-based schemes. Tradeoff: Relative methods extrapolate differently, while implementation and scaling choices affect long context. Watch for: Extending context without validating the position scheme can collapse quality.}
\paragraph{MCQ-0930 [hard]} The team proposes Positional encoding for a real workload. Which finding gives the correct decision boundary?
\begin{enumerate}[label=\Alph*.]
\item It can optimize human preferences but adds reward hacking and training complexity.
\item Relative methods extrapolate differently, while implementation and scaling choices affect long context.
\item Serving cost falls but rare capabilities and calibration may be lost.
\item It is direct and stable but can overfit style and inherit annotation errors.
\end{enumerate}
\noindent\textbf{Answer.} Relative methods extrapolate differently, while implementation and scaling choices affect long context.
\par\textit{Position signals break permutation invariance using absolute, relative, rotary, or bias-based schemes. Tradeoff: Relative methods extrapolate differently, while implementation and scaling choices affect long context. Watch for: Extending context without validating the position scheme can collapse quality.}
\paragraph{MCQ-0931 [hard]} Before releasing Positional encoding, which failure deserves a dedicated regression test?
\begin{enumerate}[label=\Alph*.]
\item Extending context without validating the position scheme can collapse quality.
\item Distilling only easy teacher outputs creates a deceptively strong narrow model.
\item A biased reward model can produce confidently optimized failure modes.
\item Training on conflicting instructions without provenance makes behavior unstable.
\end{enumerate}
\noindent\textbf{Answer.} Extending context without validating the position scheme can collapse quality.
\par\textit{Position signals break permutation invariance using absolute, relative, rotary, or bias-based schemes. Tradeoff: Relative methods extrapolate differently, while implementation and scaling choices affect long context. Watch for: Extending context without validating the position scheme can collapse quality.}
\paragraph{MCQ-0932 [medium]} A design review reveals this mistake: Extending context without validating the position scheme can collapse quality. Which mechanism should be traced first?
\begin{enumerate}[label=\Alph*.]
\item Knowledge distillation
\item SFT
\item Positional encoding
\item RLHF and PPO
\end{enumerate}
\noindent\textbf{Answer.} Positional encoding
\par\textit{Position signals break permutation invariance using absolute, relative, rotary, or bias-based schemes. Tradeoff: Relative methods extrapolate differently, while implementation and scaling choices affect long context. Watch for: Extending context without validating the position scheme can collapse quality.}
\paragraph{MCQ-0933 [hard]} Which causal explanation of Positional encoding connects what it does to when it is worth its cost?
\begin{enumerate}[label=\Alph*.]
\item Mechanism: Position signals break permutation invariance using absolute, relative, rotary, or bias-based schemes. Consequence: Relative methods extrapolate differently, while implementation and scaling choices affect long context.
\item Mechanism: A student learns from teacher probabilities, generated data, rationales, or intermediate representations. Consequence: Serving cost falls but rare capabilities and calibration may be lost.
\item Mechanism: A preference model supplies rewards while PPO constrains policy updates relative to a reference model. Consequence: It can optimize human preferences but adds reward hacking and training complexity.
\item Mechanism: Supervised fine-tuning teaches instruction-response behavior from curated demonstrations. Consequence: It is direct and stable but can overfit style and inherit annotation errors.
\end{enumerate}
\noindent\textbf{Answer.} Mechanism: Position signals break permutation invariance using absolute, relative, rotary, or bias-based schemes. Consequence: Relative methods extrapolate differently, while implementation and scaling choices affect long context.
\par\textit{Position signals break permutation invariance using absolute, relative, rotary, or bias-based schemes. Tradeoff: Relative methods extrapolate differently, while implementation and scaling choices affect long context. Watch for: Extending context without validating the position scheme can collapse quality.}
\paragraph{FC-0934 [medium]} Define Positional encoding in interview-ready language.
\noindent\textbf{Answer.} Position signals break permutation invariance using absolute, relative, rotary, or bias-based schemes.
\par\textit{Position signals break permutation invariance using absolute, relative, rotary, or bias-based schemes.}
\paragraph{FC-0935 [hard]} State the main tradeoff and production pitfall for Positional encoding.
\noindent\textbf{Answer.} Tradeoff: Relative methods extrapolate differently, while implementation and scaling choices affect long context. Pitfall: Extending context without validating the position scheme can collapse quality.
\par\textit{Tradeoff: Relative methods extrapolate differently, while implementation and scaling choices affect long context. Pitfall: Extending context without validating the position scheme can collapse quality.}
\paragraph{LA-0936 [hard]} Explain Positional encoding from first principles, then show how you would evaluate and operate it in production.
\noindent\textbf{Answer.} Start with the mechanism: Position signals break permutation invariance using absolute, relative, rotary, or bias-based schemes. Make the tradeoff explicit: Relative methods extrapolate differently, while implementation and scaling choices affect long context. Define workload-specific offline metrics and a latency/cost budget, test important slices, instrument the critical path, and create a rollback criterion. The production review must explicitly guard against this failure: Extending context without validating the position scheme can collapse quality.
\par\textit{A strong answer connects mechanism, measurable tradeoffs, failure isolation, observability, and rollback rather than listing product features.}
\paragraph{MCQ-0937 [medium]} A design requires this behavior: Rotary position embeddings rotate query and key coordinates so dot products encode relative offsets. Which mechanism should the team study?
\begin{enumerate}[label=\Alph*.]
\item Causal language modeling
\item Multi-head attention
\item RoPE
\item Positional encoding
\end{enumerate}
\noindent\textbf{Answer.} RoPE
\par\textit{Rotary position embeddings rotate query and key coordinates so dot products encode relative offsets. Tradeoff: RoPE integrates cleanly with attention but long-context scaling needs careful frequency treatment. Watch for: Naively increasing the maximum length does not preserve trained frequency behavior.}
\paragraph{MCQ-0938 [hard]} The team proposes RoPE for a real workload. Which finding gives the correct decision boundary?
\begin{enumerate}[label=\Alph*.]
\item It provides a simple scalable objective but does not directly optimize instruction usefulness or truth.
\item RoPE integrates cleanly with attention but long-context scaling needs careful frequency treatment.
\item Diverse interactions improve capacity but add projections and KV memory.
\item Relative methods extrapolate differently, while implementation and scaling choices affect long context.
\end{enumerate}
\noindent\textbf{Answer.} RoPE integrates cleanly with attention but long-context scaling needs careful frequency treatment.
\par\textit{Rotary position embeddings rotate query and key coordinates so dot products encode relative offsets. Tradeoff: RoPE integrates cleanly with attention but long-context scaling needs careful frequency treatment. Watch for: Naively increasing the maximum length does not preserve trained frequency behavior.}
\paragraph{MCQ-0939 [hard]} Before releasing RoPE, which failure deserves a dedicated regression test?
\begin{enumerate}[label=\Alph*.]
\item Assuming heads are independently interpretable can overstate mechanistic conclusions.
\item Extending context without validating the position scheme can collapse quality.
\item Interpreting low perplexity as safe task performance is invalid.
\item Naively increasing the maximum length does not preserve trained frequency behavior.
\end{enumerate}
\noindent\textbf{Answer.} Naively increasing the maximum length does not preserve trained frequency behavior.
\par\textit{Rotary position embeddings rotate query and key coordinates so dot products encode relative offsets. Tradeoff: RoPE integrates cleanly with attention but long-context scaling needs careful frequency treatment. Watch for: Naively increasing the maximum length does not preserve trained frequency behavior.}
\paragraph{MCQ-0940 [medium]} A design review reveals this mistake: Naively increasing the maximum length does not preserve trained frequency behavior. Which mechanism should be traced first?
\begin{enumerate}[label=\Alph*.]
\item RoPE
\item Causal language modeling
\item Positional encoding
\item Multi-head attention
\end{enumerate}
\noindent\textbf{Answer.} RoPE
\par\textit{Rotary position embeddings rotate query and key coordinates so dot products encode relative offsets. Tradeoff: RoPE integrates cleanly with attention but long-context scaling needs careful frequency treatment. Watch for: Naively increasing the maximum length does not preserve trained frequency behavior.}
\paragraph{MCQ-0941 [hard]} Which causal explanation of RoPE connects what it does to when it is worth its cost?
\begin{enumerate}[label=\Alph*.]
\item Mechanism: Next-token training factorizes sequence likelihood from left to right under a causal mask. Consequence: It provides a simple scalable objective but does not directly optimize instruction usefulness or truth.
\item Mechanism: Position signals break permutation invariance using absolute, relative, rotary, or bias-based schemes. Consequence: Relative methods extrapolate differently, while implementation and scaling choices affect long context.
\item Mechanism: Multiple heads project attention into separate subspaces before concatenation and output projection. Consequence: Diverse interactions improve capacity but add projections and KV memory.
\item Mechanism: Rotary position embeddings rotate query and key coordinates so dot products encode relative offsets. Consequence: RoPE integrates cleanly with attention but long-context scaling needs careful frequency treatment.
\end{enumerate}
\noindent\textbf{Answer.} Mechanism: Rotary position embeddings rotate query and key coordinates so dot products encode relative offsets. Consequence: RoPE integrates cleanly with attention but long-context scaling needs careful frequency treatment.
\par\textit{Rotary position embeddings rotate query and key coordinates so dot products encode relative offsets. Tradeoff: RoPE integrates cleanly with attention but long-context scaling needs careful frequency treatment. Watch for: Naively increasing the maximum length does not preserve trained frequency behavior.}
\paragraph{FC-0942 [medium]} Define RoPE in interview-ready language.
\noindent\textbf{Answer.} Rotary position embeddings rotate query and key coordinates so dot products encode relative offsets.
\par\textit{Rotary position embeddings rotate query and key coordinates so dot products encode relative offsets.}
\paragraph{FC-0943 [hard]} State the main tradeoff and production pitfall for RoPE.
\noindent\textbf{Answer.} Tradeoff: RoPE integrates cleanly with attention but long-context scaling needs careful frequency treatment. Pitfall: Naively increasing the maximum length does not preserve trained frequency behavior.
\par\textit{Tradeoff: RoPE integrates cleanly with attention but long-context scaling needs careful frequency treatment. Pitfall: Naively increasing the maximum length does not preserve trained frequency behavior.}
\paragraph{LA-0944 [hard]} Explain RoPE from first principles, then show how you would evaluate and operate it in production.
\noindent\textbf{Answer.} Start with the mechanism: Rotary position embeddings rotate query and key coordinates so dot products encode relative offsets. Make the tradeoff explicit: RoPE integrates cleanly with attention but long-context scaling needs careful frequency treatment. Define workload-specific offline metrics and a latency/cost budget, test important slices, instrument the critical path, and create a rollback criterion. The production review must explicitly guard against this failure: Naively increasing the maximum length does not preserve trained frequency behavior.
\par\textit{A strong answer connects mechanism, measurable tradeoffs, failure isolation, observability, and rollback rather than listing product features.}
\paragraph{MCQ-0945 [medium]} A design requires this behavior: Transformer MLP blocks expand hidden dimensions, apply nonlinear gates, and project back per token. Which mechanism should the team study?
\begin{enumerate}[label=\Alph*.]
\item RoPE
\item Feed-forward networks
\item Tokenization
\item Layer normalization and RMSNorm
\end{enumerate}
\noindent\textbf{Answer.} Feed-forward networks
\par\textit{Transformer MLP blocks expand hidden dimensions, apply nonlinear gates, and project back per token. Tradeoff: They hold substantial parameters and compute while offering parallel token processing. Watch for: Ignoring activation memory underestimates training capacity.}
\paragraph{MCQ-0946 [hard]} The team proposes Feed-forward networks for a real workload. Which finding gives the correct decision boundary?
\begin{enumerate}[label=\Alph*.]
\item They hold substantial parameters and compute while offering parallel token processing.
\item Simpler normalization can be faster but architecture placement changes optimization dynamics.
\item Larger vocabularies shorten sequences but increase embedding parameters and sparse coverage.
\item RoPE integrates cleanly with attention but long-context scaling needs careful frequency treatment.
\end{enumerate}
\noindent\textbf{Answer.} They hold substantial parameters and compute while offering parallel token processing.
\par\textit{Transformer MLP blocks expand hidden dimensions, apply nonlinear gates, and project back per token. Tradeoff: They hold substantial parameters and compute while offering parallel token processing. Watch for: Ignoring activation memory underestimates training capacity.}
\paragraph{MCQ-0947 [hard]} Before releasing Feed-forward networks, which failure deserves a dedicated regression test?
\begin{enumerate}[label=\Alph*.]
\item Comparing context lengths across tokenizers as if tokens were equal misstates capacity and cost.
\item Naively increasing the maximum length does not preserve trained frequency behavior.
\item Moving pre-norm to post-norm without retuning can destabilize deep training.
\item Ignoring activation memory underestimates training capacity.
\end{enumerate}
\noindent\textbf{Answer.} Ignoring activation memory underestimates training capacity.
\par\textit{Transformer MLP blocks expand hidden dimensions, apply nonlinear gates, and project back per token. Tradeoff: They hold substantial parameters and compute while offering parallel token processing. Watch for: Ignoring activation memory underestimates training capacity.}
\paragraph{MCQ-0948 [medium]} A design review reveals this mistake: Ignoring activation memory underestimates training capacity. Which mechanism should be traced first?
\begin{enumerate}[label=\Alph*.]
\item RoPE
\item Layer normalization and RMSNorm
\item Tokenization
\item Feed-forward networks
\end{enumerate}
\noindent\textbf{Answer.} Feed-forward networks
\par\textit{Transformer MLP blocks expand hidden dimensions, apply nonlinear gates, and project back per token. Tradeoff: They hold substantial parameters and compute while offering parallel token processing. Watch for: Ignoring activation memory underestimates training capacity.}
\paragraph{MCQ-0949 [hard]} Which causal explanation of Feed-forward networks connects what it does to when it is worth its cost?
\begin{enumerate}[label=\Alph*.]
\item Mechanism: Rotary position embeddings rotate query and key coordinates so dot products encode relative offsets. Consequence: RoPE integrates cleanly with attention but long-context scaling needs careful frequency treatment.
\item Mechanism: Tokenizers map text or bytes into model vocabulary IDs using algorithms such as BPE, WordPiece, or unigram models. Consequence: Larger vocabularies shorten sequences but increase embedding parameters and sparse coverage.
\item Mechanism: Transformer MLP blocks expand hidden dimensions, apply nonlinear gates, and project back per token. Consequence: They hold substantial parameters and compute while offering parallel token processing.
\item Mechanism: Normalization stabilizes activations; RMSNorm removes mean centering and normalizes by root-mean-square magnitude. Consequence: Simpler normalization can be faster but architecture placement changes optimization dynamics.
\end{enumerate}
\noindent\textbf{Answer.} Mechanism: Transformer MLP blocks expand hidden dimensions, apply nonlinear gates, and project back per token. Consequence: They hold substantial parameters and compute while offering parallel token processing.
\par\textit{Transformer MLP blocks expand hidden dimensions, apply nonlinear gates, and project back per token. Tradeoff: They hold substantial parameters and compute while offering parallel token processing. Watch for: Ignoring activation memory underestimates training capacity.}
\paragraph{FC-0950 [medium]} Define Feed-forward networks in interview-ready language.
\noindent\textbf{Answer.} Transformer MLP blocks expand hidden dimensions, apply nonlinear gates, and project back per token.
\par\textit{Transformer MLP blocks expand hidden dimensions, apply nonlinear gates, and project back per token.}
\paragraph{FC-0951 [hard]} State the main tradeoff and production pitfall for Feed-forward networks.
\noindent\textbf{Answer.} Tradeoff: They hold substantial parameters and compute while offering parallel token processing. Pitfall: Ignoring activation memory underestimates training capacity.
\par\textit{Tradeoff: They hold substantial parameters and compute while offering parallel token processing. Pitfall: Ignoring activation memory underestimates training capacity.}
\paragraph{LA-0952 [hard]} Explain Feed-forward networks from first principles, then show how you would evaluate and operate it in production.
\noindent\textbf{Answer.} Start with the mechanism: Transformer MLP blocks expand hidden dimensions, apply nonlinear gates, and project back per token. Make the tradeoff explicit: They hold substantial parameters and compute while offering parallel token processing. Define workload-specific offline metrics and a latency/cost budget, test important slices, instrument the critical path, and create a rollback criterion. The production review must explicitly guard against this failure: Ignoring activation memory underestimates training capacity.
\par\textit{A strong answer connects mechanism, measurable tradeoffs, failure isolation, observability, and rollback rather than listing product features.}
\paragraph{MCQ-0953 [medium]} A design requires this behavior: Normalization stabilizes activations; RMSNorm removes mean centering and normalizes by root-mean-square magnitude. Which mechanism should the team study?
\begin{enumerate}[label=\Alph*.]
\item QLoRA
\item LoRA
\item Layer normalization and RMSNorm
\item RoPE
\end{enumerate}
\noindent\textbf{Answer.} Layer normalization and RMSNorm
\par\textit{Normalization stabilizes activations; RMSNorm removes mean centering and normalizes by root-mean-square magnitude. Tradeoff: Simpler normalization can be faster but architecture placement changes optimization dynamics. Watch for: Moving pre-norm to post-norm without retuning can destabilize deep training.}
\paragraph{MCQ-0954 [hard]} The team proposes Layer normalization and RMSNorm for a real workload. Which finding gives the correct decision boundary?
\begin{enumerate}[label=\Alph*.]
\item Fine-tuning large models becomes cheaper, while kernels and quantization settings affect stability.
\item Simpler normalization can be faster but architecture placement changes optimization dynamics.
\item Memory and storage fall, while rank and target-module choices limit expressiveness.
\item RoPE integrates cleanly with attention but long-context scaling needs careful frequency treatment.
\end{enumerate}
\noindent\textbf{Answer.} Simpler normalization can be faster but architecture placement changes optimization dynamics.
\par\textit{Normalization stabilizes activations; RMSNorm removes mean centering and normalizes by root-mean-square magnitude. Tradeoff: Simpler normalization can be faster but architecture placement changes optimization dynamics. Watch for: Moving pre-norm to post-norm without retuning can destabilize deep training.}
\paragraph{MCQ-0955 [hard]} Before releasing Layer normalization and RMSNorm, which failure deserves a dedicated regression test?
\begin{enumerate}[label=\Alph*.]
\item Moving pre-norm to post-norm without retuning can destabilize deep training.
\item Merging incompatible adapters without evaluation creates interference.
\item Confusing training quantization with final serving quantization produces wrong expectations.
\item Naively increasing the maximum length does not preserve trained frequency behavior.
\end{enumerate}
\noindent\textbf{Answer.} Moving pre-norm to post-norm without retuning can destabilize deep training.
\par\textit{Normalization stabilizes activations; RMSNorm removes mean centering and normalizes by root-mean-square magnitude. Tradeoff: Simpler normalization can be faster but architecture placement changes optimization dynamics. Watch for: Moving pre-norm to post-norm without retuning can destabilize deep training.}
\paragraph{MCQ-0956 [medium]} A design review reveals this mistake: Moving pre-norm to post-norm without retuning can destabilize deep training. Which mechanism should be traced first?
\begin{enumerate}[label=\Alph*.]
\item RoPE
\item LoRA
\item Layer normalization and RMSNorm
\item QLoRA
\end{enumerate}
\noindent\textbf{Answer.} Layer normalization and RMSNorm
\par\textit{Normalization stabilizes activations; RMSNorm removes mean centering and normalizes by root-mean-square magnitude. Tradeoff: Simpler normalization can be faster but architecture placement changes optimization dynamics. Watch for: Moving pre-norm to post-norm without retuning can destabilize deep training.}
\paragraph{MCQ-0957 [hard]} Which causal explanation of Layer normalization and RMSNorm connects what it does to when it is worth its cost?
\begin{enumerate}[label=\Alph*.]
\item Mechanism: LoRA trains low-rank update matrices while freezing base weights for parameter-efficient adaptation. Consequence: Memory and storage fall, while rank and target-module choices limit expressiveness.
\item Mechanism: QLoRA backpropagates through a frozen quantized base model into LoRA adapters using memory-saving quantization techniques. Consequence: Fine-tuning large models becomes cheaper, while kernels and quantization settings affect stability.
\item Mechanism: Normalization stabilizes activations; RMSNorm removes mean centering and normalizes by root-mean-square magnitude. Consequence: Simpler normalization can be faster but architecture placement changes optimization dynamics.
\item Mechanism: Rotary position embeddings rotate query and key coordinates so dot products encode relative offsets. Consequence: RoPE integrates cleanly with attention but long-context scaling needs careful frequency treatment.
\end{enumerate}
\noindent\textbf{Answer.} Mechanism: Normalization stabilizes activations; RMSNorm removes mean centering and normalizes by root-mean-square magnitude. Consequence: Simpler normalization can be faster but architecture placement changes optimization dynamics.
\par\textit{Normalization stabilizes activations; RMSNorm removes mean centering and normalizes by root-mean-square magnitude. Tradeoff: Simpler normalization can be faster but architecture placement changes optimization dynamics. Watch for: Moving pre-norm to post-norm without retuning can destabilize deep training.}
\paragraph{FC-0958 [medium]} Define Layer normalization and RMSNorm in interview-ready language.
\noindent\textbf{Answer.} Normalization stabilizes activations; RMSNorm removes mean centering and normalizes by root-mean-square magnitude.
\par\textit{Normalization stabilizes activations; RMSNorm removes mean centering and normalizes by root-mean-square magnitude.}
\paragraph{FC-0959 [hard]} State the main tradeoff and production pitfall for Layer normalization and RMSNorm.
\noindent\textbf{Answer.} Tradeoff: Simpler normalization can be faster but architecture placement changes optimization dynamics. Pitfall: Moving pre-norm to post-norm without retuning can destabilize deep training.
\par\textit{Tradeoff: Simpler normalization can be faster but architecture placement changes optimization dynamics. Pitfall: Moving pre-norm to post-norm without retuning can destabilize deep training.}
\paragraph{LA-0960 [hard]} Explain Layer normalization and RMSNorm from first principles, then show how you would evaluate and operate it in production.
\noindent\textbf{Answer.} Start with the mechanism: Normalization stabilizes activations; RMSNorm removes mean centering and normalizes by root-mean-square magnitude. Make the tradeoff explicit: Simpler normalization can be faster but architecture placement changes optimization dynamics. Define workload-specific offline metrics and a latency/cost budget, test important slices, instrument the critical path, and create a rollback criterion. The production review must explicitly guard against this failure: Moving pre-norm to post-norm without retuning can destabilize deep training.
\par\textit{A strong answer connects mechanism, measurable tradeoffs, failure isolation, observability, and rollback rather than listing product features.}
\paragraph{MCQ-0961 [medium]} A design requires this behavior: Residual paths add transformed activations back to the stream to improve gradient flow and feature reuse. Which mechanism should the team study?
\begin{enumerate}[label=\Alph*.]
\item Residual connections
\item Cross-entropy and perplexity
\item RoPE
\item SFT
\end{enumerate}
\noindent\textbf{Answer.} Residual connections
\par\textit{Residual paths add transformed activations back to the stream to improve gradient flow and feature reuse. Tradeoff: Deep optimization becomes feasible, while residual scale can grow and require controls. Watch for: In-place mutation of residual tensors can break autograd or checkpointing.}
\paragraph{MCQ-0962 [hard]} The team proposes Residual connections for a real workload. Which finding gives the correct decision boundary?
\begin{enumerate}[label=\Alph*.]
\item RoPE integrates cleanly with attention but long-context scaling needs careful frequency treatment.
\item It is direct and stable but can overfit style and inherit annotation errors.
\item Deep optimization becomes feasible, while residual scale can grow and require controls.
\item It is stable for model comparison on the same tokenization and data, not across arbitrary setups.
\end{enumerate}
\noindent\textbf{Answer.} Deep optimization becomes feasible, while residual scale can grow and require controls.
\par\textit{Residual paths add transformed activations back to the stream to improve gradient flow and feature reuse. Tradeoff: Deep optimization becomes feasible, while residual scale can grow and require controls. Watch for: In-place mutation of residual tensors can break autograd or checkpointing.}
\paragraph{MCQ-0963 [hard]} Before releasing Residual connections, which failure deserves a dedicated regression test?
\begin{enumerate}[label=\Alph*.]
\item Comparing perplexity across vocabularies or tokenizers is misleading.
\item In-place mutation of residual tensors can break autograd or checkpointing.
\item Training on conflicting instructions without provenance makes behavior unstable.
\item Naively increasing the maximum length does not preserve trained frequency behavior.
\end{enumerate}
\noindent\textbf{Answer.} In-place mutation of residual tensors can break autograd or checkpointing.
\par\textit{Residual paths add transformed activations back to the stream to improve gradient flow and feature reuse. Tradeoff: Deep optimization becomes feasible, while residual scale can grow and require controls. Watch for: In-place mutation of residual tensors can break autograd or checkpointing.}
\paragraph{MCQ-0964 [medium]} A design review reveals this mistake: In-place mutation of residual tensors can break autograd or checkpointing. Which mechanism should be traced first?
\begin{enumerate}[label=\Alph*.]
\item Residual connections
\item Cross-entropy and perplexity
\item RoPE
\item SFT
\end{enumerate}
\noindent\textbf{Answer.} Residual connections
\par\textit{Residual paths add transformed activations back to the stream to improve gradient flow and feature reuse. Tradeoff: Deep optimization becomes feasible, while residual scale can grow and require controls. Watch for: In-place mutation of residual tensors can break autograd or checkpointing.}
\paragraph{MCQ-0965 [hard]} Which causal explanation of Residual connections connects what it does to when it is worth its cost?
\begin{enumerate}[label=\Alph*.]
\item Mechanism: Residual paths add transformed activations back to the stream to improve gradient flow and feature reuse. Consequence: Deep optimization becomes feasible, while residual scale can grow and require controls.
\item Mechanism: Supervised fine-tuning teaches instruction-response behavior from curated demonstrations. Consequence: It is direct and stable but can overfit style and inherit annotation errors.
\item Mechanism: Rotary position embeddings rotate query and key coordinates so dot products encode relative offsets. Consequence: RoPE integrates cleanly with attention but long-context scaling needs careful frequency treatment.
\item Mechanism: Cross-entropy is negative log likelihood; perplexity exponentiates average token loss. Consequence: It is stable for model comparison on the same tokenization and data, not across arbitrary setups.
\end{enumerate}
\noindent\textbf{Answer.} Mechanism: Residual paths add transformed activations back to the stream to improve gradient flow and feature reuse. Consequence: Deep optimization becomes feasible, while residual scale can grow and require controls.
\par\textit{Residual paths add transformed activations back to the stream to improve gradient flow and feature reuse. Tradeoff: Deep optimization becomes feasible, while residual scale can grow and require controls. Watch for: In-place mutation of residual tensors can break autograd or checkpointing.}
\paragraph{FC-0966 [medium]} Define Residual connections in interview-ready language.
\noindent\textbf{Answer.} Residual paths add transformed activations back to the stream to improve gradient flow and feature reuse.
\par\textit{Residual paths add transformed activations back to the stream to improve gradient flow and feature reuse.}
\paragraph{FC-0967 [hard]} State the main tradeoff and production pitfall for Residual connections.
\noindent\textbf{Answer.} Tradeoff: Deep optimization becomes feasible, while residual scale can grow and require controls. Pitfall: In-place mutation of residual tensors can break autograd or checkpointing.
\par\textit{Tradeoff: Deep optimization becomes feasible, while residual scale can grow and require controls. Pitfall: In-place mutation of residual tensors can break autograd or checkpointing.}
\paragraph{LA-0968 [hard]} Explain Residual connections from first principles, then show how you would evaluate and operate it in production.
\noindent\textbf{Answer.} Start with the mechanism: Residual paths add transformed activations back to the stream to improve gradient flow and feature reuse. Make the tradeoff explicit: Deep optimization becomes feasible, while residual scale can grow and require controls. Define workload-specific offline metrics and a latency/cost budget, test important slices, instrument the critical path, and create a rollback criterion. The production review must explicitly guard against this failure: In-place mutation of residual tensors can break autograd or checkpointing.
\par\textit{A strong answer connects mechanism, measurable tradeoffs, failure isolation, observability, and rollback rather than listing product features.}
\paragraph{MCQ-0969 [medium]} A design requires this behavior: MoE routes each token to a subset of expert networks, increasing parameters without proportional per-token compute. Which mechanism should the team study?
\begin{enumerate}[label=\Alph*.]
\item Self-attention math
\item RLHF and PPO
\item Cross-entropy and perplexity
\item Mixture of Experts
\end{enumerate}
\noindent\textbf{Answer.} Mixture of Experts
\par\textit{MoE routes each token to a subset of expert networks, increasing parameters without proportional per-token compute. Tradeoff: Capacity scales efficiently but routing balance, communication, and memory are hard. Watch for: Expert collapse and uneven routing create stragglers and weak specialization.}
\paragraph{MCQ-0970 [hard]} The team proposes Mixture of Experts for a real workload. Which finding gives the correct decision boundary?
\begin{enumerate}[label=\Alph*.]
\item It can optimize human preferences but adds reward hacking and training complexity.
\item Capacity scales efficiently but routing balance, communication, and memory are hard.
\item Every token can interact directly, while standard sequence cost grows quadratically.
\item It is stable for model comparison on the same tokenization and data, not across arbitrary setups.
\end{enumerate}
\noindent\textbf{Answer.} Capacity scales efficiently but routing balance, communication, and memory are hard.
\par\textit{MoE routes each token to a subset of expert networks, increasing parameters without proportional per-token compute. Tradeoff: Capacity scales efficiently but routing balance, communication, and memory are hard. Watch for: Expert collapse and uneven routing create stragglers and weak specialization.}
\paragraph{MCQ-0971 [hard]} Before releasing Mixture of Experts, which failure deserves a dedicated regression test?
\begin{enumerate}[label=\Alph*.]
\item Comparing perplexity across vocabularies or tokenizers is misleading.
\item A biased reward model can produce confidently optimized failure modes.
\item Omitting the square-root head-dimension scale saturates softmax gradients.
\item Expert collapse and uneven routing create stragglers and weak specialization.
\end{enumerate}
\noindent\textbf{Answer.} Expert collapse and uneven routing create stragglers and weak specialization.
\par\textit{MoE routes each token to a subset of expert networks, increasing parameters without proportional per-token compute. Tradeoff: Capacity scales efficiently but routing balance, communication, and memory are hard. Watch for: Expert collapse and uneven routing create stragglers and weak specialization.}
\paragraph{MCQ-0972 [medium]} A design review reveals this mistake: Expert collapse and uneven routing create stragglers and weak specialization. Which mechanism should be traced first?
\begin{enumerate}[label=\Alph*.]
\item Mixture of Experts
\item RLHF and PPO
\item Self-attention math
\item Cross-entropy and perplexity
\end{enumerate}
\noindent\textbf{Answer.} Mixture of Experts
\par\textit{MoE routes each token to a subset of expert networks, increasing parameters without proportional per-token compute. Tradeoff: Capacity scales efficiently but routing balance, communication, and memory are hard. Watch for: Expert collapse and uneven routing create stragglers and weak specialization.}
\paragraph{MCQ-0973 [hard]} Which causal explanation of Mixture of Experts connects what it does to when it is worth its cost?
\begin{enumerate}[label=\Alph*.]
\item Mechanism: Attention computes scaled query-key scores, applies masking and softmax, then mixes value vectors. Consequence: Every token can interact directly, while standard sequence cost grows quadratically.
\item Mechanism: MoE routes each token to a subset of expert networks, increasing parameters without proportional per-token compute. Consequence: Capacity scales efficiently but routing balance, communication, and memory are hard.
\item Mechanism: Cross-entropy is negative log likelihood; perplexity exponentiates average token loss. Consequence: It is stable for model comparison on the same tokenization and data, not across arbitrary setups.
\item Mechanism: A preference model supplies rewards while PPO constrains policy updates relative to a reference model. Consequence: It can optimize human preferences but adds reward hacking and training complexity.
\end{enumerate}
\noindent\textbf{Answer.} Mechanism: MoE routes each token to a subset of expert networks, increasing parameters without proportional per-token compute. Consequence: Capacity scales efficiently but routing balance, communication, and memory are hard.
\par\textit{MoE routes each token to a subset of expert networks, increasing parameters without proportional per-token compute. Tradeoff: Capacity scales efficiently but routing balance, communication, and memory are hard. Watch for: Expert collapse and uneven routing create stragglers and weak specialization.}
\paragraph{FC-0974 [medium]} Define Mixture of Experts in interview-ready language.
\noindent\textbf{Answer.} MoE routes each token to a subset of expert networks, increasing parameters without proportional per-token compute.
\par\textit{MoE routes each token to a subset of expert networks, increasing parameters without proportional per-token compute.}
\paragraph{FC-0975 [hard]} State the main tradeoff and production pitfall for Mixture of Experts.
\noindent\textbf{Answer.} Tradeoff: Capacity scales efficiently but routing balance, communication, and memory are hard. Pitfall: Expert collapse and uneven routing create stragglers and weak specialization.
\par\textit{Tradeoff: Capacity scales efficiently but routing balance, communication, and memory are hard. Pitfall: Expert collapse and uneven routing create stragglers and weak specialization.}
\paragraph{LA-0976 [hard]} Explain Mixture of Experts from first principles, then show how you would evaluate and operate it in production.
\noindent\textbf{Answer.} Start with the mechanism: MoE routes each token to a subset of expert networks, increasing parameters without proportional per-token compute. Make the tradeoff explicit: Capacity scales efficiently but routing balance, communication, and memory are hard. Define workload-specific offline metrics and a latency/cost budget, test important slices, instrument the critical path, and create a rollback criterion. The production review must explicitly guard against this failure: Expert collapse and uneven routing create stragglers and weak specialization.
\par\textit{A strong answer connects mechanism, measurable tradeoffs, failure isolation, observability, and rollback rather than listing product features.}
\paragraph{MCQ-0977 [medium]} A design requires this behavior: Next-token training factorizes sequence likelihood from left to right under a causal mask. Which mechanism should the team study?
\begin{enumerate}[label=\Alph*.]
\item QLoRA
\item SFT
\item Causal language modeling
\item Residual connections
\end{enumerate}
\noindent\textbf{Answer.} Causal language modeling
\par\textit{Next-token training factorizes sequence likelihood from left to right under a causal mask. Tradeoff: It provides a simple scalable objective but does not directly optimize instruction usefulness or truth. Watch for: Interpreting low perplexity as safe task performance is invalid.}
\paragraph{MCQ-0978 [hard]} The team proposes Causal language modeling for a real workload. Which finding gives the correct decision boundary?
\begin{enumerate}[label=\Alph*.]
\item It is direct and stable but can overfit style and inherit annotation errors.
\item It provides a simple scalable objective but does not directly optimize instruction usefulness or truth.
\item Deep optimization becomes feasible, while residual scale can grow and require controls.
\item Fine-tuning large models becomes cheaper, while kernels and quantization settings affect stability.
\end{enumerate}
\noindent\textbf{Answer.} It provides a simple scalable objective but does not directly optimize instruction usefulness or truth.
\par\textit{Next-token training factorizes sequence likelihood from left to right under a causal mask. Tradeoff: It provides a simple scalable objective but does not directly optimize instruction usefulness or truth. Watch for: Interpreting low perplexity as safe task performance is invalid.}
\paragraph{MCQ-0979 [hard]} Before releasing Causal language modeling, which failure deserves a dedicated regression test?
\begin{enumerate}[label=\Alph*.]
\item Interpreting low perplexity as safe task performance is invalid.
\item In-place mutation of residual tensors can break autograd or checkpointing.
\item Training on conflicting instructions without provenance makes behavior unstable.
\item Confusing training quantization with final serving quantization produces wrong expectations.
\end{enumerate}
\noindent\textbf{Answer.} Interpreting low perplexity as safe task performance is invalid.
\par\textit{Next-token training factorizes sequence likelihood from left to right under a causal mask. Tradeoff: It provides a simple scalable objective but does not directly optimize instruction usefulness or truth. Watch for: Interpreting low perplexity as safe task performance is invalid.}
\paragraph{MCQ-0980 [medium]} A design review reveals this mistake: Interpreting low perplexity as safe task performance is invalid. Which mechanism should be traced first?
\begin{enumerate}[label=\Alph*.]
\item Residual connections
\item QLoRA
\item SFT
\item Causal language modeling
\end{enumerate}
\noindent\textbf{Answer.} Causal language modeling
\par\textit{Next-token training factorizes sequence likelihood from left to right under a causal mask. Tradeoff: It provides a simple scalable objective but does not directly optimize instruction usefulness or truth. Watch for: Interpreting low perplexity as safe task performance is invalid.}
\paragraph{MCQ-0981 [hard]} Which causal explanation of Causal language modeling connects what it does to when it is worth its cost?
\begin{enumerate}[label=\Alph*.]
\item Mechanism: QLoRA backpropagates through a frozen quantized base model into LoRA adapters using memory-saving quantization techniques. Consequence: Fine-tuning large models becomes cheaper, while kernels and quantization settings affect stability.
\item Mechanism: Next-token training factorizes sequence likelihood from left to right under a causal mask. Consequence: It provides a simple scalable objective but does not directly optimize instruction usefulness or truth.
\item Mechanism: Residual paths add transformed activations back to the stream to improve gradient flow and feature reuse. Consequence: Deep optimization becomes feasible, while residual scale can grow and require controls.
\item Mechanism: Supervised fine-tuning teaches instruction-response behavior from curated demonstrations. Consequence: It is direct and stable but can overfit style and inherit annotation errors.
\end{enumerate}
\noindent\textbf{Answer.} Mechanism: Next-token training factorizes sequence likelihood from left to right under a causal mask. Consequence: It provides a simple scalable objective but does not directly optimize instruction usefulness or truth.
\par\textit{Next-token training factorizes sequence likelihood from left to right under a causal mask. Tradeoff: It provides a simple scalable objective but does not directly optimize instruction usefulness or truth. Watch for: Interpreting low perplexity as safe task performance is invalid.}
\paragraph{FC-0982 [medium]} Define Causal language modeling in interview-ready language.
\noindent\textbf{Answer.} Next-token training factorizes sequence likelihood from left to right under a causal mask.
\par\textit{Next-token training factorizes sequence likelihood from left to right under a causal mask.}
\paragraph{FC-0983 [hard]} State the main tradeoff and production pitfall for Causal language modeling.
\noindent\textbf{Answer.} Tradeoff: It provides a simple scalable objective but does not directly optimize instruction usefulness or truth. Pitfall: Interpreting low perplexity as safe task performance is invalid.
\par\textit{Tradeoff: It provides a simple scalable objective but does not directly optimize instruction usefulness or truth. Pitfall: Interpreting low perplexity as safe task performance is invalid.}
\paragraph{LA-0984 [hard]} Explain Causal language modeling from first principles, then show how you would evaluate and operate it in production.
\noindent\textbf{Answer.} Start with the mechanism: Next-token training factorizes sequence likelihood from left to right under a causal mask. Make the tradeoff explicit: It provides a simple scalable objective but does not directly optimize instruction usefulness or truth. Define workload-specific offline metrics and a latency/cost budget, test important slices, instrument the critical path, and create a rollback criterion. The production review must explicitly guard against this failure: Interpreting low perplexity as safe task performance is invalid.
\par\textit{A strong answer connects mechanism, measurable tradeoffs, failure isolation, observability, and rollback rather than listing product features.}
\paragraph{MCQ-0985 [medium]} A design requires this behavior: Cross-entropy is negative log likelihood; perplexity exponentiates average token loss. Which mechanism should the team study?
\begin{enumerate}[label=\Alph*.]
\item Cross-entropy and perplexity
\item Positional encoding
\item Knowledge distillation
\item Layer normalization and RMSNorm
\end{enumerate}
\noindent\textbf{Answer.} Cross-entropy and perplexity
\par\textit{Cross-entropy is negative log likelihood; perplexity exponentiates average token loss. Tradeoff: It is stable for model comparison on the same tokenization and data, not across arbitrary setups. Watch for: Comparing perplexity across vocabularies or tokenizers is misleading.}
\paragraph{MCQ-0986 [hard]} The team proposes Cross-entropy and perplexity for a real workload. Which finding gives the correct decision boundary?
\begin{enumerate}[label=\Alph*.]
\item It is stable for model comparison on the same tokenization and data, not across arbitrary setups.
\item Simpler normalization can be faster but architecture placement changes optimization dynamics.
\item Relative methods extrapolate differently, while implementation and scaling choices affect long context.
\item Serving cost falls but rare capabilities and calibration may be lost.
\end{enumerate}
\noindent\textbf{Answer.} It is stable for model comparison on the same tokenization and data, not across arbitrary setups.
\par\textit{Cross-entropy is negative log likelihood; perplexity exponentiates average token loss. Tradeoff: It is stable for model comparison on the same tokenization and data, not across arbitrary setups. Watch for: Comparing perplexity across vocabularies or tokenizers is misleading.}
\paragraph{MCQ-0987 [hard]} Before releasing Cross-entropy and perplexity, which failure deserves a dedicated regression test?
\begin{enumerate}[label=\Alph*.]
\item Distilling only easy teacher outputs creates a deceptively strong narrow model.
\item Extending context without validating the position scheme can collapse quality.
\item Comparing perplexity across vocabularies or tokenizers is misleading.
\item Moving pre-norm to post-norm without retuning can destabilize deep training.
\end{enumerate}
\noindent\textbf{Answer.} Comparing perplexity across vocabularies or tokenizers is misleading.
\par\textit{Cross-entropy is negative log likelihood; perplexity exponentiates average token loss. Tradeoff: It is stable for model comparison on the same tokenization and data, not across arbitrary setups. Watch for: Comparing perplexity across vocabularies or tokenizers is misleading.}
\paragraph{MCQ-0988 [medium]} A design review reveals this mistake: Comparing perplexity across vocabularies or tokenizers is misleading. Which mechanism should be traced first?
\begin{enumerate}[label=\Alph*.]
\item Cross-entropy and perplexity
\item Knowledge distillation
\item Positional encoding
\item Layer normalization and RMSNorm
\end{enumerate}
\noindent\textbf{Answer.} Cross-entropy and perplexity
\par\textit{Cross-entropy is negative log likelihood; perplexity exponentiates average token loss. Tradeoff: It is stable for model comparison on the same tokenization and data, not across arbitrary setups. Watch for: Comparing perplexity across vocabularies or tokenizers is misleading.}
\paragraph{MCQ-0989 [hard]} Which causal explanation of Cross-entropy and perplexity connects what it does to when it is worth its cost?
\begin{enumerate}[label=\Alph*.]
\item Mechanism: A student learns from teacher probabilities, generated data, rationales, or intermediate representations. Consequence: Serving cost falls but rare capabilities and calibration may be lost.
\item Mechanism: Normalization stabilizes activations; RMSNorm removes mean centering and normalizes by root-mean-square magnitude. Consequence: Simpler normalization can be faster but architecture placement changes optimization dynamics.
\item Mechanism: Position signals break permutation invariance using absolute, relative, rotary, or bias-based schemes. Consequence: Relative methods extrapolate differently, while implementation and scaling choices affect long context.
\item Mechanism: Cross-entropy is negative log likelihood; perplexity exponentiates average token loss. Consequence: It is stable for model comparison on the same tokenization and data, not across arbitrary setups.
\end{enumerate}
\noindent\textbf{Answer.} Mechanism: Cross-entropy is negative log likelihood; perplexity exponentiates average token loss. Consequence: It is stable for model comparison on the same tokenization and data, not across arbitrary setups.
\par\textit{Cross-entropy is negative log likelihood; perplexity exponentiates average token loss. Tradeoff: It is stable for model comparison on the same tokenization and data, not across arbitrary setups. Watch for: Comparing perplexity across vocabularies or tokenizers is misleading.}
\paragraph{FC-0990 [medium]} Define Cross-entropy and perplexity in interview-ready language.
\noindent\textbf{Answer.} Cross-entropy is negative log likelihood; perplexity exponentiates average token loss.
\par\textit{Cross-entropy is negative log likelihood; perplexity exponentiates average token loss.}
\paragraph{FC-0991 [hard]} State the main tradeoff and production pitfall for Cross-entropy and perplexity.
\noindent\textbf{Answer.} Tradeoff: It is stable for model comparison on the same tokenization and data, not across arbitrary setups. Pitfall: Comparing perplexity across vocabularies or tokenizers is misleading.
\par\textit{Tradeoff: It is stable for model comparison on the same tokenization and data, not across arbitrary setups. Pitfall: Comparing perplexity across vocabularies or tokenizers is misleading.}
\paragraph{LA-0992 [hard]} Explain Cross-entropy and perplexity from first principles, then show how you would evaluate and operate it in production.
\noindent\textbf{Answer.} Start with the mechanism: Cross-entropy is negative log likelihood; perplexity exponentiates average token loss. Make the tradeoff explicit: It is stable for model comparison on the same tokenization and data, not across arbitrary setups. Define workload-specific offline metrics and a latency/cost budget, test important slices, instrument the critical path, and create a rollback criterion. The production review must explicitly guard against this failure: Comparing perplexity across vocabularies or tokenizers is misleading.
\par\textit{A strong answer connects mechanism, measurable tradeoffs, failure isolation, observability, and rollback rather than listing product features.}
\paragraph{MCQ-0993 [medium]} A design requires this behavior: Supervised fine-tuning teaches instruction-response behavior from curated demonstrations. Which mechanism should the team study?
\begin{enumerate}[label=\Alph*.]
\item Feed-forward networks
\item Mixture of Experts
\item SFT
\item Embedding layer
\end{enumerate}
\noindent\textbf{Answer.} SFT
\par\textit{Supervised fine-tuning teaches instruction-response behavior from curated demonstrations. Tradeoff: It is direct and stable but can overfit style and inherit annotation errors. Watch for: Training on conflicting instructions without provenance makes behavior unstable.}
\paragraph{MCQ-0994 [hard]} The team proposes SFT for a real workload. Which finding gives the correct decision boundary?
\begin{enumerate}[label=\Alph*.]
\item Weight tying saves parameters but constrains input-output geometry.
\item They hold substantial parameters and compute while offering parallel token processing.
\item It is direct and stable but can overfit style and inherit annotation errors.
\item Capacity scales efficiently but routing balance, communication, and memory are hard.
\end{enumerate}
\noindent\textbf{Answer.} It is direct and stable but can overfit style and inherit annotation errors.
\par\textit{Supervised fine-tuning teaches instruction-response behavior from curated demonstrations. Tradeoff: It is direct and stable but can overfit style and inherit annotation errors. Watch for: Training on conflicting instructions without provenance makes behavior unstable.}
\paragraph{MCQ-0995 [hard]} Before releasing SFT, which failure deserves a dedicated regression test?
\begin{enumerate}[label=\Alph*.]
\item Ignoring activation memory underestimates training capacity.
\item Forgetting vocabulary changes invalidates checkpoint embedding shapes.
\item Training on conflicting instructions without provenance makes behavior unstable.
\item Expert collapse and uneven routing create stragglers and weak specialization.
\end{enumerate}
\noindent\textbf{Answer.} Training on conflicting instructions without provenance makes behavior unstable.
\par\textit{Supervised fine-tuning teaches instruction-response behavior from curated demonstrations. Tradeoff: It is direct and stable but can overfit style and inherit annotation errors. Watch for: Training on conflicting instructions without provenance makes behavior unstable.}
\paragraph{MCQ-0996 [medium]} A design review reveals this mistake: Training on conflicting instructions without provenance makes behavior unstable. Which mechanism should be traced first?
\begin{enumerate}[label=\Alph*.]
\item Feed-forward networks
\item Embedding layer
\item Mixture of Experts
\item SFT
\end{enumerate}
\noindent\textbf{Answer.} SFT
\par\textit{Supervised fine-tuning teaches instruction-response behavior from curated demonstrations. Tradeoff: It is direct and stable but can overfit style and inherit annotation errors. Watch for: Training on conflicting instructions without provenance makes behavior unstable.}
\paragraph{MCQ-0997 [hard]} Which causal explanation of SFT connects what it does to when it is worth its cost?
\begin{enumerate}[label=\Alph*.]
\item Mechanism: Transformer MLP blocks expand hidden dimensions, apply nonlinear gates, and project back per token. Consequence: They hold substantial parameters and compute while offering parallel token processing.
\item Mechanism: Supervised fine-tuning teaches instruction-response behavior from curated demonstrations. Consequence: It is direct and stable but can overfit style and inherit annotation errors.
\item Mechanism: MoE routes each token to a subset of expert networks, increasing parameters without proportional per-token compute. Consequence: Capacity scales efficiently but routing balance, communication, and memory are hard.
\item Mechanism: Learned token vectors map discrete IDs into continuous model space and are often tied to output weights. Consequence: Weight tying saves parameters but constrains input-output geometry.
\end{enumerate}
\noindent\textbf{Answer.} Mechanism: Supervised fine-tuning teaches instruction-response behavior from curated demonstrations. Consequence: It is direct and stable but can overfit style and inherit annotation errors.
\par\textit{Supervised fine-tuning teaches instruction-response behavior from curated demonstrations. Tradeoff: It is direct and stable but can overfit style and inherit annotation errors. Watch for: Training on conflicting instructions without provenance makes behavior unstable.}
\paragraph{FC-0998 [medium]} Define SFT in interview-ready language.
\noindent\textbf{Answer.} Supervised fine-tuning teaches instruction-response behavior from curated demonstrations.
\par\textit{Supervised fine-tuning teaches instruction-response behavior from curated demonstrations.}
\paragraph{FC-0999 [hard]} State the main tradeoff and production pitfall for SFT.
\noindent\textbf{Answer.} Tradeoff: It is direct and stable but can overfit style and inherit annotation errors. Pitfall: Training on conflicting instructions without provenance makes behavior unstable.
\par\textit{Tradeoff: It is direct and stable but can overfit style and inherit annotation errors. Pitfall: Training on conflicting instructions without provenance makes behavior unstable.}
\paragraph{LA-1000 [hard]} Explain SFT from first principles, then show how you would evaluate and operate it in production.
\noindent\textbf{Answer.} Start with the mechanism: Supervised fine-tuning teaches instruction-response behavior from curated demonstrations. Make the tradeoff explicit: It is direct and stable but can overfit style and inherit annotation errors. Define workload-specific offline metrics and a latency/cost budget, test important slices, instrument the critical path, and create a rollback criterion. The production review must explicitly guard against this failure: Training on conflicting instructions without provenance makes behavior unstable.
\par\textit{A strong answer connects mechanism, measurable tradeoffs, failure isolation, observability, and rollback rather than listing product features.}
\paragraph{MCQ-1001 [medium]} A design requires this behavior: A preference model supplies rewards while PPO constrains policy updates relative to a reference model. Which mechanism should the team study?
\begin{enumerate}[label=\Alph*.]
\item Cross-entropy and perplexity
\item Mixture of Experts
\item LoRA
\item RLHF and PPO
\end{enumerate}
\noindent\textbf{Answer.} RLHF and PPO
\par\textit{A preference model supplies rewards while PPO constrains policy updates relative to a reference model. Tradeoff: It can optimize human preferences but adds reward hacking and training complexity. Watch for: A biased reward model can produce confidently optimized failure modes.}
\paragraph{MCQ-1002 [hard]} The team proposes RLHF and PPO for a real workload. Which finding gives the correct decision boundary?
\begin{enumerate}[label=\Alph*.]
\item Memory and storage fall, while rank and target-module choices limit expressiveness.
\item It is stable for model comparison on the same tokenization and data, not across arbitrary setups.
\item Capacity scales efficiently but routing balance, communication, and memory are hard.
\item It can optimize human preferences but adds reward hacking and training complexity.
\end{enumerate}
\noindent\textbf{Answer.} It can optimize human preferences but adds reward hacking and training complexity.
\par\textit{A preference model supplies rewards while PPO constrains policy updates relative to a reference model. Tradeoff: It can optimize human preferences but adds reward hacking and training complexity. Watch for: A biased reward model can produce confidently optimized failure modes.}
\paragraph{MCQ-1003 [hard]} Before releasing RLHF and PPO, which failure deserves a dedicated regression test?
\begin{enumerate}[label=\Alph*.]
\item Expert collapse and uneven routing create stragglers and weak specialization.
\item Merging incompatible adapters without evaluation creates interference.
\item Comparing perplexity across vocabularies or tokenizers is misleading.
\item A biased reward model can produce confidently optimized failure modes.
\end{enumerate}
\noindent\textbf{Answer.} A biased reward model can produce confidently optimized failure modes.
\par\textit{A preference model supplies rewards while PPO constrains policy updates relative to a reference model. Tradeoff: It can optimize human preferences but adds reward hacking and training complexity. Watch for: A biased reward model can produce confidently optimized failure modes.}
\paragraph{MCQ-1004 [medium]} A design review reveals this mistake: A biased reward model can produce confidently optimized failure modes. Which mechanism should be traced first?
\begin{enumerate}[label=\Alph*.]
\item Cross-entropy and perplexity
\item RLHF and PPO
\item LoRA
\item Mixture of Experts
\end{enumerate}
\noindent\textbf{Answer.} RLHF and PPO
\par\textit{A preference model supplies rewards while PPO constrains policy updates relative to a reference model. Tradeoff: It can optimize human preferences but adds reward hacking and training complexity. Watch for: A biased reward model can produce confidently optimized failure modes.}
\paragraph{MCQ-1005 [hard]} Which causal explanation of RLHF and PPO connects what it does to when it is worth its cost?
\begin{enumerate}[label=\Alph*.]
\item Mechanism: MoE routes each token to a subset of expert networks, increasing parameters without proportional per-token compute. Consequence: Capacity scales efficiently but routing balance, communication, and memory are hard.
\item Mechanism: Cross-entropy is negative log likelihood; perplexity exponentiates average token loss. Consequence: It is stable for model comparison on the same tokenization and data, not across arbitrary setups.
\item Mechanism: LoRA trains low-rank update matrices while freezing base weights for parameter-efficient adaptation. Consequence: Memory and storage fall, while rank and target-module choices limit expressiveness.
\item Mechanism: A preference model supplies rewards while PPO constrains policy updates relative to a reference model. Consequence: It can optimize human preferences but adds reward hacking and training complexity.
\end{enumerate}
\noindent\textbf{Answer.} Mechanism: A preference model supplies rewards while PPO constrains policy updates relative to a reference model. Consequence: It can optimize human preferences but adds reward hacking and training complexity.
\par\textit{A preference model supplies rewards while PPO constrains policy updates relative to a reference model. Tradeoff: It can optimize human preferences but adds reward hacking and training complexity. Watch for: A biased reward model can produce confidently optimized failure modes.}
\paragraph{FC-1006 [medium]} Define RLHF and PPO in interview-ready language.
\noindent\textbf{Answer.} A preference model supplies rewards while PPO constrains policy updates relative to a reference model.
\par\textit{A preference model supplies rewards while PPO constrains policy updates relative to a reference model.}
\paragraph{FC-1007 [hard]} State the main tradeoff and production pitfall for RLHF and PPO.
\noindent\textbf{Answer.} Tradeoff: It can optimize human preferences but adds reward hacking and training complexity. Pitfall: A biased reward model can produce confidently optimized failure modes.
\par\textit{Tradeoff: It can optimize human preferences but adds reward hacking and training complexity. Pitfall: A biased reward model can produce confidently optimized failure modes.}
\paragraph{LA-1008 [hard]} Explain RLHF and PPO from first principles, then show how you would evaluate and operate it in production.
\noindent\textbf{Answer.} Start with the mechanism: A preference model supplies rewards while PPO constrains policy updates relative to a reference model. Make the tradeoff explicit: It can optimize human preferences but adds reward hacking and training complexity. Define workload-specific offline metrics and a latency/cost budget, test important slices, instrument the critical path, and create a rollback criterion. The production review must explicitly guard against this failure: A biased reward model can produce confidently optimized failure modes.
\par\textit{A strong answer connects mechanism, measurable tradeoffs, failure isolation, observability, and rollback rather than listing product features.}
\paragraph{MCQ-1009 [medium]} A design requires this behavior: Direct Preference Optimization fits preferred over rejected responses through a classification-style objective relative to a reference policy. Which mechanism should the team study?
\begin{enumerate}[label=\Alph*.]
\item Mixture of Experts
\item Feed-forward networks
\item DPO
\item Embedding layer
\end{enumerate}
\noindent\textbf{Answer.} DPO
\par\textit{Direct Preference Optimization fits preferred over rejected responses through a classification-style objective relative to a reference policy. Tradeoff: It simplifies preference training but still depends on pair quality and distribution coverage. Watch for: Treating implicit user choices as clean preferences introduces confounding.}
\paragraph{MCQ-1010 [hard]} The team proposes DPO for a real workload. Which finding gives the correct decision boundary?
\begin{enumerate}[label=\Alph*.]
\item It simplifies preference training but still depends on pair quality and distribution coverage.
\item They hold substantial parameters and compute while offering parallel token processing.
\item Weight tying saves parameters but constrains input-output geometry.
\item Capacity scales efficiently but routing balance, communication, and memory are hard.
\end{enumerate}
\noindent\textbf{Answer.} It simplifies preference training but still depends on pair quality and distribution coverage.
\par\textit{Direct Preference Optimization fits preferred over rejected responses through a classification-style objective relative to a reference policy. Tradeoff: It simplifies preference training but still depends on pair quality and distribution coverage. Watch for: Treating implicit user choices as clean preferences introduces confounding.}
\paragraph{MCQ-1011 [hard]} Before releasing DPO, which failure deserves a dedicated regression test?
\begin{enumerate}[label=\Alph*.]
\item Expert collapse and uneven routing create stragglers and weak specialization.
\item Treating implicit user choices as clean preferences introduces confounding.
\item Forgetting vocabulary changes invalidates checkpoint embedding shapes.
\item Ignoring activation memory underestimates training capacity.
\end{enumerate}
\noindent\textbf{Answer.} Treating implicit user choices as clean preferences introduces confounding.
\par\textit{Direct Preference Optimization fits preferred over rejected responses through a classification-style objective relative to a reference policy. Tradeoff: It simplifies preference training but still depends on pair quality and distribution coverage. Watch for: Treating implicit user choices as clean preferences introduces confounding.}
\paragraph{MCQ-1012 [medium]} A design review reveals this mistake: Treating implicit user choices as clean preferences introduces confounding. Which mechanism should be traced first?
\begin{enumerate}[label=\Alph*.]
\item DPO
\item Embedding layer
\item Mixture of Experts
\item Feed-forward networks
\end{enumerate}
\noindent\textbf{Answer.} DPO
\par\textit{Direct Preference Optimization fits preferred over rejected responses through a classification-style objective relative to a reference policy. Tradeoff: It simplifies preference training but still depends on pair quality and distribution coverage. Watch for: Treating implicit user choices as clean preferences introduces confounding.}
\paragraph{MCQ-1013 [hard]} Which causal explanation of DPO connects what it does to when it is worth its cost?
\begin{enumerate}[label=\Alph*.]
\item Mechanism: MoE routes each token to a subset of expert networks, increasing parameters without proportional per-token compute. Consequence: Capacity scales efficiently but routing balance, communication, and memory are hard.
\item Mechanism: Transformer MLP blocks expand hidden dimensions, apply nonlinear gates, and project back per token. Consequence: They hold substantial parameters and compute while offering parallel token processing.
\item Mechanism: Learned token vectors map discrete IDs into continuous model space and are often tied to output weights. Consequence: Weight tying saves parameters but constrains input-output geometry.
\item Mechanism: Direct Preference Optimization fits preferred over rejected responses through a classification-style objective relative to a reference policy. Consequence: It simplifies preference training but still depends on pair quality and distribution coverage.
\end{enumerate}
\noindent\textbf{Answer.} Mechanism: Direct Preference Optimization fits preferred over rejected responses through a classification-style objective relative to a reference policy. Consequence: It simplifies preference training but still depends on pair quality and distribution coverage.
\par\textit{Direct Preference Optimization fits preferred over rejected responses through a classification-style objective relative to a reference policy. Tradeoff: It simplifies preference training but still depends on pair quality and distribution coverage. Watch for: Treating implicit user choices as clean preferences introduces confounding.}
\paragraph{FC-1014 [medium]} Define DPO in interview-ready language.
\noindent\textbf{Answer.} Direct Preference Optimization fits preferred over rejected responses through a classification-style objective relative to a reference policy.
\par\textit{Direct Preference Optimization fits preferred over rejected responses through a classification-style objective relative to a reference policy.}
\paragraph{FC-1015 [hard]} State the main tradeoff and production pitfall for DPO.
\noindent\textbf{Answer.} Tradeoff: It simplifies preference training but still depends on pair quality and distribution coverage. Pitfall: Treating implicit user choices as clean preferences introduces confounding.
\par\textit{Tradeoff: It simplifies preference training but still depends on pair quality and distribution coverage. Pitfall: Treating implicit user choices as clean preferences introduces confounding.}
\paragraph{LA-1016 [hard]} Explain DPO from first principles, then show how you would evaluate and operate it in production.
\noindent\textbf{Answer.} Start with the mechanism: Direct Preference Optimization fits preferred over rejected responses through a classification-style objective relative to a reference policy. Make the tradeoff explicit: It simplifies preference training but still depends on pair quality and distribution coverage. Define workload-specific offline metrics and a latency/cost budget, test important slices, instrument the critical path, and create a rollback criterion. The production review must explicitly guard against this failure: Treating implicit user choices as clean preferences introduces confounding.
\par\textit{A strong answer connects mechanism, measurable tradeoffs, failure isolation, observability, and rollback rather than listing product features.}
\paragraph{MCQ-1017 [medium]} A design requires this behavior: LoRA trains low-rank update matrices while freezing base weights for parameter-efficient adaptation. Which mechanism should the team study?
\begin{enumerate}[label=\Alph*.]
\item Cross-entropy and perplexity
\item Positional encoding
\item LoRA
\item Knowledge distillation
\end{enumerate}
\noindent\textbf{Answer.} LoRA
\par\textit{LoRA trains low-rank update matrices while freezing base weights for parameter-efficient adaptation. Tradeoff: Memory and storage fall, while rank and target-module choices limit expressiveness. Watch for: Merging incompatible adapters without evaluation creates interference.}
\paragraph{MCQ-1018 [hard]} The team proposes LoRA for a real workload. Which finding gives the correct decision boundary?
\begin{enumerate}[label=\Alph*.]
\item Memory and storage fall, while rank and target-module choices limit expressiveness.
\item It is stable for model comparison on the same tokenization and data, not across arbitrary setups.
\item Serving cost falls but rare capabilities and calibration may be lost.
\item Relative methods extrapolate differently, while implementation and scaling choices affect long context.
\end{enumerate}
\noindent\textbf{Answer.} Memory and storage fall, while rank and target-module choices limit expressiveness.
\par\textit{LoRA trains low-rank update matrices while freezing base weights for parameter-efficient adaptation. Tradeoff: Memory and storage fall, while rank and target-module choices limit expressiveness. Watch for: Merging incompatible adapters without evaluation creates interference.}
\paragraph{MCQ-1019 [hard]} Before releasing LoRA, which failure deserves a dedicated regression test?
\begin{enumerate}[label=\Alph*.]
\item Distilling only easy teacher outputs creates a deceptively strong narrow model.
\item Extending context without validating the position scheme can collapse quality.
\item Merging incompatible adapters without evaluation creates interference.
\item Comparing perplexity across vocabularies or tokenizers is misleading.
\end{enumerate}
\noindent\textbf{Answer.} Merging incompatible adapters without evaluation creates interference.
\par\textit{LoRA trains low-rank update matrices while freezing base weights for parameter-efficient adaptation. Tradeoff: Memory and storage fall, while rank and target-module choices limit expressiveness. Watch for: Merging incompatible adapters without evaluation creates interference.}
\paragraph{MCQ-1020 [medium]} A design review reveals this mistake: Merging incompatible adapters without evaluation creates interference. Which mechanism should be traced first?
\begin{enumerate}[label=\Alph*.]
\item Positional encoding
\item LoRA
\item Cross-entropy and perplexity
\item Knowledge distillation
\end{enumerate}
\noindent\textbf{Answer.} LoRA
\par\textit{LoRA trains low-rank update matrices while freezing base weights for parameter-efficient adaptation. Tradeoff: Memory and storage fall, while rank and target-module choices limit expressiveness. Watch for: Merging incompatible adapters without evaluation creates interference.}
\paragraph{MCQ-1021 [hard]} Which causal explanation of LoRA connects what it does to when it is worth its cost?
\begin{enumerate}[label=\Alph*.]
\item Mechanism: Cross-entropy is negative log likelihood; perplexity exponentiates average token loss. Consequence: It is stable for model comparison on the same tokenization and data, not across arbitrary setups.
\item Mechanism: A student learns from teacher probabilities, generated data, rationales, or intermediate representations. Consequence: Serving cost falls but rare capabilities and calibration may be lost.
\item Mechanism: LoRA trains low-rank update matrices while freezing base weights for parameter-efficient adaptation. Consequence: Memory and storage fall, while rank and target-module choices limit expressiveness.
\item Mechanism: Position signals break permutation invariance using absolute, relative, rotary, or bias-based schemes. Consequence: Relative methods extrapolate differently, while implementation and scaling choices affect long context.
\end{enumerate}
\noindent\textbf{Answer.} Mechanism: LoRA trains low-rank update matrices while freezing base weights for parameter-efficient adaptation. Consequence: Memory and storage fall, while rank and target-module choices limit expressiveness.
\par\textit{LoRA trains low-rank update matrices while freezing base weights for parameter-efficient adaptation. Tradeoff: Memory and storage fall, while rank and target-module choices limit expressiveness. Watch for: Merging incompatible adapters without evaluation creates interference.}
\paragraph{FC-1022 [medium]} Define LoRA in interview-ready language.
\noindent\textbf{Answer.} LoRA trains low-rank update matrices while freezing base weights for parameter-efficient adaptation.
\par\textit{LoRA trains low-rank update matrices while freezing base weights for parameter-efficient adaptation.}
\paragraph{FC-1023 [hard]} State the main tradeoff and production pitfall for LoRA.
\noindent\textbf{Answer.} Tradeoff: Memory and storage fall, while rank and target-module choices limit expressiveness. Pitfall: Merging incompatible adapters without evaluation creates interference.
\par\textit{Tradeoff: Memory and storage fall, while rank and target-module choices limit expressiveness. Pitfall: Merging incompatible adapters without evaluation creates interference.}
\paragraph{LA-1024 [hard]} Explain LoRA from first principles, then show how you would evaluate and operate it in production.
\noindent\textbf{Answer.} Start with the mechanism: LoRA trains low-rank update matrices while freezing base weights for parameter-efficient adaptation. Make the tradeoff explicit: Memory and storage fall, while rank and target-module choices limit expressiveness. Define workload-specific offline metrics and a latency/cost budget, test important slices, instrument the critical path, and create a rollback criterion. The production review must explicitly guard against this failure: Merging incompatible adapters without evaluation creates interference.
\par\textit{A strong answer connects mechanism, measurable tradeoffs, failure isolation, observability, and rollback rather than listing product features.}
\paragraph{MCQ-1025 [medium]} A design requires this behavior: QLoRA backpropagates through a frozen quantized base model into LoRA adapters using memory-saving quantization techniques. Which mechanism should the team study?
\begin{enumerate}[label=\Alph*.]
\item Embedding layer
\item RLHF and PPO
\item Mixture of Experts
\item QLoRA
\end{enumerate}
\noindent\textbf{Answer.} QLoRA
\par\textit{QLoRA backpropagates through a frozen quantized base model into LoRA adapters using memory-saving quantization techniques. Tradeoff: Fine-tuning large models becomes cheaper, while kernels and quantization settings affect stability. Watch for: Confusing training quantization with final serving quantization produces wrong expectations.}
\paragraph{MCQ-1026 [hard]} The team proposes QLoRA for a real workload. Which finding gives the correct decision boundary?
\begin{enumerate}[label=\Alph*.]
\item Weight tying saves parameters but constrains input-output geometry.
\item It can optimize human preferences but adds reward hacking and training complexity.
\item Capacity scales efficiently but routing balance, communication, and memory are hard.
\item Fine-tuning large models becomes cheaper, while kernels and quantization settings affect stability.
\end{enumerate}
\noindent\textbf{Answer.} Fine-tuning large models becomes cheaper, while kernels and quantization settings affect stability.
\par\textit{QLoRA backpropagates through a frozen quantized base model into LoRA adapters using memory-saving quantization techniques. Tradeoff: Fine-tuning large models becomes cheaper, while kernels and quantization settings affect stability. Watch for: Confusing training quantization with final serving quantization produces wrong expectations.}
\paragraph{MCQ-1027 [hard]} Before releasing QLoRA, which failure deserves a dedicated regression test?
\begin{enumerate}[label=\Alph*.]
\item Forgetting vocabulary changes invalidates checkpoint embedding shapes.
\item Confusing training quantization with final serving quantization produces wrong expectations.
\item A biased reward model can produce confidently optimized failure modes.
\item Expert collapse and uneven routing create stragglers and weak specialization.
\end{enumerate}
\noindent\textbf{Answer.} Confusing training quantization with final serving quantization produces wrong expectations.
\par\textit{QLoRA backpropagates through a frozen quantized base model into LoRA adapters using memory-saving quantization techniques. Tradeoff: Fine-tuning large models becomes cheaper, while kernels and quantization settings affect stability. Watch for: Confusing training quantization with final serving quantization produces wrong expectations.}
\paragraph{MCQ-1028 [medium]} A design review reveals this mistake: Confusing training quantization with final serving quantization produces wrong expectations. Which mechanism should be traced first?
\begin{enumerate}[label=\Alph*.]
\item Embedding layer
\item QLoRA
\item RLHF and PPO
\item Mixture of Experts
\end{enumerate}
\noindent\textbf{Answer.} QLoRA
\par\textit{QLoRA backpropagates through a frozen quantized base model into LoRA adapters using memory-saving quantization techniques. Tradeoff: Fine-tuning large models becomes cheaper, while kernels and quantization settings affect stability. Watch for: Confusing training quantization with final serving quantization produces wrong expectations.}
\paragraph{MCQ-1029 [hard]} Which causal explanation of QLoRA connects what it does to when it is worth its cost?
\begin{enumerate}[label=\Alph*.]
\item Mechanism: A preference model supplies rewards while PPO constrains policy updates relative to a reference model. Consequence: It can optimize human preferences but adds reward hacking and training complexity.
\item Mechanism: QLoRA backpropagates through a frozen quantized base model into LoRA adapters using memory-saving quantization techniques. Consequence: Fine-tuning large models becomes cheaper, while kernels and quantization settings affect stability.
\item Mechanism: MoE routes each token to a subset of expert networks, increasing parameters without proportional per-token compute. Consequence: Capacity scales efficiently but routing balance, communication, and memory are hard.
\item Mechanism: Learned token vectors map discrete IDs into continuous model space and are often tied to output weights. Consequence: Weight tying saves parameters but constrains input-output geometry.
\end{enumerate}
\noindent\textbf{Answer.} Mechanism: QLoRA backpropagates through a frozen quantized base model into LoRA adapters using memory-saving quantization techniques. Consequence: Fine-tuning large models becomes cheaper, while kernels and quantization settings affect stability.
\par\textit{QLoRA backpropagates through a frozen quantized base model into LoRA adapters using memory-saving quantization techniques. Tradeoff: Fine-tuning large models becomes cheaper, while kernels and quantization settings affect stability. Watch for: Confusing training quantization with final serving quantization produces wrong expectations.}
\paragraph{FC-1030 [medium]} Define QLoRA in interview-ready language.
\noindent\textbf{Answer.} QLoRA backpropagates through a frozen quantized base model into LoRA adapters using memory-saving quantization techniques.
\par\textit{QLoRA backpropagates through a frozen quantized base model into LoRA adapters using memory-saving quantization techniques.}
\paragraph{FC-1031 [hard]} State the main tradeoff and production pitfall for QLoRA.
\noindent\textbf{Answer.} Tradeoff: Fine-tuning large models becomes cheaper, while kernels and quantization settings affect stability. Pitfall: Confusing training quantization with final serving quantization produces wrong expectations.
\par\textit{Tradeoff: Fine-tuning large models becomes cheaper, while kernels and quantization settings affect stability. Pitfall: Confusing training quantization with final serving quantization produces wrong expectations.}
\paragraph{LA-1032 [hard]} Explain QLoRA from first principles, then show how you would evaluate and operate it in production.
\noindent\textbf{Answer.} Start with the mechanism: QLoRA backpropagates through a frozen quantized base model into LoRA adapters using memory-saving quantization techniques. Make the tradeoff explicit: Fine-tuning large models becomes cheaper, while kernels and quantization settings affect stability. Define workload-specific offline metrics and a latency/cost budget, test important slices, instrument the critical path, and create a rollback criterion. The production review must explicitly guard against this failure: Confusing training quantization with final serving quantization produces wrong expectations.
\par\textit{A strong answer connects mechanism, measurable tradeoffs, failure isolation, observability, and rollback rather than listing product features.}
\paragraph{MCQ-1033 [medium]} A design requires this behavior: A student learns from teacher probabilities, generated data, rationales, or intermediate representations. Which mechanism should the team study?
\begin{enumerate}[label=\Alph*.]
\item Self-attention math
\item Feed-forward networks
\item Knowledge distillation
\item Embedding layer
\end{enumerate}
\noindent\textbf{Answer.} Knowledge distillation
\par\textit{A student learns from teacher probabilities, generated data, rationales, or intermediate representations. Tradeoff: Serving cost falls but rare capabilities and calibration may be lost. Watch for: Distilling only easy teacher outputs creates a deceptively strong narrow model.}
\paragraph{MCQ-1034 [hard]} The team proposes Knowledge distillation for a real workload. Which finding gives the correct decision boundary?
\begin{enumerate}[label=\Alph*.]
\item Every token can interact directly, while standard sequence cost grows quadratically.
\item They hold substantial parameters and compute while offering parallel token processing.
\item Serving cost falls but rare capabilities and calibration may be lost.
\item Weight tying saves parameters but constrains input-output geometry.
\end{enumerate}
\noindent\textbf{Answer.} Serving cost falls but rare capabilities and calibration may be lost.
\par\textit{A student learns from teacher probabilities, generated data, rationales, or intermediate representations. Tradeoff: Serving cost falls but rare capabilities and calibration may be lost. Watch for: Distilling only easy teacher outputs creates a deceptively strong narrow model.}
\paragraph{MCQ-1035 [hard]} Before releasing Knowledge distillation, which failure deserves a dedicated regression test?
\begin{enumerate}[label=\Alph*.]
\item Omitting the square-root head-dimension scale saturates softmax gradients.
\item Distilling only easy teacher outputs creates a deceptively strong narrow model.
\item Forgetting vocabulary changes invalidates checkpoint embedding shapes.
\item Ignoring activation memory underestimates training capacity.
\end{enumerate}
\noindent\textbf{Answer.} Distilling only easy teacher outputs creates a deceptively strong narrow model.
\par\textit{A student learns from teacher probabilities, generated data, rationales, or intermediate representations. Tradeoff: Serving cost falls but rare capabilities and calibration may be lost. Watch for: Distilling only easy teacher outputs creates a deceptively strong narrow model.}
\paragraph{MCQ-1036 [medium]} A design review reveals this mistake: Distilling only easy teacher outputs creates a deceptively strong narrow model. Which mechanism should be traced first?
\begin{enumerate}[label=\Alph*.]
\item Knowledge distillation
\item Embedding layer
\item Self-attention math
\item Feed-forward networks
\end{enumerate}
\noindent\textbf{Answer.} Knowledge distillation
\par\textit{A student learns from teacher probabilities, generated data, rationales, or intermediate representations. Tradeoff: Serving cost falls but rare capabilities and calibration may be lost. Watch for: Distilling only easy teacher outputs creates a deceptively strong narrow model.}
\paragraph{MCQ-1037 [hard]} Which causal explanation of Knowledge distillation connects what it does to when it is worth its cost?
\begin{enumerate}[label=\Alph*.]
\item Mechanism: Attention computes scaled query-key scores, applies masking and softmax, then mixes value vectors. Consequence: Every token can interact directly, while standard sequence cost grows quadratically.
\item Mechanism: Transformer MLP blocks expand hidden dimensions, apply nonlinear gates, and project back per token. Consequence: They hold substantial parameters and compute while offering parallel token processing.
\item Mechanism: Learned token vectors map discrete IDs into continuous model space and are often tied to output weights. Consequence: Weight tying saves parameters but constrains input-output geometry.
\item Mechanism: A student learns from teacher probabilities, generated data, rationales, or intermediate representations. Consequence: Serving cost falls but rare capabilities and calibration may be lost.
\end{enumerate}
\noindent\textbf{Answer.} Mechanism: A student learns from teacher probabilities, generated data, rationales, or intermediate representations. Consequence: Serving cost falls but rare capabilities and calibration may be lost.
\par\textit{A student learns from teacher probabilities, generated data, rationales, or intermediate representations. Tradeoff: Serving cost falls but rare capabilities and calibration may be lost. Watch for: Distilling only easy teacher outputs creates a deceptively strong narrow model.}
\paragraph{FC-1038 [medium]} Define Knowledge distillation in interview-ready language.
\noindent\textbf{Answer.} A student learns from teacher probabilities, generated data, rationales, or intermediate representations.
\par\textit{A student learns from teacher probabilities, generated data, rationales, or intermediate representations.}
\paragraph{FC-1039 [hard]} State the main tradeoff and production pitfall for Knowledge distillation.
\noindent\textbf{Answer.} Tradeoff: Serving cost falls but rare capabilities and calibration may be lost. Pitfall: Distilling only easy teacher outputs creates a deceptively strong narrow model.
\par\textit{Tradeoff: Serving cost falls but rare capabilities and calibration may be lost. Pitfall: Distilling only easy teacher outputs creates a deceptively strong narrow model.}
\paragraph{LA-1040 [hard]} Explain Knowledge distillation from first principles, then show how you would evaluate and operate it in production.
\noindent\textbf{Answer.} Start with the mechanism: A student learns from teacher probabilities, generated data, rationales, or intermediate representations. Make the tradeoff explicit: Serving cost falls but rare capabilities and calibration may be lost. Define workload-specific offline metrics and a latency/cost budget, test important slices, instrument the critical path, and create a rollback criterion. The production review must explicitly guard against this failure: Distilling only easy teacher outputs creates a deceptively strong narrow model.
\par\textit{A strong answer connects mechanism, measurable tradeoffs, failure isolation, observability, and rollback rather than listing product features.}
